\documentclass[11pt,letterpaper]{article}
\usepackage[T1]{fontenc}
\usepackage[utf8]{inputenc}
\usepackage{lmodern,microtype}
\usepackage[margin=1in,headheight=15pt]{geometry}
\usepackage{amsmath,amssymb,amsthm,mathtools,mathrsfs}
\usepackage{enumitem,booktabs,array,longtable}
\usepackage{xcolor,fancyhdr,xurl}
\usepackage[colorlinks=true,linkcolor=blue!40!black,citecolor=blue!40!black,urlcolor=blue!40!black]{hyperref}
\usepackage[nameinlink,noabbrev]{cleveref}
\hypersetup{pdftitle={Independent Surreal Copies and Arithmetic Intersections},pdfauthor={AI-assisted research draft},pdfsubject={Prescribed intersections, Hahn fields, linear disjointness, omnific integers, and transcendental descent obstructions}}
\newtheorem{theorem}{Theorem}[section]
\newtheorem{lemma}[theorem]{Lemma}
\newtheorem{proposition}[theorem]{Proposition}
\newtheorem{corollary}[theorem]{Corollary}
\newtheorem{mainthm}{Theorem}
\renewcommand{\themainthm}{\Alph{mainthm}}
\theoremstyle{definition}
\newtheorem{definition}[theorem]{Definition}
\newtheorem{example}[theorem]{Example}
\newtheorem{question}{Research question}
\theoremstyle{remark}
\newtheorem{remark}[theorem]{Remark}
\newtheorem{warning}[theorem]{Warning}
\crefname{mainthm}{Theorem}{Theorems}
\Crefname{mainthm}{Theorem}{Theorems}
\crefname{warning}{warning}{warnings}% [merge] used with the cleveref type hints on labels
\Crefname{warning}{Warning}{Warnings}
\newcommand{\No}{\mathbf{No}}
\newcommand{\Oz}{\mathbf{Oz}}
\newcommand{\On}{\mathbf{On}}
\newcommand{\SC}{\No(i)}
\newcommand{\GOz}{\Oz[i]}
\newcommand{\R}{\mathbb R}
\newcommand{\C}{\mathbb C}
\newcommand{\Q}{\mathbb Q}
\newcommand{\Z}{\mathbb Z}
\newcommand{\N}{\mathbb N}
\newcommand{\K}[2]{\mathcal K_{#1}(#2)}
\newcommand{\A}[2]{\mathcal A_{#1}(#2)}
\newcommand{\B}[2]{\mathcal B_{#1}(#2)}
\newcommand{\PiG}[2]{\Pi_{#1}(#2)}
\newcommand{\supp}{\operatorname{supp}}
\newcommand{\Frac}{\operatorname{Frac}}
\newcommand{\Span}{\operatorname{span}_{\Q}}
\newcommand{\st}{\operatorname{st}}
\newcommand{\ct}{\operatorname{ct}}
\newcommand{\id}{\operatorname{id}}
\newcommand{\trdeg}{\operatorname{trdeg}}
\newcommand{\cf}{\operatorname{cf}}
\newcommand{\im}{\operatorname{im}}
\newcommand{\Mon}{\mathfrak M}
\newcommand{\Alg}{\operatorname{Alg}}
\newcommand{\rk}{\operatorname{rank}}
\newcommand{\val}{\operatorname{v}}
\newcommand{\comp}{\mathbin{\cdot}}
\newcommand{\snapshot}{\texttt{bcac55ae6fd3f2b354e568ba1d2e94d496a2cc59}}
\newcommand{\repo}[1]{\href{https://github.com/VladimirReshetnikov/Surreal/blob/bcac55ae6fd3f2b354e568ba1d2e94d496a2cc59/#1}{\nolinkurl{#1}}}
\newcommand{\rpath}[1]{\nolinkurl{#1}}% a path in the collection, relative to docs/
\newcommand{\lbl}[1]{\nolinkurl{#1}}% a label of another report, cited by name
\newcommand{\wtag}{\textup{[write]}}% text added when the report joined the collection
\newcommand{\mtag}{\textup{[merge]}}% text added when source 02 (batch 34) or source 03 (batch 35) was merged
\newcommand{\crk}{\operatorname{crk}}% coefficient rank (source 02)
\newcommand{\cc}{\mathfrak c}
\newcommand{\dtrdeg}{\operatorname{dtrdeg}}
\newcommand{\Cfg}{\mathcal C_{\mathrm{fg}}}
\newcommand{\ordt}{\operatorname{ord}_t}
\newcommand{\jet}{\mathcal J}
\newcommand{\Ord}{\operatorname{ord}}% ordinal operations (source 03)
\newcommand{\Ad}[3]{\mathcal A^{#1}_{#2}(#3)}% D + purely infinite part (source 03)
\numberwithin{equation}{section}
\setlength{\parskip}{3pt}
\setlength{\parindent}{1.15em}
\setlength{\emergencystretch}{2em}
\setlist[enumerate]{leftmargin=2em,itemsep=3pt,topsep=4pt}
\pagestyle{fancy}
\fancyhf{}
\fancyhead[L]{\small Independent surreal copies}
\fancyhead[R]{\small Arithmetic intersections}
\fancyfoot[C]{\thepage}
\renewcommand{\headrulewidth}{0.3pt}
\setcounter{tocdepth}{2}

\begin{document}
\hypersetup{pageanchor=false}
\begin{titlepage}
\thispagestyle{empty}
\vspace*{0.05in}
{\large\scshape Research report built from three manuscripts\par}
\vspace{0.1in}
{\Huge\bfseries Independent Surreal Copies\par}
\vspace{0.12in}
{\Huge\bfseries and Arithmetic Intersections\par}
\vspace{0.16in}
{\Large Prescribed Hahn cores, algebraic independence,\par
two transcendental descent obstructions, composita,\par
and maximal transcendence of Hahn joins\par}
\vspace{0.12in}
{\large AI-assisted research draft, 23 September 2026\par}
\vspace{0.04in}
{\small Built from manuscripts of batches~31, 34 and~35 of the collection,
pinned to \texttt{bcac55a}, \texttt{7b256e0}, \texttt{f388f95}\par}
\vspace{0.1in}
\hrule
\vspace{0.1in}
\begin{abstract}
Let $H$ be any set-sized divisible additive subgroup of the surreal numbers,
and let $L=\K{\R}{H}$ be the full Conway--Hahn field with exponents in $H$.
We construct an ordinal-indexed proper class of strong, real-fixing,
monomial-preserving self-embeddings of $\No$ whose images are jointly
linearly disjoint over $L$ and have pairwise intersection exactly $L$.
All embeddings fix $L$ pointwise and preserve and reflect the omnific
integers. The same construction gives complex-fixing, conjugation-compatible
copies of $\No(i)$ and its Gaussian omnific ring. Every set of parameters
can be fixed by choosing a suitable $H$.

The proof combines a cut-realization lemma that avoids an arbitrary set of
previously chosen values with a coset-slice criterion for linear disjointness
of full Hahn fields. It yields exact Boolean meet diagrams and descent of
all polynomial relations across their squares. Two consequences distinguish
this construction from mere universality. First, the fraction field of the
intersection of two full omnific copies is a proper subfield of their common
field, with a transcendental gap for every set-sized $H$. For $H=\Q$ a
witness is $1+\sum_{n\ge1}\omega^{-n!}$; real constants already belong to the
smaller fraction field. Second, explicit independent copies with intersection
$\R$ have a full Hahn join containing an explicitly described series
transcendental over their ordinary compositum. Both obstructions persist
on complexification.
\end{abstract}
\vfill
\noindent\textbf{Status and scope.} The article provides complete mathematical
arguments relative to standard Conway normal form and Hahn-field foundations.
The proposed new contribution is the combined prescribed-intersection theorem
and its arithmetic and compositum applications. The ingredients and the
bounded-support transcendence mechanism are not claimed as new. This is an
unrefereed, non-formalized draft, not a certified priority claim or a claimed
solution of a named longstanding conjecture.
\wtag{} Independent proof review and formalization are pending. Labels
carry the prefix \texttt{isc:}. Conventions, the place in the collection,
the non-claims and the provenance are in
\cref{isc:sec:conventions,isc:sec:collection,isc:sec:nonclaims} and
Appendix~\ref{isc:app:report}; added text is marked \wtag.
\mtag{} Sections~15--27 and~28--41 add two more manuscripts
(\texttt{isc:cp:}, \texttt{isc:hj:}).

\smallskip
\noindent\small Repository comparison is pinned to commit\par
\noindent\small\snapshot.\par
\noindent\small Sources, limitations, and a concrete research program are included.
\end{titlepage}
\hypersetup{pageanchor=true}
\setcounter{page}{1}
{\small\setlength{\parskip}{0pt}\tableofcontents}
\clearpage

\section{Three questions and their answers}\label{isc:sec:introduction}

The existence of an embedding into the surreal numbers says little about
how two embedded copies meet. An intersection question asks for much more:
which elements must be shared, which algebraic relations can pass from one
copy to another, and whether their distinguished integer rings reconstruct
the same common field. This article treats these questions for \emph{full
proper-class copies} of the surreal field, rather than only for small
subfields.

The ambient presentation matters. A surreal has a unique Conway normal form
\[
 x=\sum_{g\in S}r_g\omega^g,
\]
where $S\subseteq\No$ is a reverse-well-ordered \emph{set}. For an additive
subgroup $G\subseteq\No$, put
\[
 \K{\R}{G}=\{x\in\No:\supp(x)\subseteq G\},\qquad
 \K{\C}{G}=\K{\R}{G}(i).
\]
These are full Hahn fields: there is no additional bound on the cardinality
or order type of a set-sized support. If $G$ is divisible, the real field is
real closed and the complex field is algebraically closed. These foundations
are classical; see \cite{Gonshor,vdDE,LM}.

\subsection{Prescribing the common field}

\begin{mainthm}[Independent full copies over a prescribed Hahn core]\label{isc:main:copies}
Let $H\subseteq\No$ be a set-sized divisible additive subgroup and write
$L_k=\K{k}{H}$ for $k=\R,\C$. There are class maps
\[
 \Phi_\alpha:\No\longrightarrow\No\qquad(\alpha\in\On)
\]
with the following properties.
\begin{enumerate}[label=\textup{(\roman*)}]
\item Each $\Phi_\alpha$ is a proper ordered field embedding, fixes $L_\R$
pointwise, sends Conway monomials to Conway monomials, and preserves and
reflects Hahn summability of set-indexed families.
\item With $K_\alpha=\Phi_\alpha(\No)$,
\[
 K_\alpha\cap K_\beta=L_\R\quad(\alpha\ne\beta).
\]
Every finite subfamily is jointly linearly disjoint over $L_\R$.
\item For $R_\alpha=\Phi_\alpha(\Oz)$,
\[
 R_\alpha=K_\alpha\cap\Oz,
 \qquad R_\alpha\cap R_\beta=L_\R\cap\Oz.
\]
Moreover $\Phi_\alpha^{-1}(\Oz)=\Oz$.
\item The coefficientwise complexifications have the same assertions over
$L_\C$, commute with conjugation, and preserve and reflect $\GOz$.
\end{enumerate}
\end{mainthm}

The proof is given in \cref{isc:sec:group,isc:sec:lifts,isc:sec:slices,isc:sec:copies}.
``Jointly linearly disjoint'' is stronger than pairwise intersection being
small: finitely many independent choices from different copies satisfy no
new polynomial relations over the common core. Proper-class terminology is
made precise in \cref{isc:sec:foundations}; no tensor product indexed by a proper
class is required.

A useful immediate specialization is $H=\{0\}$. It produces full copies of
$\No$ intersecting exactly in $\R$, while their omnific rings intersect
exactly in $\Z$. A second specialization, $H=\Q$, fixes not just $\omega$
and its rational powers but every Hahn series with rational exponents.
In either case each copy is as large as the whole surreal field.

\subsection{Arithmetic intersection is not field reconstruction}

\begin{mainthm}[Transcendental failure of fraction-field descent]\label{isc:main:fractions}
Choose two copies in \cref{isc:main:copies}. Then
\[
 \Frac(R_\alpha)=K_\alpha,\qquad
 \Frac(R_\alpha\cap R_\beta)
 \subsetneq K_\alpha\cap K_\beta=L_\R.
\]
The latter extension is transcendental for every set-sized $H$, including
$H=\{0\}$. For $H\ne0$, the smaller field already contains every real
constant. For $H=\Q$, the element
\[
 F=1+\sum_{n=1}^{\infty}\omega^{-n!}
\]
is transcendental over $\Frac(R_\alpha\cap R_\beta)$.
The corresponding assertions hold for the Gaussian omnific copies inside
$\No(i)$.
\end{mainthm}

This is a failure of a particular descent operation, not a failure of
$\Frac(\Oz)=\No$. Indeed that classical identity is indispensable to the
argument. Each copy separately supplies denominators for every element of
its field, but the \emph{shared} denominators need not suffice for their
shared field. The phenomenon remains after the missing-real-constants
explanation has been removed. \Cref{isc:sec:fractions,isc:sec:gaps} prove the result.

\subsection{An ordinary compositum is smaller than a full Hahn join}

For full Hahn fields in the same normal-form presentation, their \emph{Hahn
join} is the full Hahn field on the sum of their exponent groups. Their
\emph{ordinary compositum} is the field generated using finitely many field
operations. These should not be conflated.

\begin{mainthm}[A transcendental gap above the compositum]\label{isc:main:join}
There are explicit strong embeddings $\Phi_0,\Phi_1$ as in
\cref{isc:main:copies} with $H=0$ for which, writing
\[
 K_j=\Phi_j(\No),\qquad
 E=\K{\R}{G_0+G_1},\qquad M=K_0K_1,
\]
the field $E$ contains an explicit countably supported element $S$
transcendental over $M$. In particular even the relative real algebraic
closure of the ordinary compositum is a proper subfield of the Hahn join.
The same $S$ is transcendental over the compositum of the complexified
copies inside $\K{\C}{G_0+G_1}$.
\end{mainthm}

The two-level normal-form construction is explicit in \cref{isc:sec:explicit}.
The formula for $S$ and the coefficient-field obstruction are proved in
\cref{isc:sec:join}. The latter argument is independent of the lacunary-series
argument used for \cref{isc:main:fractions}.

\subsection{What is new here, and what is not}\label{isc:subsec:provenance}

The classical inputs are the set-cut property of $\No$, its normal-form
presentation, Hahn multiplication and inversion, and the real/algebraic
closedness criteria. The group extension argument is a concrete saturation
construction, not a new saturation principle. The monomial lifting argument
is standard Hahn functoriality. The coset-slice proof is presented as an
elementary reusable lemma rather than as a certified new theorem of Hahn
field theory.

The relevant repository comparison used the documentation catalogue and the
reports on omnific-preserving automorphisms and bounded-support Hahn
arithmetic at the pinned snapshot \cite{RepoCatalogue,RepoAuto,RepoBounded}.
The former concerns stabilizers and rigidity; the latter already proposes
much larger algebraically independent families over bounded-support fraction
fields. We neither repackage those results as discoveries nor use their
unrefereed main theorems as black boxes. A self-contained one-witness gap
argument is supplied here. In the published literature, the failure of
fraction-field equality at set-sized cofinality already appears in
L'Innocente--Mantova, Proposition~2.4.5 \cite{LM}.

The proposed contribution is instead the \emph{simultaneous arrangement}:
full copies, a prescribed set-sized full Hahn intersection, strong
arithmetic compatibility, independence of every finite subfamily, exact meet
diagrams, and the two transcendental descent obstructions inside those
arrangements. Targeted literature searches do not certify that this combined
formulation is absent from all published or unpublished work. No claim of
having audited every report or every later repository revision is made.

\wtag{} This subsection is the source's, written at its pin. Its
description of the omnific-preserving report as concerning stabilizers and
rigidity was accurate then: at the pin that report had no part on
embeddings. Its Part~III, written afterwards, classifies the strong
exact-monomial embeddings that preserve $\Oz$, and the lifts used here are
the real-fixing case of that classification. The bounded-support report was
unchanged between the pin and batch~31; batch~32 added its Galois part (its
Sections~11--19). The attribution to L'Innocente--Mantova was checked
(Appendix~\ref{isc:app:literature}). The relation to the collection as it
now stands, with exact labels, is in \cref{isc:sec:collection}.

\subsection{Notation and conventions in this report}
\label{isc:sec:conventions}

\wtag{} This subsection and the next were written when the manuscript
joined the collection. No symbol of the source was renamed and no
normalization was changed. The meanings below are fixed once for the whole
report, and the collisions with other reports are disclosed here rather
than removed. Paths such as \rpath{surreal/omnific-preserving-automorphisms}
are relative to \rpath{docs/}, and labels of other reports are cited by
name.
\begin{itemize}[leftmargin=*,itemsep=3pt]
\item \emph{Two exponent conventions.} \Cref{isc:sec:introduction} to
\cref{isc:sec:join}, \cref{isc:sec:setversion} and the beginning of
\cref{isc:sec:fractions} use Conway's decreasing large-$\omega$ convention
$x=\sum r_g\omega^g$. There $\Oz$ consists of the series supported in
$\No^{>0}\cup\{0\}$ with integral constant coefficient: the purely infinite
part sits at \emph{positive} exponents. From \cref{isc:sub:bounded} to the end of
\cref{isc:sec:gaps} the source switches to $t^h=\omega^{-h}$, the convention
of the omnific-preserving report's Part~III (\lbl{opa:as:sub:notation}) and of
the bounded-support report (\lbl{bst:sec:setup}); there the purely infinite
part has \emph{negative} $t$-exponents and $v$ is the least $t$-exponent.
The tempting false reading of $\B{k}{H}$ is ``support bounded above in the
Conway exponents''. The true reading is ``bounded above in the
$t$-exponents'', that is, Conway exponents bounded \emph{below}.
\item \emph{Fields.} $\K{k}{G}$ is the \emph{full} Hahn field on $G$,
support-closed and not generated by monomials; for $k\ne\R,\C$ it is
formal. $L_k=\K{k}{H}$, with $H$ a set-sized divisible subgroup of $\No$.
This $H$ need not be convex. It is not a convex subgroup $H_g$ of the
omnific-preserving report (\lbl{opa:def:convex}), and not the large-cardinal
report's embedding $H_j$.
\item \emph{Rings.} $\A{k}{H}=D_k\oplus\PiG{k}{H}$ with $D_\R=\Z$,
$D_\C=\Z[i]$ restricts the constant coefficient to $D_k$. It is the ring
$\mathcal R_\Gamma$ (or $\mathcal R^\C_\Gamma$) of the omnific-preserving
report's Part~III and the ring $\mathcal R_{D_k}(k,\Gamma)$ of the
Diophantine report's Section~11. It is \emph{not} that Part's
$\mathcal A_\Gamma=\R\oplus\Pi_\Gamma$. Nor is it the Diophantine report's
$\mathcal A_k(\Gamma)$, whose constant coefficient ranges over all of $k$.
$D_k$ is a ring of integers and never a derivation; the omnific-preserving
report reserves $d$ and $D$ for derivations. $\B{k}{H}$ is the
bounded-support report's $\mathcal B_H(k)$, with the two arguments in the
opposite order, and $\Frac\B{k}{H}$ is its $\mathcal F_H(k)$.
\item \emph{Maps.} $\Phi_h$ of \eqref{isc:eq:lift} is the monomial map
$M_{1,h}=J_{h,\id}$ of \lbl{opa:as:eq:classification}. It is not the
Taylor automorphism $\Phi_d$ of the surcomplex-automorphisms report
(the omnific-preserving report's $\mathcal T_{d,\delta}$), and not the
large-cardinal report's sign-sequence embedding $J$. The letters $h_\alpha$,
$G_\alpha$, $\Phi_\alpha$, $K_\alpha$ denote the family of
\cref{isc:thm:group} in \cref{isc:sec:group,isc:sec:lifts,isc:sec:slices,isc:sec:copies}
and the explicit family in \cref{isc:sec:explicit,isc:sec:join}. There
$h_\alpha$ is the \emph{inner} map \eqref{isc:eq:inner}, which is the inner
lift $\iota_{\sigma_\alpha}$ of \lbl{opa:as:eq:innerlift}. The maps
$h_\alpha$ fix $H$ as a set of exponents, and $\Phi_\alpha$ fixes
$\K{\R}{H}$ (\cref{isc:rem:base}). ``Monomial-preserving'' is the
omnific-preserving report's ``exact-monomial'': a Conway monomial goes to a
Conway monomial with coefficient one. It does not mean commuting with the
omega-map. The order isomorphism $\psi$ of \cref{isc:lem:compression} is
not that report's $\mu_{12}$ (its source~12's $\psi$), nor $\mu_{11}$ or
$\mu_{13}$. The slice $\lambda_{g,H}$ of \eqref{isc:eq:slice} is the
bounded-support report's shifted coset projection $\pi_g$
(\lbl{bst:eq:projection}); it is not the class functional $\lambda_{\On}$.
\item \emph{Local letters.} $L,R$ are the sides of a cut in
\eqref{isc:eq:cut} and \cref{isc:lem:avoid,isc:lem:extend}; $L$ is also a
generic base field in the definition of linear disjointness and in
\cref{isc:lem:matrix}. $M=K_0K_1$ is an ordinary compositum, not a
monomial map $M_{\chi,\tau}$. $E$ is a span in the proof of
\cref{isc:thm:group}, the field $E_J$ of \eqref{isc:eq:diagram}, and the
Hahn join of \cref{isc:main:join}. $F$ is the witness of
\cref{isc:main:fractions} and \eqref{isc:eq:cofinalF}, and a polynomial in
\cref{isc:cor:relations}. It is not the image field $F$ of
\lbl{opa:as:prop:overlap}, and not the fraction field $\mathcal F_G(K)$.
$A$ is a set of parameters in \cref{isc:cor:parameters}, a matrix in
\cref{isc:lem:matrix} and a subgroup in \cref{isc:lem:iterated}, where $B$
is a subgroup too; the calligraphic $\mathcal A$, $\mathcal B$ are the
rings. $D$ is a domain or subspace in \cref{isc:sec:group} and a
coefficient field in \cref{isc:sec:join}. $S$ is a support set and the
witness \eqref{isc:eq:explicitS}. $I$ is an index set and $I_\alpha$ an
interval.
\item \emph{Cardinals.} In \cref{isc:sec:gaps}, $\kappa=\cf(H)$. In
\cref{isc:sec:setversion}, $\kappa$ is an arbitrary uncountable regular
cardinal. Neither is the critical point $\kappa$ of the large-cardinal
report.
\item \emph{Strong} means preserving Hahn summability of set-indexed
families; it is not convergence of partial sums.
\end{itemize}

\subsection{Place in the collection}\label{isc:sec:collection}

\wtag{} The source compared itself with the catalogue, the
omnific-preserving report and the bounded-support report at its pin
(Appendix~\ref{isc:app:pinned}). Checked against the collection at the
placement commit, the relations are these.

\paragraph{Omnific-preserving automorphisms}
(\rpath{surreal/omnific-preserving-automorphisms}, prefix \texttt{opa:}).
Its Part~III (sub-prefix \texttt{opa:as:}) was written after the pin, so the
source did not see it. Its bottom-gap classification
\lbl{opa:as:thm:classification} lists the strongly additive field embeddings
$J$ with $J(t^\gamma)=t^{\tau(\gamma)}$ and $J(\mathcal R_\Gamma)\subseteq
\mathcal R_\Delta$ as the maps $J_{\tau,\rho}$. Here $\rho$ ranges over
the ordered embeddings of $\R$ into the Hahn field of the bottom gap of
$\tau$.
\begin{enumerate}[leftmargin=*,itemsep=3pt]
\item \emph{Lifts.} With $\Gamma=\Delta=\No$, $t^\gamma=\omega^{-\gamma}$
and $\tau=h$, the lift $\Phi_h$ of \cref{isc:thm:lift} is $J_{h,\id}$, the
case $\rho=\id$. The same exponent map serves in both conventions because
$h$ is additive. The preservation and reflection of $\Oz$ in
\cref{isc:prop:arithmeticlift} is Step~5 of that proof and the display
\lbl{opa:as:eq:classprops}. \Cref{isc:prop:classification} is the
uniqueness clause of Step~2, for real-fixing maps; here the order of $h$ is
derived from real fixing. The floor, standard-part and modulus clauses of
\cref{isc:prop:arithmeticlift} are not stated there.
\item \emph{Real fixing is a choice.} The explicit ranges $G_\alpha$ of
\cref{isc:sec:explicit} have a nonzero bottom gap: $\omega^{3\alpha}$ lies
below every positive element of $G_\alpha$, and for $\alpha=0$ the gap is
the group of finite surreals. By \lbl{opa:as:thm:taylorcriterion}, the
composite of a Taylor automorphism $\mathcal T_{d,\delta}$ with $\Phi_\alpha$,
for $0<\delta$ in that gap, therefore also preserves and reflects $\Oz$ and
moves a transcendental real. The source takes $\rho=\id$ throughout, which
is consistent. For the family of \cref{isc:thm:group} the construction does
not determine the gap. \Cref{isc:thm:intersectionclassification} depends on
the real-fixing hypothesis: in the larger exact-monomial class, images meet
$\R$ in the fixed field of $\rho$, for example $\ker d$
(\lbl{opa:as:prop:overlap}). No classification of intersections in that
class is claimed.
\item \emph{Inner and outer levels.} The inner map \eqref{isc:eq:inner} is
the inner lift $\iota_{\sigma_\alpha}$ of \lbl{opa:as:eq:innerlift} and
\lbl{opa:as:lem:inner}, and $\Phi_\alpha=M_{1,\iota_{\sigma_\alpha}}$. The
order isomorphisms of \lbl{opa:as:eq:mus} have unbounded images; the
explicit embeddings \lbl{opa:as:thm:explicit} combine them with Taylor
sections in order to move a real. The maps $\sigma_\alpha$ here have
bounded, pairwise disjoint images, which is what gives independence. The
failure of omega-map commutation, $\Phi_0(\omega^1)\ne\omega^{\Phi_0(1)}$,
is the same observation as the one recorded after
\lbl{opa:as:cor:homogeneity}.
\item \emph{Parameters.} \Cref{isc:cor:parameters} extends
\lbl{opa:as:thm:parameters}, which gives two proper strong exact-monomial
copies fixing $A\cup\R$ and preserving and reflecting $\Oz$. Here there is an
$\On$-indexed family whose pairwise intersection is exactly the full Hahn
field $\K{\R}{H}\supseteq A\cup\R$, jointly linearly disjoint over it.
The mechanisms differ. That theorem fixes the inner supports through an
order embedding (\lbl{opa:as:lem:parameters}); here the outer exponent
group $H$ is fixed by $h_\alpha$. The topological assertions of that
theorem (closed, discrete or nondiscrete images) are not claimed here, and
that report does not study intersections.
\item \emph{Complex version.} The complex clauses of \cref{isc:main:copies}
and \cref{isc:thm:explicit} are the case $\varepsilon=1$, $\rho=\id$ of
\lbl{opa:as:thm:complexembed}, parallel to \lbl{opa:as:cor:complexcopies}.
\item \emph{Elementarity.} The ordered-field statement of \cref{isc:sub:model} is the first assertion of
\lbl{opa:as:thm:elementary}. That theorem's nonelementarity criterion
needs a moved real, and the maps here fix $\R$.
\item \emph{Questions.} None is answered. Research question~2 is the pair
version, at a fixed Hahn core, of \lbl{opa:as:q:conjugacy}. Research
question~9 asks \lbl{opa:as:q:elementary} for the maps here, with a
prescribed intersection added. Research questions~6 and~7 lie inside
\lbl{opa:as:q:exp}. \Cref{isc:thm:boolean} decides image inclusion and
intersection inside one independent family with $\rho=\id$. That touches
\lbl{opa:as:q:composition}, which asks for image inclusion of classified
maps in general, but does not answer it. Research question~1 is related to
\lbl{opa:as:q:units}, which asks to remove exact monomial preservation, but
it asks about intersections and is new.
\end{enumerate}

\paragraph{Bounded-support transcendence}
(\rpath{surreal/transcendence-over-bounded-support}, prefix \texttt{bst:}).
The source cites it and says that \cref{isc:sec:gaps} duplicates its
one-witness case. That report was unchanged between the pin and batch~31;
batch~32 added its Galois part (its Sections~11--19), which is not used here.
\begin{enumerate}[leftmargin=*,itemsep=3pt]
\item \Cref{isc:prop:boundedfrac} is the localization of
\lbl{bst:prop:localdensity}, with the constant coefficient restricted to
$D_k$. Since batch~32 that report also prints the general form, for every
field $K$ and every unital subring $D\subseteq K$, as
\lbl{bst:gr:prop:fractions}.
\item \Cref{isc:thm:cofinalgap} is \lbl{bst:cor:onegap} when
$\cf(H)=\aleph_0$; that corollary's hypothesis forces a cofinal support of
countable cofinality. In characteristic zero and at every cofinality it is
the one-series case of \lbl{bst:thm:mainintro} (\lbl{bst:thm:independent}),
which gives $2^\kappa$ independent series on one support of order type
$\kappa$ and proves that $\kappa$ is optimal (\lbl{bst:cor:optimal}). The
single-series form at uncountable cofinality over an arbitrary coefficient
field, finite fields included, is not stated there; that report says that
finite coefficient fields are not covered by its independence theorem. Its
proof is that report's tail estimate and root count.
\Cref{isc:lem:gapscale} is the second inequality of a separated exponent
system (\lbl{bst:lem:skeleton}).
\item The factorial witness of \cref{isc:cor:factorial} is
\lbl{bst:cor:onegap} with $a_n=n!$ and all coefficients one; the constant
$1$ changes nothing. \lbl{bst:thm:actualreal} uses the same exponents, with
integer codes, for a continuum of independent actual infinitesimals over the
complex bounded-support fraction field. In the omnific-preserving report,
\lbl{opa:as:ex:smallfrac} uses $\sum t^{n!}$ at $\Gamma=\Z$ to show that
$\sum t^{n!}$ is not in $\Frac\mathcal R_\Z=\R(t^{-1})$: non-membership,
not transcendence.
\item Research question~8 is that report's programme
(\lbl{bst:cor:exactcard}, \lbl{bst:cor:realgroup}), as the source says.
Only its uniformity over realized diagrams is new.
\item The coset slice $\lambda_{g,H}$ is that report's projection $\pi_g$,
and the linearity in \cref{isc:lem:slice} is \lbl{bst:lem:projection}.
There it serves descent over a cofinal subgroup
(\lbl{bst:thm:descent}). \Cref{isc:thm:sliceddisjoint} is not stated there:
two full Hahn fields are linearly disjoint over the full Hahn field of the
intersection of their groups, with no cofinality hypothesis. The source
presents it as an elementary reusable lemma, not as a certified new theorem.
\end{enumerate}

\paragraph{Omnific Diophantine geometry}
(\rpath{surreal/omnific-diophantine-geometry}, prefix \texttt{odg:}).
\Cref{isc:lem:denominator} is the global monomial clearing
\lbl{odg:thm:fractions}, with the same proof.

\paragraph{Large-cardinal embeddings}
(\rpath{foundations-and-computation/large-cardinal-embeddings-and-normal-forms},
title ``Critical-Point Defects in Surreal Arithmetic''). It was placed after
the pin, so the source did not see it. Its write phase is in progress, so
it is cited here by name and question title, not by label. An elementary
embedding with critical point $\kappa$ gives there two ordered-field
embeddings $J$ and $H_j$ of $\No$. The strong one, $H_j$, is the lift
$\Phi_h$ of \cref{isc:thm:lift} with $h$ the restriction of $J$ to
$(\No,+,<)$, since $H_j(\omega^a)=\omega^{J(a)}$. The map $J$ is not strongly
additive, and $J(\No)$ is not a full Hahn field on a subgroup. The common
field there (surreals with fewer than $\kappa$ terms and exponents in
$J(\No)$) is a proper class, not a full Hahn field. That report's question
``Joint algebraic dependence of the images'' asks whether $J(\No)$ and
$H_j(\No)$ are linearly disjoint over that common field. It is related to
\cref{isc:thm:sliceddisjoint} but \emph{not answered} by it: the slice
theorem needs full Hahn fields on subgroups on both sides.

\paragraph{Formalization.} No \texttt{isc:} label has a Lean mapping in
\rpath{FORMALIZATION.md}. No theorem of this report is Lean-verified, and
the identifiers proposed in \cref{isc:sub:formal} do not exist in the
collection's Lean library.

\paragraph{Status of the research questions.} None of the collection's
named questions is answered. Research questions~1, 3, 4, 5 and~10 are new.
Research question~2 is a pair version of \lbl{opa:as:q:conjugacy}, 6 and~7
lie inside \lbl{opa:as:q:exp}, 8 is the bounded-support programme, and 9 is
an instance of \lbl{opa:as:q:elementary}. Research question~10 overlaps
\lbl{opa:as:q:formal} in its Hahn-field kernel.
\mtag{} \emph{Status (batch 34).} Research question~3 is now answered in
the separated case with trivial common core, the case of
\cref{isc:main:join}, and in the natural Laurent-series cardinal cases, by
the second manuscript printed in
Sections~\ref{isc:cp:sec:intro}--\ref{isc:cp:sec:verify}
(\cref{isc:cp:sec:answers}); the non-separated case and prescribed nonzero
common cores stay open. Research questions~11--22 come from that
manuscript and are new. The relations of the addition to the three-duals,
tail-spans, bounded-support and dilation reports are in
\cref{isc:cp:sec:collection}.
\mtag{} \emph{Status (batch 35).} A third manuscript, printed in
Sections~\ref{isc:hj:sec:intro}--\ref{isc:hj:sec:verify}, answers the
transcendence-degree part of Research questions~3
and~\ref{isc:cp:q:amalgams} for explicit \emph{non-separated} examples
(interleaved groups at every infinite cardinal, dense Archimedean groups,
and a class example); the description of the compositum in the
non-separated case, and prescribed nonzero cores, stay open
(\cref{isc:hj:sec:answers}). Research questions~23--32 come from that
manuscript and are new. Its relations to the tail-spans, Diophantine,
set-sized-quotient, bounded-support and omnific-preserving reports are in
\cref{isc:hj:sec:collection}.

\section{Foundations, support, and size}\label{isc:sec:foundations}

\subsection{A class framework with set-sized constructions}

We work in von Neumann--Bernays--G\"odel class theory with Global Choice,
with sets satisfying ZFC, as in the foundational convention of
\cite{vdDE}. Classes are not members of other classes. Thus an
ordinal-indexed family of class functions means one class relation whose
sections are the functions. A field with carrier $\No$ is a \emph{class
field}; its elements, finite tuples, polynomial coefficient lists, and the
supports of individual elements are sets.

The one surreal-specific order principle needed for the construction is:
\begin{equation}\label{isc:eq:cut}
 L,R\subseteq\No\text{ sets},\quad L<R
 \quad\Longrightarrow\quad
 \text{there is }y\in\No\text{ with }L<y<R.
\end{equation}
Empty sides are permitted. The simplest such $y$, denoted $\{L\mid R\}$,
is uniquely determined. We never apply \eqref{isc:eq:cut} to proper-class sides.
This restriction is central, not cosmetic.

Global Choice supplies a set-like class well-order of the surreals: order
first by birthday and well-order each set-sized birthday level. It similarly
supplies a set-like enumeration of $\On\times\No$. For the latter, group
pairs by a bound on both their ordinal coordinate and the birthday of the
surreal coordinate, then well-order each new set layer. Each initial segment
is a set, and the resulting class order has order type $\On$.

The recursion below has a \emph{set as its entire state at every ordinal
stage}. For any fixed ordinal, ordinary set transfinite recursion with the
fixed class parameters constructs that initial segment. Uniqueness makes
these initial segments coherent, and elementary class comprehension collects
their stages. This is not an appeal to recursion with a proper-class-valued
state at each stage, nor to an unmentioned class truth predicate.

\subsection{Full Hahn fields in decreasing-exponent convention}

If $G$ is an ordered additive subgroup of $\No$ and $k$ is a set-sized
field, let $\K{k}{G}$ denote the formal series
\[
 \sum_{g\in S}a_g\omega^g,\qquad a_g\in k,
\]
with reverse-well-ordered set support $S\subseteq G$. For $k=\R$ these
are actual surreals. For $k=\C$, taking real and imaginary coefficients
identifies them with actual elements of $\SC$. The union of the two real
supports is again reverse well ordered, so this complexification is exact.
For other $k$ the notation is purely formal.

A nonzero series has a largest exponent, and its leading coefficient is
nonzero. Addition is coefficientwise and multiplication is convolution.
The Hahn support lemma guarantees that the support of a product is allowable
and that each product coefficient is a finite sum. Normalizing a nonzero
series as $a\omega^g(1+\varepsilon)$ with all exponents of $\varepsilon$
negative gives its inverse by a summable geometric family. Consequently
$\K{k}{G}$ is a field. These statements are understood locally in set-sized
subgroups when $G$ is a proper class.

The standard closedness theorem says that a Hahn field of characteristic
zero with divisible exponent group is real closed when its coefficients are
real closed, and algebraically closed when its coefficients are algebraically
closed \cite{LM}. The same conclusion holds for our class fields: the
supports of finitely many coefficients of a polynomial generate a set-sized
divisible subgroup, and the standard set-sized theorem provides the required
roots there. In particular,
\[
 \K{\R}{G}\text{ is real closed},\qquad
 \K{\C}{G}\text{ is algebraically closed}
 \quad\text{if }G\text{ is divisible}.
\]

If $H$ is a set, then $\K{k}{H}$ is a set: it is a subclass of the set
of all functions $H\to k$. In contrast, if $G$ is class-isomorphic to
$(\No,+,<)$, then $\K{k}{G}$ is a proper class, because it contains the
proper class of monomials $\omega^g$.

\subsection{Strong summation and finite algebraic independence}

\begin{definition}[Hahn summability]
A set-indexed family $(x_j)_{j\in J}$ of Hahn series is summable if the union
of its supports is reverse well ordered and every exponent belongs to the
support of only finitely many $x_j$. Its sum is the resulting coefficientwise
finite sum at every exponent. A map is \emph{strongly additive} if it sends
such families to summable families and commutes with their sums.
\end{definition}

The sum of a normal form is strong in this sense. Ordinary analytic
convergence of a sequence of real constants is a different notion; it is
not used in any proof. In particular strong additivity cannot be replaced by
an unqualified assertion of continuity.

\begin{definition}[Linear disjointness for class fields]
Subfields $P,Q$ of one class field are linearly disjoint over $L\subseteq
P\cap Q$ if every finite $L$-linearly independent tuple from $P$ remains
linearly independent over $Q$. A finite family is jointly linearly disjoint
if its multiplication map on finite tensors over $L$ is injective.
Equivalently, this injectivity holds after restricting each factor to any
set-sized subfield containing the finitely many elements under consideration
and the base when the base is a set.
\end{definition}

All tensor arguments below can instead be read as finite matrix arguments.
For a class-sized base, linear independence still quantifies only over a
finite tuple of coefficients at a time. ``Algebraically independent'' also
has its usual finite-polynomial meaning. We make no assertion about a
proper-class-dimensional tensor object in the category of sets.

\section{Independent ordered-group embeddings}\label{isc:sec:group}

The core construction is carried out on the \emph{additive exponent group},
not initially on the field. Its flexibility comes from the ability to realize
a set-sized cut while avoiding all previously used values.

\begin{lemma}[A cut avoiding a set]\label[lemma]{isc:lem:avoid}
Let $L,R,E\subseteq\No$ be sets and suppose $L<R$. There is a surreal
$y\notin E$ with $L<y<R$.
\end{lemma}
\begin{proof}
Let $T=\{e\in E:L<e<R\}$. This is a set and $L\cup T<R$.
Choose
\[
 y=\{L\cup T\mid R\}.
\]
Then $L<y<R$. Were $y\in E$, it would belong to $T$, contradicting
$T<y$. This also covers empty $L$, $R$, or $T$.
\end{proof}

\begin{lemma}[Ordered one-point extension with avoidance]\label[lemma]{isc:lem:extend}
Let $D$ be a set-sized $\Q$-vector subspace of $\No$, let
$f:D\to\No$ be an ordered additive embedding, and let $E$ be a set-sized
$\Q$-vector subspace containing $f(D)$. For $x\notin D$, there is an
ordered additive embedding
\[
 f^+:D+\Q x\longrightarrow\No
\]
extending $f$ and such that $f^+(x)\notin E$.
\end{lemma}
\begin{proof}
Set
\[
 L=\{f(d):d\in D,\ d<x\},\qquad
 R=\{f(d):d\in D,\ x<d\}.
\]
These are sets with $L<R$. By \cref{isc:lem:avoid}, choose $y\notin E$
realizing this cut. Define
\[
 f^+(d+qx)=f(d)+qy\qquad(d\in D,\ q\in\Q).
\]
The representation $d+qx$ is unique because $x\notin D$, so the formula is
well defined and additive. If $q>0$, positivity of $d+qx$ is equivalent to
$x>-d/q$; by the defining cut this is equivalent to $y>f(-d/q)$, hence to
$f(d)+qy>0$. For $q<0$, the same argument reverses the comparison. The
case $q=0$ follows from the order preservation of $f$. Thus $f^+$ preserves
strict positivity and is an ordered embedding. Its new value is $y\notin E$.
\end{proof}

\begin{definition}[Independence over an additive subgroup]
Subgroups $(G_i)_{i\in I}$ containing $H$ are \emph{independent over $H$}
if, whenever $J\subseteq I$ is finite and $g_i\in G_i$,
\[
 \sum_{i\in J}g_i\in H\quad\Longrightarrow\quad
 g_i\in H\text{ for every }i\in J.
\]
Equivalently the images $G_i/H$ have internal direct sum in the quotient
additive group. The quotient here is algebraic; $H$ need not be convex.
\end{definition}

\begin{theorem}[Independent group copies over any set-sized base]
\label{isc:thm:group}
Let $H\subseteq\No$ be a set-sized divisible additive subgroup. There is
an ordinal-indexed class family of ordered additive embeddings
\[
 h_\alpha:\No\longrightarrow\No
\]
such that $h_\alpha|_H=\id_H$ and the ranges $G_\alpha=h_\alpha(\No)$
are independent over $H$.
\end{theorem}
\begin{proof}
A divisible subgroup of the torsion-free additive group of $\No$ is a
$\Q$-vector subspace. Fix the set-like enumeration of $\On\times\No$
described in \cref{isc:sec:foundations}. We process each requirement $(\alpha,x)$
in that order, meaning that $x$ must eventually be in the domain of the
$\alpha$th partial embedding.

At a stage the state consists of partial ordered $\Q$-linear embeddings
$f_\alpha:D_\alpha\to\No$ for a set of active indices. Their domains and
graphs are sets, they contain the identity on $H$, and their ranges are
independent over $H$. An inactive index has the implicit partial map
$\id_H$; these implicit copies of $H$ need not be stored separately.

Suppose the next requirement is $(\alpha,x)$. Activate $\alpha$ if needed.
If $x\in D_\alpha$, do nothing. Otherwise let
\[
 E=\Span\left(H\cup
     \bigcup_{\beta\text{ active}}f_\beta(D_\beta)\right).
\]
The state is a set, so $E$ is a set-sized vector space. Apply
\cref{isc:lem:extend} to extend $f_\alpha$ to $D_\alpha+\Q x$, choosing its
new value $y$ outside $E$.

We verify that independence survives. A new finite relation modulo $H$
can involve at most one changed component. Write that component as
$f_\alpha(d)+qy$. If $q\ne0$, the relation would express $y$ as a rational
linear combination of old range elements and elements of $H$, contradicting
$y\notin E$. Therefore $q=0$ and the old independence assertion applies.
This proves the invariant at successor stages.

At a limit ordinal, take unions of the compatible graphs. There have been
only a set of preceding stages, so the resulting state is still a set.
Every finite relation already occurred at an earlier stage, so independence
is retained. The set-valued recursion is legitimate in the stated class
framework.

Finally define $h_\alpha$ to be the union of its partial graphs over all
stages. Each requirement was processed, hence its domain is all of $\No$.
Every ordered-additive identity is witnessed at some stage, so $h_\alpha$
is an ordered additive embedding fixing $H$. Any finite relation among
final range elements is also witnessed at one stage. Consequently the final
ranges are independent over $H$.
\end{proof}

\begin{corollary}\label[corollary]{isc:cor:groupintersection}
For distinct $\alpha,\beta$, one has $G_\alpha\cap G_\beta=H$. More
generally, for $\alpha\notin J$,
\[
 G_\alpha\cap\left(H+\sum_{\beta\in J}G_\beta\right)=H,
\]
where the sum means finite sums, even when $J$ is infinite. Every $G_\alpha$
is a proper subgroup of $\No$.
\end{corollary}
\begin{proof}
The first two assertions are exactly the independence condition. To prove
properness, choose $\beta\ne\alpha$ and $x\notin H$. Injectivity and the
identity on $H$ imply $h_\beta(x)\notin H$. Thus it cannot lie in
$G_\alpha$, although it lies in $\No$.
\end{proof}

\begin{remark}[The role of the base]\label[remark]{isc:rem:base}
The maps $h_\alpha$ fix $H$ as a set of \emph{exponent values}. The field
maps constructed next fix $\K{\R}{H}$. These are different assertions:
for a general subgroup $H$, it need not be true that $H\subseteq\K{\R}{H}$.
No identification of those two roles is used.
\end{remark}

\section{Strong field lifts and omnific compatibility}\label{isc:sec:lifts}

\begin{theorem}[Monomial lift]\label{isc:thm:lift}
Let $h:\No\to\No$ be an ordered additive embedding. Then
\begin{equation}\label{isc:eq:lift}
 \Phi_h\left(\sum_g r_g\omega^g\right)
   =\sum_g r_g\omega^{h(g)}
\end{equation}
defines a real-fixing ordered field embedding of $\No$ with image exactly
$\K{\R}{h(\No)}$. It preserves and reflects Hahn summability, and its
complexification has image $\K{\C}{h(\No)}$.
\end{theorem}
\begin{proof}
An order embedding sends a reverse-well-ordered set to a reverse-well-ordered
set, so the formula defines a surreal. It preserves the leading coefficient
and therefore the sign of a nonzero real-coefficient series. Coefficientwise
addition is preserved. For multiplication, the identity
$h(g+g')=h(g)+h(g')$ identifies each convolution coefficient on the two
sides; the support lemma ensures that only finitely many terms contribute
to each such coefficient. Since $h(0)=0$, real constants are fixed. The
map is injective by uniqueness of normal form, hence is a field embedding.

For a set-indexed family, the union of image supports is the $h$-image of
the union of original supports. Reverse well-ordering and pointwise
finiteness are equivalent under this injection. Thus summability is both
preserved and reflected, and the sums agree coefficientwise.

Finally, a series supported in $h(\No)$ can be pulled back term by term
along the inverse of $h$ on its image. Its support remains a set and is
reverse well ordered. This proves that the image is the \emph{full} Hahn
field on $h(\No)$, not merely the field generated by its monomials. Taking
real and imaginary coefficients gives the complex assertion.
\end{proof}

\begin{proposition}[What the lift remembers]\label[proposition]{isc:prop:arithmeticlift}
The map $\Phi_h$ preserves the three normal-form projections to positive,
zero, and negative exponents. It preserves and reflects membership in
$\Oz$, and
\begin{equation}\label{isc:eq:ozimage}
 \Phi_h(\Oz)=\Phi_h(\No)\cap\Oz.
\end{equation}
It commutes with the omnific floor map and with standard part on finite
surreals. Its complexification commutes with conjugation, preserves and
reflects $\GOz$, and preserves the real-valued modulus in the sense
\[
 |\Phi_h^{\C}(z)|=\Phi_h(|z|).
\]
\end{proposition}
\begin{proof}
Because $h$ is ordered and additive, $h(g)>0$, $h(g)=0$, and $h(g)<0$
are equivalent to the respective conditions on $g$. The coefficient at zero
is unchanged. An omnific integer is precisely a series supported in
$\No^{>0}\cup\{0\}$ with integral constant coefficient, so membership is
preserved and reflected. The image formula follows by pulling back a series
in the intersection using \cref{isc:thm:lift}.

If $a=\lfloor x\rfloor_\Oz$, then $a\in\Oz$ and $a\le x<a+1$.
Applying $\Phi_h$ gives an omnific integer in the corresponding unit interval
for $\Phi_h(x)$, hence the unique omnific floor. If $x=r+\varepsilon$ is
finite with $r\in\R$ and $\varepsilon$ infinitesimal, the image has the
same decomposition with infinitesimal part $\Phi_h(\varepsilon)$. This
proves the standard-part assertion.

Complex conjugation is coefficientwise. Gaussian omnific membership is
omnific membership of both real and imaginary parts. Finally $|z|$ is the
unique nonnegative square root of $z\bar z$. The image of this equation,
order preservation on the real part, and uniqueness of that square root
prove the modulus identity.
\end{proof}

\begin{remark}[A floor boundary that must not be dropped]\label[remark]{isc:rem:floor}
The positive-exponent truncation plus the ordinary floor of the constant
coefficient is not always the omnific floor. For example, if $0<\varepsilon$
is infinitesimal and $n\in\Z$, then
$\lfloor n-\varepsilon\rfloor_\Oz=n-1$. The order-and-uniqueness proof above
retains this endpoint case automatically.
\end{remark}

\begin{proposition}[Classification within the monomial-preserving class]
\label[proposition]{isc:prop:classification}
Every strongly additive, real-fixing field embedding
$\Phi:\No\to\No$ that maps each Conway monomial to a Conway monomial
is uniquely $\Phi_h$ for an ordered additive embedding $h:\No\to\No$.
\end{proposition}
\begin{proof}
Define $h(g)$ uniquely by $\Phi(\omega^g)=\omega^{h(g)}$. Multiplication
of monomials proves additivity, and injectivity proves injectivity of $h$.
A field embedding between real closed ordered fields preserves positivity:
a positive element is a nonzero square. Hence it preserves order. For
$g>0$, the monomial $\omega^g$ exceeds every positive real number. Its image
has the same property because the reals are fixed. A Conway monomial has
this property exactly when its exponent is positive, so $h(g)>0$.
Thus $h$ is ordered. Applying strong additivity to each normal form gives
\eqref{isc:eq:lift}. Uniqueness is immediate from the monomials.
\end{proof}

\wtag{} \Cref{isc:thm:lift,isc:prop:arithmeticlift,isc:prop:classification}
are the case $\rho=\id$ of the bottom-gap classification
\lbl{opa:as:thm:classification} of the omnific-preserving report, written
after the pin. Without the real-fixing hypothesis the exact-monomial
embeddings preserving $\Oz$ are the maps $J_{h,\rho}$, which may move
reals (\cref{isc:sec:collection}).

\begin{warning}[Preserving monomials is not preserving the omega-map]\label[warning]{isc:warn:omega}
The equality $\Phi_h(\omega^g)=\omega^{h(g)}$ does not say
$\Phi_h(\omega^g)=\omega^{\Phi_h(g)}$. In general $h\ne\Phi_h$ as maps
on surreal numbers. Nor have we shown compatibility with the Gonshor
exponential, the Berarducci--Mantova derivation, birthdays, or simplicity.
These extra structures are outside the assertions of this article.
\end{warning}

\section{Coset slices and linear disjointness}\label{isc:sec:slices}

The next argument upgrades an exponent-group intersection calculation to
linear disjointness of \emph{full} series fields. The subgroups need not be
convex, and no quotient ordering is needed.

\begin{definition}[Normalized coset slice]
Let $H\subseteq G$ be additive subgroups and $g\in G$. For
$f=\sum_e a_e\omega^e\in\K{k}{G}$ set
\begin{equation}\label{isc:eq:slice}
 \lambda_{g,H}(f)
 =\omega^{-g}\sum_{e\in(g+H)\cap\supp(f)}a_e\omega^e
 \ \in\K{k}{H}.
\end{equation}
The value depends on the chosen representative $g$ by the expected monomial
normalization. We do not choose a global transversal of $G/H$.
\end{definition}

\begin{lemma}[Slice linearity and separation]\label[lemma]{isc:lem:slice}
Every $\lambda_{g,H}$ is $\K{k}{H}$-linear, and the family of slices
separates points of $\K{k}{G}$.
\end{lemma}
\begin{proof}
A subset of a reverse-well-ordered support is reverse well ordered, and
translation preserves that property. Thus the definition is valid even
when $H$ is not convex. If $a\in\K{k}{H}$, a product term with exponent
$e+h$ belongs to $g+H$ exactly when $e$ belongs to $g+H$. Therefore slicing
the product $af$ gives $a$ times the slice of $f$. The equality is verified
coefficientwise in finite convolution sums. Additivity follows directly.
If $f\ne0$, choose any $g$ in its support. Its slice contains the nonzero
coefficient $a_g$ at exponent zero, so the slice is nonzero.
\end{proof}

\begin{lemma}[A finite separating matrix]\label[lemma]{isc:lem:matrix}
If $f_1,\ldots,f_m\in\K{k}{G}$ are linearly independent over
$L=\K{k}{H}$, there are $g_1,\ldots,g_m\in G$ such that
\[
 A=\bigl(\lambda_{g_a,H}(f_b)\bigr)_{1\le a,b\le m}
 \quad\text{is invertible over }L.
\]
\end{lemma}
\begin{proof}
On the $m$-dimensional space $V=\operatorname{span}_L(f_1,\ldots,f_m)$,
the restrictions of the slices have common kernel zero by
\cref{isc:lem:slice}. Start with $V$. If the current common kernel is nonzero,
choose a nonzero vector in it and a slice not vanishing on that vector.
Intersecting with this new kernel strictly lowers dimension. After at most
$m$ steps the common kernel is zero. The resulting linear map from $V$ to
at most $m$ copies of $L$ is injective, so exactly $m$ independent coordinate
functionals can be used, and their matrix in the stated basis is invertible.
This is a finite argument, also for class-sized $L$.
\end{proof}

\begin{theorem}[Full Hahn fields are disjoint over their full intersection]
\label{isc:thm:sliceddisjoint}
Let $G_1,G_2$ be additive subgroups of one ordered abelian group and put
$H=G_1\cap G_2$. Inside $\K{k}{G_1+G_2}$,
\[
 \K{k}{G_1}\cap\K{k}{G_2}=\K{k}{H},
\]
and the two fields are linearly disjoint over $\K{k}{H}$.
The assertion applies to set or class subgroups, with set-sized individual
supports, and to any coefficient field $k$.
\end{theorem}
\begin{proof}
The intersection identity follows from uniqueness of coefficients and
support. Put $L=\K{k}{H}$ and let $f_1,\ldots,f_m\in\K{k}{G_1}$ be
$L$-linearly independent. Choose the invertible slice matrix of
\cref{isc:lem:matrix} with representatives $g_a\in G_1$.

On the larger field $\K{k}{G_1+G_2}$ take slices
$\lambda_{g_a,G_2}$ with values in $\K{k}{G_2}$. They are
$\K{k}{G_2}$-linear. For $f\in\K{k}{G_1}$, the identity
\[
 G_1\cap(g_a+G_2)=g_a+(G_1\cap G_2)=g_a+H
\]
shows that $\lambda_{g_a,G_2}(f)=\lambda_{g_a,H}(f)$.
Suppose $\sum_b b_b f_b=0$ with $b_b\in\K{k}{G_2}$. Applying all of
these enlarged slices gives $A(b_1,\ldots,b_m)^{\mathsf T}=0$. Since $A$
is invertible already over $L$, all $b_b$ vanish. This is precisely linear
disjointness.
\end{proof}

\begin{corollary}[Independence of any finite number of full Hahn fields]
\label[corollary]{isc:cor:joint}
If $(G_i)_{i\in I}$ are independent over $H$, then the full fields
$\K{k}{G_i}$ are jointly linearly disjoint over $\K{k}{H}$.
For $i\notin J$, the field $\K{k}{G_i}$ is linearly disjoint over
$\K{k}{H}$ even from the full field
$\K{k}{H+\sum_{j\in J}G_j}$.
\end{corollary}
\begin{proof}
The second assertion follows from \cref{isc:thm:sliceddisjoint} and the
corresponding group intersection. For a finite family, induct on its size.
The compositum of the preceding fields is contained in the full Hahn field
on the sum of their groups, and the next field is linearly disjoint from
that larger field. Extension of scalars over a field preserves an injection
of vector spaces. Applying this to the already injective finite-tensor
multiplication map proves injectivity for the enlarged family. Equivalently,
choose a basis of the finite span of the preceding products and apply linear
disjointness to that basis.
\end{proof}

\begin{corollary}[Polynomial-relation descent]\label[corollary]{isc:cor:relations}
Let $P=\K{k}{G_1}$, $Q=\K{k}{G_2}$, and $L=\K{k}{G_1\cap G_2}$.
For a finite tuple $a$ from $P$, let $I_L(a)\subseteq L[X_1,\ldots,X_n]$
be its polynomial relation ideal. Then
\[
 I_Q(a)=I_L(a)\,Q[X_1,\ldots,X_n].
\]
\end{corollary}
\begin{proof}
Take a polynomial relation $F(a)=0$ with coefficients in $Q$. A finite
$L$-basis $b_1,\ldots,b_r$ for the span of its coefficients gives
$F=\sum_j b_jF_j$ with $F_j\in L[X_1,\ldots,X_n]$. Linear disjointness
is symmetric (as is the finite-tensor criterion), so the $b_j$ remain
independent over $P$. It follows that every $F_j(a)=0$. Thus $F$ belongs
to the extended ideal. The other inclusion is immediate.
\end{proof}

\section{Prescribed intersections and Boolean meet diagrams}\label{isc:sec:copies}

\begin{proof}[Proof of \cref{isc:main:copies}]
Apply \cref{isc:thm:group} to obtain $(h_\alpha)$ and their independent ranges
$G_\alpha$, and set $\Phi_\alpha=\Phi_{h_\alpha}$. By
\cref{isc:thm:lift}, the image is $K_\alpha=\K{\R}{G_\alpha}$.
Since $h_\alpha$ is the identity on $H$, it fixes every monomial with
exponent in $H$, and strong additivity makes $\Phi_\alpha$ the identity
on all of $L_\R$. Each $G_\alpha$ is proper by
\cref{isc:cor:groupintersection}; a monomial with exponent outside it witnesses
properness of the corresponding field.

The intersection and independence assertions follow from
\cref{isc:thm:sliceddisjoint,isc:cor:joint}. \Cref{isc:prop:arithmeticlift} gives the
omnific assertions. The same groups and the same slice argument with
coefficients in $\C$ give the entire complex package.
\end{proof}

\begin{corollary}[Fixing every element of a prescribed set]\label[corollary]{isc:cor:parameters}
For every set $A\subseteq\No$ there is a set-sized real closed full Hahn
field $L\supseteq\R\cup A$ which is the exact pairwise intersection of
an ordinal-indexed family of independent strong omnific-preserving copies
of $\No$ fixing $L$ pointwise. For a set $A\subseteq\SC$, the analogous
statement holds with an algebraically closed $L\supseteq\C\cup A$.
\end{corollary}
\begin{proof}
Take $H$ to be the $\Q$-span of the union of the normal-form supports of
all elements of $A$. In the complex case take the real and imaginary
supports, equivalently the complex support. Replacement and union make this
a set, and rational span remains a set. Then $A\subseteq\K{k}{H}$.
Apply \cref{isc:main:copies}.
\end{proof}

This corollary does \emph{not} assert that the intersection can be prescribed
to be the ordinary field generated by $A$, its real closure, or an arbitrary
non-Hahn subfield. It prescribes the full support closure just constructed.
The distinction is essential to both the proof and the classification below.
\wtag{} For two maps, and without intersection statements, proper strong
exact-monomial copies fixing $A\cup\R$ and preserving and reflecting $\Oz$
are \lbl{opa:as:thm:parameters} of the omnific-preserving report; the
corollary extends that result to an $\On$-indexed family with an exact
common field (\cref{isc:sec:collection}).
\mtag{} The same support-closure step localizes a set of parameters in two
later results. In the dilation report, \lbl{adr:sr:thm:tail} places a set of
surcomplex parameters in the Hahn field of the real span of their supports
and finds a tail of an explicit ordinal family whose positive real dilation
orbits are algebraically independent over that whole field, and
\lbl{adr:sr:thm:kappa} gives set-sized families of independent orbits over
any set-sized base. In this report, \cref{isc:cp:lem:localizeP} applies the
step one normal-form level deeper, to the outer coefficients of a parameter
set, and \cref{isc:cp:thm:resilience} keeps a tail of a fixed ordinal family
differentially independent over the compositum with any set of parameters
adjoined.

\begin{theorem}[Exact classification of set-sized intersections in this class]
\label{isc:thm:intersectionclassification}
The set-sized fields that occur as intersections of two images of
strong, real-fixing, monomial-preserving self-embeddings of $\No$ are
exactly the fields $\K{\R}{H}$ with $H\subseteq\No$ a set-sized divisible
additive subgroup.
\end{theorem}
\begin{proof}
For necessity, \cref{isc:prop:classification} writes the two images as full
Hahn fields on divisible groups $G_1,G_2$. Their intersection is
$\K{\R}{G_1\cap G_2}$. In a torsion-free divisible ambient group, each
divisible subgroup is closed under the uniquely defined rational division;
therefore their intersection $H$ is divisible. If the intersection field
is a set, its monomials form a set and their exponents, recovered by the
inverse of the omega-map on monomials, form a set. Thus $H$ is set-sized.
Sufficiency is \cref{isc:main:copies}.
\end{proof}

\wtag{} The real-fixing hypothesis is used through
\cref{isc:prop:classification}. Exact-monomial embeddings that preserve
$\Oz$ but move reals exist whenever the bottom gap is nonzero
(\lbl{opa:as:thm:taylorcriterion}); their images meet $\R$ in a proper
real closed subfield (\lbl{opa:as:prop:overlap}). The theorem says nothing
about intersections in that larger class.

\subsection{All intersections in an entire diagram}

Fix a set of indices $I$ from the family of \cref{isc:main:copies}. For
$J\subseteq I$ put
\begin{equation}\label{isc:eq:diagram}
 G_J=H+\sum_{j\in J}G_j,\qquad E_J=\K{k}{G_J},
 \qquad G_\varnothing=H.
\end{equation}
Every sum in $G_J$ is finite; no infinite additive closure is intended.

\begin{theorem}[Boolean meet realization]\label{isc:thm:boolean}
The assignment $J\mapsto E_J$ is injective and order preserving. For
$J,T\subseteq I$,
\[
 E_J\cap E_T=E_{J\cap T}.
\]
The analogous identity holds for every nonempty set-indexed intersection of
subsets of $I$. Moreover $E_J$ and $E_T$ are linearly disjoint over
$E_{J\cap T}$. Their full Hahn join is $E_{J\cup T}$. All $E_J$ are real
closed for $k=\R$ and algebraically closed for $k=\C$.
\end{theorem}
\begin{proof}
Independence over $H$ means that an element of $G_I/H$ has a unique finite
set of nonzero components in the direct sum of the groups $G_i/H$.
It lies in $G_J/H$ exactly when all of these components have indices in
$J$. This proves all the group intersection identities and
$G_J+G_T=G_{J\cup T}$. Uniqueness of normal form transfers the
intersection identities to full Hahn fields, and
\cref{isc:thm:sliceddisjoint} proves the asserted linear disjointness.

If $j\in J\setminus T$, choose $g\in G_j\setminus H$; then
$\omega^g\in E_J\setminus E_T$. This proves injectivity and reflection
of containment. The full Hahn join is, by definition, the full Hahn field
on the sum of the two exponent groups, hence is $E_{J\cup T}$.
Every $G_J$ is divisible, which gives the closedness assertions.
\end{proof}

\begin{remark}[The word ``join'']\label[remark]{isc:rem:join}
The theorem realizes a Boolean meet diagram with \emph{Hahn} joins. It does
not identify ordinary field composita with those joins. \Cref{isc:main:join}
shows that the difference can already be transcendental for two factors.
\end{remark}

\begin{corollary}[Independent choices and absence of new relations]
\label[corollary]{isc:cor:choices}
For distinct indices $i_1,\ldots,i_m$, choose
$a_j\in K_{i_j}\setminus L_\R$. Then $(a_1,\ldots,a_m)$ is algebraically
independent over $L_\R$. The same assertion holds over $L_\C$ for the
complexified copies. More generally all polynomial relations across any
square of \cref{isc:thm:boolean} descend to its intersection field.
\end{corollary}
\begin{proof}
A real closed subfield has no proper algebraic extension inside an ordered
field, so each $a_j$ is transcendental over $L_\R$. In the complex case
use algebraic closedness. The polynomial algebras $L[a_j]$ therefore have
no relations individually, and their finite tensor product is a polynomial
algebra on $m$ variables. Joint linear disjointness makes its multiplication
map injective. The last statement is \cref{isc:cor:relations}.
\end{proof}

\subsection{A limited model-theoretic conclusion}\label{isc:sub:model}

For each fixed first-order formula in the language of ordered fields, the
inclusion of any $E_J$ into $\No$ is elementary, by quantifier elimination
for real closed fields. The corresponding complex statement follows from
quantifier elimination for algebraically closed fields of characteristic
zero~\cite{Marker}. This is a formula-by-formula use of ordinary field theory, not the
construction of a class satisfaction predicate.

Nothing here asserts elementary embedding in the expanded language with a
predicate for $\Oz$, a symbol for the omega-map, or a derivation. The
independence proved is algebraic independence and linear disjointness; it
is not being silently identified with a model-theoretic independence notion
in any richer theory.

\section{Explicit copies with intersection exactly the constants}\label{isc:sec:explicit}

The general construction can fix an arbitrary set-sized Hahn core, but uses
class bookkeeping. When $H=0$ there is a direct formula. It also furnishes
the separated scales needed for \cref{isc:main:join}.

\begin{lemma}[Compression of the surreal order]\label[lemma]{isc:lem:compression}
The function
\[
 \psi(x)=\frac12+\frac{x}{2(1+|x|)}
\]
is an order isomorphism from $\No$ onto the open interval $(0,1)$.
\end{lemma}
\begin{proof}
On the nonnegative half-line, $x/(1+x)$ is strictly increasing and belongs
to $[0,1)$. On the nonpositive half-line, $x/(1-x)$ is strictly increasing
and belongs to $(-1,0]$. The formulas agree at zero, proving strict
monotonicity and the stated range containment. For $u\in(0,1)$ put
$s=2u-1$ and $x=s/(1-|s|)$. Direct substitution on the two signs of $s$
gives $\psi(x)=u$. All denominators are positive in the ordered field.
\end{proof}

For every ordinal $\alpha$ let
\[
 I_\alpha=(3\alpha,3\alpha+1),\qquad
 \sigma_\alpha(a)=3\alpha+\psi(a),
\]
where the arithmetic in $3\alpha$ is surreal field arithmetic. These
intervals are pairwise disjoint. Indeed $\alpha<\beta$ for ordinals implies
$\alpha+1\le\beta$ in the surreal order, hence $3\alpha+1<3\beta$.

Define an \emph{inner} map on normal forms by
\begin{equation}\label{isc:eq:inner}
 h_\alpha\left(\sum_a r_a\omega^a\right)
   =\sum_a r_a\omega^{\sigma_\alpha(a)}.
\end{equation}
This map is an ordered, real-linear, strongly additive \emph{group}
embedding. Additivity only uses the injective relabeling of coefficients;
no additivity of $\sigma_\alpha$ is needed. It is not a unital field map:
\begin{equation}\label{isc:eq:innerone}
 h_\alpha(1)=\omega^{3\alpha+1/2}.
\end{equation}
Its image is precisely the additive class of all surreal normal forms
supported in $I_\alpha$. Pulling back the support along the order
isomorphism $\sigma_\alpha$ proves fullness of that image.

Now apply the \emph{outer} construction of \cref{isc:thm:lift}:
\begin{equation}\label{isc:eq:explicitouter}
 \Phi_\alpha\left(\sum_g c_g\omega^g\right)
   =\sum_g c_g\omega^{h_\alpha(g)}.
\end{equation}
The distinction between \eqref{isc:eq:inner} and \eqref{isc:eq:explicitouter} is
crucial: the outer exponents $g$ are themselves surreal numbers with their
own normal forms. \wtag{} The inner map is the inner lift
$\iota_{\sigma_\alpha}$ of the omnific-preserving report
(\lbl{opa:as:eq:innerlift}, \lbl{opa:as:lem:inner}), and
\eqref{isc:eq:explicitouter} is its monomial map $M_{1,\iota_{\sigma_\alpha}}$.
That report's order isomorphisms $\mu_{11},\mu_{12},\mu_{13}$ have unbounded
images; the bounded, pairwise disjoint images of the $\sigma_\alpha$ are
what makes the ranges independent here.

\begin{theorem}[An explicit independent ordinal family]\label{isc:thm:explicit}
The maps \eqref{isc:eq:explicitouter} are strong, real-fixing,
monomial-preserving proper field embeddings. Their image fields are jointly
linearly disjoint over $\R$, have pairwise intersection $\R$, and their
omnific image rings have pairwise intersection $\Z$. Their complexifications
have the corresponding intersections $\C$ and $\Z[i]$.
\end{theorem}
\begin{proof}
The inner ranges $G_\alpha$ have disjoint normal-form support intervals.
In a finite sum of elements from different ranges, no nonzero coefficient
from one interval can be canceled by another interval. Therefore these
ranges are independent over zero. Apply
\cref{isc:thm:lift,isc:prop:arithmeticlift,isc:cor:joint}.
\end{proof}

For example,
\[
 \Phi_0(\omega)=\omega^{\omega^{1/2}},\qquad
 \Phi_1(\omega)=\omega^{\omega^{7/2}}.
\]
Both maps fix every real number. Since $\Phi_0(1)=1$, the first identity
also explicitly shows
\[
 \Phi_0(\omega^1)\ne\omega^{\Phi_0(1)}.
\]
Thus preservation of the set of Conway monomials is much weaker than
commutation with the omega-map as a unary operation on all of $\No$.

\begin{remark}[A tempting formula that does not give a proper copy]\label[remark]{isc:rem:scaling}
Replacing every exponent $g$ by $cg$ with a fixed $c>0$ gives an
automorphism, not a proper self-embedding, because $c\No=\No$. The
compression above is applied one normal-form level below the field lift;
that is what makes the exponent-group image genuinely proper.
\end{remark}

\section{A transcendental gap above an ordinary compositum}\label{isc:sec:join}

Take $G_0,G_1$ from \cref{isc:sec:explicit}. These groups are not merely
independent. Their scales are separated:
\begin{equation}\label{isc:eq:separated}
 |a|<b/n\quad
 (a\in G_0,\ 0<b\in G_1,\ n\in\N_{>0}).
\end{equation}
Indeed the inner leading exponent of a nonzero $a$ lies in $(0,1)$,
whereas that of $b$ lies in $(3,4)$. Leading-monomial comparison proves
\eqref{isc:eq:separated}. Thus $G_0$ is convex in the ordered group
$G_0+G_1$, and the sign of $a+b$ for $b\ne0$ is the sign of $b$.

\subsection{Iterated Hahn coordinates}

\begin{lemma}[Re-expansion at separated scales]\label[lemma]{isc:lem:iterated}
Suppose $A,B$ are additive subgroups with $A\cap B=0$ and every element
of $A$ is infinitesimal relative to every positive element of $B$, in the
sense of \eqref{isc:eq:separated}. Then regrouping by $B$-coordinate gives a
field isomorphism
\begin{equation}\label{isc:eq:iterated}
 \K{k}{A+B}\ \simeq\ \K{\K{k}{A}}{B},
 \qquad
 \sum_{a+b}c_{a,b}\omega^{a+b}
 \longmapsto\sum_b\left(\sum_a c_{a,b}\omega^a\right)\omega^b.
\end{equation}
On the right the notation means iterated Hahn series; the coefficient
field may be a class, but each individual series uses only a set of
coefficients. The isomorphism is order preserving for ordered $k$.
\end{lemma}
\begin{proof}
The decomposition $a+b$ is unique, and the order on $A+B$ is lexicographic
with $B$ the dominant coordinate. The image in $B$ of a reverse-well-ordered
support is reverse well ordered: the inverse image of a nonempty subset of
that image has a greatest element, whose projection is greatest in the
subset. Each fiber is reverse well ordered in $A$.

Conversely, a reverse-well-ordered set of $B$-coordinates, with a
reverse-well-ordered support in each fiber, yields a reverse-well-ordered
lexicographic union. Its size is a set by replacement and union. The two
regroupings are inverse. Addition and multiplication agree coefficientwise:
there are finitely many pairs of $B$-coordinates contributing to each
outer product coefficient, and within each resulting coefficient Hahn
multiplication in $A$ applies. The dominant $B$-coordinate and then the
leading coefficient in $\K{k}{A}$ also give exactly the original order.
\end{proof}

For $k=\R$ or $\C$, put
\[
 K_0=\K{k}{G_0},\quad K_1=\K{k}{G_1},\quad
 E=\K{k}{G_0+G_1},\quad M=K_0K_1.
\]
In the iterated coordinates, $K_0$ is the coefficient field and $K_1$ has
coefficients in the original $k$.

\begin{lemma}[Finite coefficient-field control]\label[lemma]{isc:lem:finitecoeff}
Every finite subset of the ordinary compositum $M$ is contained, in the
iterated coordinates, in a field $\K{D}{G_1}$ where $D\subseteq K_0$ is
a field finitely generated over $k$.
\end{lemma}
\begin{proof}
Each of the finitely many chosen elements is obtained by finitely many
field operations from finitely many elements of $K_0$ and $K_1$. Let
$a_1,\ldots,a_r\in K_0$ include all the first-kind elements used, and put
$D=k(a_1,\ldots,a_r)$. The elements of $K_1$ have coefficients in
$k\subseteq D$, while each $a_j$ is a constant in the outer $G_1$-series
coordinate. The field $\K{D}{G_1}$ is closed under the required field
operations, proving the assertion.
\end{proof}

This lemma does not claim that every series with coefficients in such a
$D$ is in $M$. It gives a necessary coefficient-field constraint, which is
all that is needed.
\mtag{} The constraint is far from sufficient: a series whose coefficients
all lie in a field $\R(x)$ generated by one element can be differentially
transcendental over $M$ (\cref{isc:cp:cor:fginsufficient}). The exact
description of $M$ in these coordinates is the finite coefficient-rank
denominator test of \cref{isc:cp:cor:surrealrank}.

\begin{theorem}[A coefficient-transcendence obstruction]\label{isc:thm:joinobstruction}
Let $s_1,s_2,\ldots\in K_0$ be algebraically independent over $k$, where
$k=\R$ or $\C$, and choose $0<g\in G_1$. The well-defined series
\[
 S=\sum_{n=1}^{\infty}s_n\omega^{-ng}\in E
\]
is transcendental over $M$.
\end{theorem}
\begin{proof}
The outer support $\{-ng:n\ge1\}$ is strictly decreasing, so the series
is valid by \cref{isc:lem:iterated}. Suppose a nonzero polynomial over $M$
annihilates $S$. By \cref{isc:lem:finitecoeff}, all its coefficients lie in
$\K{D}{G_1}$ for one finitely generated field $D\subseteq K_0$ over $k$.

For $k=\R$, let $D'$ be the relative algebraic closure of $D$ in the real
closed field $K_0$. It is a set-sized real closed field. To see closedness,
a square root of a positive element of $D'$ and a root of an odd-degree
polynomial over $D'$ can be taken in $K_0$; each is algebraic over $D$,
so belongs to $D'$. It is a set because a set of polynomials has only a
set of finite root sets. For $k=\C$, the same argument with arbitrary
polynomial roots gives a set-sized algebraically closed relative closure
$D'\subseteq K_0$.

Since $G_1$ is divisible, the iterated Hahn field $\K{D'}{G_1}$ is real
closed in the real case, and algebraically closed in the complex case. For
a proper-class $G_1$ this follows, as before, by working with the set-sized
divisible subgroup generated by any finite collection of coefficient
supports. Thus an element of $E$ algebraic over $\K{D}{G_1}$ must belong
to $\K{D'}{G_1}$: use relative algebraic closedness in an ordered field
in the real case, or algebraic closedness in the complex case.

Applied to $S$, uniqueness of its outer expansion now gives $s_n\in D'$
for every $n$. But
\[
 \trdeg_kD'=\trdeg_kD<\infty,
\]
whereas the $s_n$ are an infinite algebraically independent family over
$k$. This contradiction proves the theorem.
\end{proof}

\begin{proof}[Explicit witness and proof of \cref{isc:main:join}]
Take
\[
 g_n=h_0(\omega^n)=\omega^{\psi(n)},\qquad
 g=h_1(1)=\omega^{7/2},\qquad s_n=\omega^{g_n}\quad(n\ge1).
\]
The $g_n$ are $\Q$-linearly independent: they are distinct inner Conway
monomials, so uniqueness of normal form precludes a finite rational linear
relation. Consequently the $s_n$ are algebraically independent over
$\R$, and also over $\C$ in the complexified field. Indeed the distinct
monomials of a polynomial in the $s_n$ acquire distinct exponents
$\sum_n m_ng_n$, and no nonzero coefficient can cancel.

The promised actual normal form is
\begin{equation}\label{isc:eq:explicitS}
 \boxed{\displaystyle
 S=\sum_{n=1}^{\infty}
   \omega^{\,\omega^{\psi(n)}-n\omega^{7/2}}.}
\end{equation}
Its exponents strictly decrease: the difference between consecutive
exponents is $g+g_n-g_{n+1}>0$ by scale separation. Thus
\eqref{isc:eq:explicitS} is a legitimate countably supported surreal, not a
formal expression awaiting an analytic interpretation. It belongs to $E$
and is transcendental over the appropriate ordinary compositum by
\cref{isc:thm:joinobstruction}. Transcendence also excludes membership in its
relative algebraic closure.
\end{proof}

\mtag{} The second manuscript extends this witness. With
$b=\omega^{15/4}\in G_1$, the element $\omega^bS$ is the first member $Y_0$ of
an ordinal family $Y_\alpha$ of purely infinite omnific integers of order
type $\omega$ (\cref{isc:cp:thm:omnific}); all their Euler jets are
algebraically independent over $M$, and a tail of the family stays
independent after any set of parameters is adjoined
(\cref{isc:cp:thm:resilience}). In the real case $S$ itself is then
differentially transcendental over $M$ for that Euler derivation.
\mtag{} A third manuscript gives a second proof that $S$ is transcendental
over $M$, real and complex, by a mixed-support annihilator that needs no
iterated coordinates (\cref{isc:hj:rem:secondroute}).

\begin{remark}[Why the infinite list of coefficients matters]\label[remark]{isc:rem:infinite}
A rational expression in the two fields can use arbitrary elements of each
field, but only finitely many of them. In the separated-scale expansion this
bounds the transcendence degree of its coefficient field over the constants.
Algebraic extension does not remove that finite bound. The series
\eqref{isc:eq:explicitS} defeats it by displaying infinitely many independent
coefficients. This argument is different from simply observing that a series
is an infinite sum; many infinite sums, such as a geometric series, already
belong to an ordinary compositum.
\end{remark}

\section{Omnific intersections and their fraction fields}\label{isc:sec:fractions}

For $k=\R,\C$ write $D_\R=\Z$ and $D_\C=\Z[i]$, and define
\[
 \PiG{k}{H}=\{x\in\K{k}{H}:\supp(x)\subseteq H^{>0}\},\qquad
 \A{k}{H}=D_k\oplus\PiG{k}{H}.
\]
The direct sum is additive; multiplication makes $\PiG{k}{H}$ a nonunital
ring and an ideal in $\A{k}{H}$. In particular,
\[
 \A{\R}{H}=\K{\R}{H}\cap\Oz,\qquad
 \A{\C}{H}=\K{\C}{H}\cap\GOz.
\]
We use the ordinary class fraction field of a class domain to mean the
subfield of the ambient field consisting of ratios of two of its elements,
with nonzero denominator. No quotient set of a proper class is invoked.

\begin{lemma}[Common denominators for a set of surreals]\label[lemma]{isc:lem:denominator}
Every set $X\subseteq\No$ has a positive monomial $d\in\Oz$ such that
$dX\subseteq\Oz$. The same holds for a set of surcomplex numbers with
$\GOz$ in place of $\Oz$, and a positive real monomial denominator.
In particular
\[
 \Frac(\Oz)=\No,\qquad\Frac(\GOz)=\SC.
\]
\end{lemma}
\begin{proof}
The union $S$ of the supports of elements of $X$ is a set. Choose a
surreal $a>0$ with $a>-s$ for every $s\in S$, using the cut property.
Then $d=\omega^a$ and every product $dx$ have exclusively positive
exponents, except that $dx$ may be zero. Thus they lie in the required
integer ring. In the complex case use the union of the real and imaginary
supports. Applying the assertion to a singleton gives the fraction-field
identities. These identities are classical; compare \cite[Proposition~2.4.5]{LM}.
\end{proof}

\wtag{} In the collection this is the global monomial clearing
\lbl{odg:thm:fractions} of the Diophantine report, with the same proof.

Transporting the lemma through an embedding gives more than merely one
denominator for each individual element.

\begin{proposition}[Separate uniform denominators, no shared uniform one]
\label[proposition]{isc:prop:uniform}
Let $K_\alpha,R_\alpha,L_\R$ be as in \cref{isc:main:copies}. Then
$\Frac(R_\alpha)=K_\alpha$, and there is a nonzero $d_\alpha\in R_\alpha$
with
\[
 d_\alpha L_\R\subseteq R_\alpha.
\]
For two distinct indices, no nonzero element of $R_\alpha\cap R_\beta$
can satisfy both such containments. The same conclusions hold in the
complex and Gaussian setting.
\end{proposition}
\begin{proof}
The image under $\Phi_\alpha$ of a ratio of omnific integers is a ratio
of elements of $R_\alpha$, and every element of $K_\alpha$ arises this
way by \cref{isc:lem:denominator}. Since $L_\R$ is a set and is fixed
pointwise, apply that lemma to $L_\R$ and transport its denominator by
$\Phi_\alpha$. This yields the uniform denominator $d_\alpha$.

A hypothetical common nonzero denominator $d$ belongs to
$R_\alpha\cap R_\beta\subseteq L_\R$. If both containments held, then
$dL_\R\subseteq\A{\R}{H}$. But $d\ne0$ is in the field $L_\R$, so
$dL_\R=L_\R$, which is not contained in its omnific part (it contains
$1/2$). This is impossible. The same proof uses $1/2\notin\Z[i]$ in the
complex case.
\end{proof}

\wtag{} The proof shows slightly more than the statement: since
$dL_\R=L_\R\not\subseteq\Oz$, a nonzero element of $R_\alpha\cap R_\beta$
satisfies \emph{neither} containment, not merely not both. The separate
denominators $d_\alpha$ lie outside $L_\R$.

This proposition already separates two notions of a denominator. To prove
\cref{isc:main:fractions}, however, we need an element with \emph{no shared
denominator at all}, not merely the absence of a shared denominator uniform
on the entire common field. That stronger assertion is proved next.

\subsection{The bounded-support field is exactly the shared fraction field}\label{isc:sub:bounded}

For this subsection and \cref{isc:sec:gaps}, use the increasing-support Hahn
convention
\[
 t^h=\omega^{-h}.
\]
As $H=-H$, this changes the presentation of $\K{k}{H}$, not its underlying
field. The support of a nonzero $t$-series is well ordered and has a least
element, its valuation $v$. Set $v(0)=+\infty$.
Let
\[
 \B{k}{H}=\{f\in k((t^H)):\supp_t(f)\text{ is bounded above in }H\}.
\]
It is a subring: finite unions of upper bounds bound sums, and the sum of
two upper bounds bounds a product. Its elements can have infinite support
inside a bounded interval. In this convention $\A{k}{H}$ consists of
series with negative exponents and a constant coefficient in $D_k$.

\begin{proposition}[Identification of the fraction fields]\label[proposition]{isc:prop:boundedfrac}
If $H\ne0$, then
\begin{equation}\label{isc:eq:boundedfrac}
 \Frac\A{k}{H}=\Frac\B{k}{H}.
\end{equation}
In particular this fraction field contains $k$ and every single monomial
$t^h$ with $h\in H$.
\end{proposition}
\begin{proof}
One has $\A{k}{H}\subseteq\B{k}{H}$. Conversely, take
$f\in\B{k}{H}$ and choose an upper bound $u\in H$ of its support.
Choose $a\in H$ with $a>u$ and $a>0$, possible in any nonzero ordered
abelian group. Both $t^{-a}$ and $t^{-a}f$ have strictly negative support,
so belong to $\A{k}{H}$. Their ratio is $f$, giving the reverse inclusion
of fraction fields. Constants and monomials have bounded support.
\end{proof}

The identity explains why a mere missing constant is not the obstruction
for nonzero $H$. For example, $r=(r\omega^h)/\omega^h$ for any $r\in\R$
and any positive $h\in H$. The unresolved issue is the unboundedness of
supports \emph{within the fixed set group $H$}.
\wtag{} \Cref{isc:prop:boundedfrac} is the localization of
\lbl{bst:prop:localdensity} of the bounded-support report, with the
constant coefficient restricted to $D_k$. Since batch~32 that report also
prints the general form, for every field $K$ and every unital subring
$D\subseteq K$, as \lbl{bst:gr:prop:fractions}.

\section{A self-contained cofinal gap argument}\label{isc:sec:gaps}

The mechanism in this section belongs to bounded-support Hahn arithmetic,
not to the new embedding construction. It is included in full so that the
transcendental conclusion in \cref{isc:main:fractions} does not depend on an
unreviewed repository result. The repository's bounded-support report gives
stronger independent-family assertions \cite{RepoBounded}, and the
non-fraction obstruction is already in \cite{LM}.
\wtag{} In that report, \cref{isc:thm:cofinalgap} is \lbl{bst:cor:onegap}
at countable cofinality and, in characteristic zero, the one-series case of
\lbl{bst:thm:independent}; the single-series form at uncountable cofinality
over an arbitrary coefficient field is not stated there, but its proof is
that report's tail estimate and root count (\cref{isc:sec:collection}). No
priority is claimed for it.

\begin{lemma}[Cofinal sequences with polynomial gaps]\label[lemma]{isc:lem:gapscale}
Let $H\ne0$ be a set-sized ordered abelian group and put $\kappa=\cf(H)$.
There is a strictly increasing cofinal family $(a_\alpha)_{\alpha<\kappa}$
in $H^{>0}$ such that for every positive integer $d$ and every $c\in H$,
\begin{equation}\label{isc:eq:gapscale}
 a_{\alpha+1}-d a_\alpha>c
 \quad\text{for all sufficiently large }\alpha<\kappa.
\end{equation}
\end{lemma}
\begin{proof}
The cofinality of a nonzero ordered group is an infinite regular cardinal:
a finite subset is bounded, and a cofinal union of fewer than $\kappa$
subsets of size less than $\kappa$ would contradict minimality of
$\kappa$ by taking bounds for those subsets. Fix a cofinal list
$(u_\alpha)_{\alpha<\kappa}$.

When $\kappa=\aleph_0$, choose recursively $a_0>0,u_0$ and
\[
 a_{n+1}>u_{n+1},\quad a_{n+1}>(n+2)a_n.
\]
The list is cofinal. For fixed $d$, once $n+2\ge d+1$ the difference in
\eqref{isc:eq:gapscale} exceeds $a_n$, and this eventually exceeds $c$.

When $\kappa>\aleph_0$, at stage $\alpha$ the set
\[
 \{0,u_\alpha\}\cup
 \{n a_\beta:\beta<\alpha,\ n\in\N_{>0}\}
\]
has cardinality less than $\kappa$ and hence is bounded above in $H$.
Choose $a_\alpha$ strictly above it. The resulting family is increasing and
cofinal. At a successor stage,
$a_{\alpha+1}>(d+1)a_\alpha$ for every positive integer $d$. Thus again
the difference in \eqref{isc:eq:gapscale} exceeds $a_\alpha$ and is eventually
larger than $c$. Since an infinite cardinal has no last ordinal below it,
all the successor indices used are still below $\kappa$.
\end{proof}

\begin{theorem}[One transcendental series over all bounded-support fractions]
\label{isc:thm:cofinalgap}
For any field $k$ and nonzero set-sized ordered abelian group $H$, choose
a sequence from \cref{isc:lem:gapscale}. Then
\begin{equation}\label{isc:eq:cofinalF}
 F=1+\sum_{\alpha<\kappa}t^{a_\alpha}
\end{equation}
is transcendental over $\Frac\B{k}{H}$.
\end{theorem}
\begin{proof}
The support is a well-ordered set, so $F$ is a Hahn series and $v(F)=0$.
Suppose a nonzero polynomial over $\Frac\B{k}{H}$ annihilates $F$.
Clearing its finitely many denominators gives
\[
 Q(Y)=\sum_{j=0}^d b_jY^j\in\B{k}{H}[Y]\setminus\{0\},\qquad Q(F)=0,
\]
where $d\ge1$. Each nonzero coefficient has support bounded both below
and above: below by its least exponent and above by the definition of
$\B{k}{H}$. Choose common bounds $b,c\in H$ for the supports of all
nonzero coefficients, with $b\le c$.

For $\alpha<\kappa$ let
\[
 P_\alpha=1+\sum_{\beta\le\alpha}t^{a_\beta}.
\]
These are distinct elements of the same Hahn field, including at limit
indices; they need not be polynomials in $t$. Their supports lie in
$[0,a_\alpha]$, and
\[
 v(P_\alpha)=0,\qquad v(F-P_\alpha)=a_{\alpha+1}.
\]
Finite Hahn multiplication gives
\begin{equation}\label{isc:eq:upperbound}
 \supp_t Q(P_\alpha)\subseteq[b,c+d a_\alpha].
\end{equation}
The elementary factorization
\[
 F^j-P_\alpha^j=(F-P_\alpha)
  \sum_{\ell=0}^{j-1}F^{j-1-\ell}P_\alpha^\ell
\]
and the nonnegative valuations of the factors in the finite sum give
\begin{equation}\label{isc:eq:tailbound}
 v\bigl(Q(F)-Q(P_\alpha)\bigr)\ge b+a_{\alpha+1}.
\end{equation}
By \cref{isc:lem:gapscale}, eventually
\[
 b+a_{\alpha+1}>c+d a_\alpha.
\]
Since $Q(F)=0$, \eqref{isc:eq:tailbound} says that every nonzero exponent of
$Q(P_\alpha)$ is above the upper bound in \eqref{isc:eq:upperbound}. Therefore
$Q(P_\alpha)=0$ for every sufficiently large $\alpha$. This gives
infinitely many distinct roots of the nonzero finite-degree polynomial
$Q$, a contradiction. The argument uses no characteristic-zero assumption.
\end{proof}

\begin{corollary}[The factorial witness]\label[corollary]{isc:cor:factorial}
In $\R((t^\Q))$ or $\C((t^\Q))$, the series
\[
 F=1+\sum_{n=1}^{\infty}t^{n!}
\]
is transcendental over the fraction field of all bounded-above-support
series, hence over $\Frac\A{k}{\Q}$.
\end{corollary}
\begin{proof}
The positive integers $n!$ are cofinal in $\Q$ and satisfy
$(n+1)!-d n!=(n+1-d)n!\to+\infty$ for each fixed $d$. Thus the proof of
\cref{isc:thm:cofinalgap} applies to this explicit sequence. For clarity, its
partial sums start with $1+t$, $1+t+t^2$, $1+t+t^2+t^6$; there is no
duplicate exponent from a $0!$ term.
\end{proof}

\wtag{} The same exponents carry the continuum of independent actual
infinitesimals of \lbl{bst:thm:actualreal}. The omnific-preserving report's
\lbl{opa:as:ex:smallfrac} uses $\sum t^{n!}$ at $\Gamma=\Z$ for the weaker
statement that it is not a fraction of $\mathcal R_\Z$.

\begin{proof}[Proof of \cref{isc:main:fractions}]
By \cref{isc:main:copies}, the common field is $L_k=\K{k}{H}$ and the
common integer ring is $\A{k}{H}$. Each full integer copy has fraction
field equal to its field copy by \cref{isc:prop:uniform}.
If $H\ne0$, combine \cref{isc:prop:boundedfrac,isc:thm:cofinalgap} to obtain an
element of $L_k$ transcendental over $\Frac\A{k}{H}$. If $H=0$, the
common ring is $\Z$ or $\Z[i]$ and its fraction field is $\Q$ or
$\Q(i)$. The ambient constant field $\R$ or $\C$ contains a
transcendental element over that countable field: its algebraic elements
are a countable union of finite root sets, whereas the constant field is
uncountable. Finally \cref{isc:cor:factorial} supplies the specified witness
for $H=\Q$. This proves every assertion.
\end{proof}

\begin{corollary}[Failure witnessed by the denominator equation]
\label[corollary]{isc:cor:noshareddenom}
For $H=\Q$ and $F$ from \cref{isc:cor:factorial}, there are nonzero
$d_\alpha\in R_\alpha$ and $d_\beta\in R_\beta$ with
$d_\alpha F\in R_\alpha$ and $d_\beta F\in R_\beta$, but there are no
$a,d\in R_\alpha\cap R_\beta$ with $d\ne0$ and $dF=a$.
\end{corollary}
\begin{proof}
The separate denominators follow even uniformly from \cref{isc:prop:uniform}.
A shared solution would express $F=a/d$ in $\Frac\A{\R}{\Q}$,
contradicting transcendence. The same proof works for the Gaussian rings.
\end{proof}

\subsection{Where the two transcendence arguments differ}

The gap above an ordinary compositum in \cref{isc:sec:join} is detected by an
infinite algebraically independent \emph{coefficient} family in a separated
scale decomposition. The gap above the shared fraction field in this section
is detected by a cofinal sequence of \emph{support exponents}, even with
all nonzero coefficients equal to one. Neither proof is a cardinality-only
argument, and neither concludes merely that an element lacks a particular
chosen representation: each excludes every finite-degree algebraic relation
over the stated base.

\section{A set-sized version of the construction}\label{isc:sec:setversion}

The class construction has a direct conditional set-sized analogue. This
both isolates the exact saturation used and gives a more conventional target
for formalization.

\begin{definition}
For an uncountable regular cardinal $\kappa$, a linear order is an
$\eta_\kappa$-order if every cut $L<R$ with
$|L\cup R|<\kappa$ has a strict intermediate point. Empty sides are
allowed.
\end{definition}

\begin{theorem}[Independent embeddings in a saturated ordered group]
\label{isc:thm:setversion}
Let $\kappa$ be uncountable regular, and let $V$ be a divisible ordered
abelian group of cardinality $\kappa$ whose order is an $\eta_\kappa$-order.
If $H\subseteq V$ is a divisible subgroup of cardinality less than
$\kappa$ and $2\le|I|\le\kappa$, then there are ordered additive proper
embeddings $h_i:V\to V$ fixing $H$ whose ranges are independent over $H$.
\end{theorem}
\begin{proof}
Enumerate the requirements $I\times V$ in $\kappa$ stages. Before each
stage, fewer than $\kappa$ elements and partial graph entries have been
introduced. Their rational span still has size less than $\kappa$;
uncountability and regularity ensure this also at limit stages below
$\kappa$.

The one-point proof works verbatim. To avoid the current span $E$, form the
set $T$ of its elements realizing the current cut and realize the enlarged
cut $L\cup T<R$. All these sets have size less than $\kappa$, so the
$\eta_\kappa$ hypothesis provides the new value. Independence is preserved
by exactly the same finite-relation argument. At the end all domains are
$V$, and another independent range witnesses properness of each one.
\end{proof}

This theorem is conditional on the existence of the stated group; it does
not assume that every cardinal admits such a structure. The uncountability
hypothesis is purposeful: in a countable group, adjoining one rational
vector already gives a countable domain, so mere realization of finite
order cuts would not justify the same proof.

Lifting these embeddings to full set-sized Hahn fields gives the same
intersection and disjointness statements. Those Hahn fields need not have
cardinality $\kappa$; the theorem places no unwarranted cardinality bound on
all their allowable supports.

\section{Further research questions}\label{isc:sec:questions}

The questions below are proposed by the present theorems. They are not
presented as a list of previously published open problems, nor as claims
that their answers must be difficult. The first four most directly extend
the proved results.

\begin{question}[Intersections outside the full Hahn category]
Which set-sized real closed subfields $L\subseteq\No$ containing $\R$
can occur as the exact intersection of two strong, omnific-preserving
self-embedding images if the monomial-preserving hypothesis is removed?
\end{question}
\noindent\emph{Why it is the next step.} \Cref{isc:thm:intersectionclassification}
answers the monomial-preserving case exactly. A non-Hahn example, or an
obstruction forcing Hahn closure for a larger class of embeddings, would
identify the precise information carried by the chosen monomial
cross-section. Strongness alone should not be assumed to settle it.
\wtag{} \emph{Status in the collection: open and new.} It is related to
\lbl{opa:as:q:units}, which asks to classify embeddings without exact
monomial preservation, but that question does not concern intersections.

\begin{question}[Relative classification of pairs of copies]
Suppose two independent pairs have the same core $\K{\R}{H}$. Which
invariants determine whether an omnific-preserving automorphism of $\No$
fixing that core carries one pair to the other?
\end{question}
\noindent\emph{Concrete first case.} Compare pairs whose exponent groups
form different ordered amalgams over $H$, even when their abstract additive
quotients agree. The ambient ordering of mixed sums is an evident invariant
for monomial lifts, but it is not proved here to be complete.
\wtag{} \emph{Status in the collection: open.} It is the pair version, at
a fixed Hahn core, of \lbl{opa:as:q:conjugacy}, whose partial answer
\lbl{opa:as:thm:continuum} concerns single copies that move reals.

\begin{question}[An exact description of ordinary composita]
For independent groups $G_1,G_2$ over $H$, characterize the subfield
$\K{k}{G_1}\K{k}{G_2}$ inside $\K{k}{G_1+G_2}$. Determine the
transcendence degree of the full Hahn join over the ordinary compositum in
natural set-sized and class-sized examples.
\end{question}
\noindent\emph{Available starting point.} In the separated case,
\cref{isc:lem:finitecoeff} gives a necessary finite coefficient-field condition,
and \cref{isc:thm:joinobstruction} supplies a transcendental obstruction even
after algebraic closure. A necessary-and-sufficient condition would be much
stronger. For a proper-class extension, formulate dimension statements as
independent families of specified set cardinalities rather than casually
assigning a set cardinal to the whole class.
\mtag{} \emph{Status (batch 34): answered in the separated case, open in
general.} For the separated explicit copies of \cref{isc:main:join} (trivial
core, $G_0$ infinitesimal relative to $G_1$), the compositum is
characterized by a finite coefficient-rank denominator test
(\cref{isc:cp:thm:rank,isc:cp:cor:surrealrank}), and the necessary condition
of \cref{isc:lem:finitecoeff} is shown not to be sufficient
(\cref{isc:cp:cor:fginsufficient}). In the natural set-sized Laurent
examples the algebraic and Euler-differential transcendence degrees of
$F((t))$ over $F k((t))$ are both $|F|^{\aleph_0}$ when $|k|\le\cc$ and
$\trdeg(F/k)>0$ (\cref{isc:cp:thm:cardinal}). In the class-sized explicit
join, as this question asks, the answer is given by independent families of
every set cardinal of countably supported omnific integers
(\cref{isc:cp:thm:omnific}), a tail independent over any set of adjoined
parameters (\cref{isc:cp:thm:resilience}), and the absence of set-sized
generation (\cref{isc:cp:cor:nosetgen}). Open: groups over a common core $H$
that are not separated, prescribed nonzero common cores, infinite algebraic
coefficient extensions, and $|k|>\cc$ with small positive $\trdeg(F/k)$
(Research questions~\ref{isc:cp:q:amalgams}, \ref{isc:cp:q:large},
\ref{isc:cp:q:algebraic} and~\ref{isc:cp:q:cores}).
\mtag{} \emph{Status (batch 35): the second part is also answered for
explicit non-separated examples.} For interleaved real closed factors
$\K{k}{A_\kappa}$, $\K{k}{B_\kappa}$ with intersection $k$, the join has
transcendence degree exactly $2^\kappa$ over the compositum at every
infinite cardinal $\kappa$, with omnific witnesses
(\cref{isc:hj:main:cardinal}); for dense Archimedean groups it is $\cc$,
with witnesses prime in $\Oz$ and $\Oz[i]$ (\cref{isc:hj:main:prime}); a
fixed pair of explicit class support subfields has an ordinal-indexed
independent family and no set-sized transcendence basis
(\cref{isc:hj:main:class}). All rest on the mixed-support criterion
\cref{isc:hj:main:mixed}, which needs no convex separation. The first part
(an exact description of the compositum) remains open outside the
separated case (Research questions~\ref{isc:cp:q:amalgams}
and~\ref{isc:hj:q:characterization}), as do prescribed nonzero cores
(Research questions~\ref{isc:cp:q:cores} and~\ref{isc:hj:q:cores}).

\begin{question}[Descent for the arithmetic intersection]
How do integral closure, complete integral closure, and finite-type
algebraic geometry over $\A{k}{H}$ compare with those over
$\Frac\A{k}{H}$ and $\K{k}{H}$? Which descent failures besides the
single denominator equation are forced by the transcendental gap?
\end{question}
\noindent\emph{Important boundary.} The elementary fraction-field failure
proved here does not by itself calculate a normalization, a Picard group,
or a nontrivial projective module. Those require additional arguments and
should be compared with the repository's dedicated arithmetic reports.

\begin{question}[Large fixed bases]
For which proper-class divisible subgroups $H\subsetneq\No$ can independent
ordered-group copies fixing $H$ be constructed? What cut-size or cofinality
condition replaces the set-sized hypothesis of \cref{isc:thm:group}?
\end{question}
\noindent\emph{Obstacle to the present proof.} Its first partial domain and
its forbidden range span must be sets. A proper-class $H$ invalidates that
step. This is a limitation of the argument, not a proof that every such
extension is impossible.

\begin{question}[Compatibility with additional surreal operations]
Can versions of the prescribed-intersection theorem be proved for embeddings
commuting with the Gonshor exponential or with a fixed surreal derivation?
What cores would necessarily be closed under those operations?
\end{question}
\noindent\emph{First necessary checks.} If an operation is preserved by both
embeddings, their images and intersection must be closed under that operation.
The explicit maps of \cref{isc:sec:explicit} already fail to commute with the
omega-map. Therefore they cannot simply be relabeled as solutions in richer
languages. Published work on exponential surreal embeddings provides a
relevant comparison \cite{EK}.
\wtag{} \emph{Status in the collection: open.} It lies inside
\lbl{opa:as:q:exp} of the omnific-preserving report, restricted to
prescribed intersections.

\begin{question}[Simplicity and initiality]
Which prescribed intersections remain possible when each image is required
to be initial for the simplicity relation, or when a specified bound on
birthdays is imposed?
\end{question}
\noindent\emph{Reason for separating this question.} The saturation proof
uses the simplest realizer of an enlarged cut to choose one point, but it
does not preserve the global simplicity relation. Local use of a simplest
number is not an initial-embedding theorem.
\wtag{} \emph{Status in the collection: open.} Simplicity relations are
also named in \lbl{opa:as:q:exp}.

\begin{question}[Uniform maximal transcendence at a fixed core]
For a fixed nonzero set-sized $H$, obtain an independently checked sharp
calculation of
$\trdeg(\K{k}{H}/\Frac\A{k}{H})$ and relate it to $\cf(H)$ and support
cardinality. How much of this calculation can be made uniform across all
realized intersection diagrams?
\end{question}
\noindent\emph{Prior work to retain.} This is not an untouched direction:
the repository's bounded-support report already proposes a
$2^{\cf(H)}$-sized independent family and optimal support thresholds
\cite{RepoBounded}. A useful next contribution is verification, extension,
or a genuinely different uniformity statement, not rediscovery of its
one-witness special case.
\wtag{} \emph{Status in the collection:} the programme of the
bounded-support report (\lbl{bst:thm:independent}, \lbl{bst:cor:exactcard},
\lbl{bst:cor:realgroup}); its uniformity over realized diagrams is open.

\begin{question}[Elementary embeddings of the field--integer pair]
Can the embeddings be chosen elementary in
$(\No,+,\cdot,<,\Oz)$ while retaining a prescribed set-sized intersection
and the stated algebraic independence? What extra restrictions does the
integer-ring predicate impose?
\end{question}
\noindent\emph{What has actually been proved.} The field embeddings preserve
and reflect membership in $\Oz$ and commute with floor, but these facts do
not establish preservation of all formulas with alternating quantifiers.
Ordinary real-closed-field elementarity is a weaker assertion.
\wtag{} \emph{Status in the collection: open.} For the maps here it is
\lbl{opa:as:q:elementary}, with a prescribed intersection added; the
nonelementarity criterion of \lbl{opa:as:thm:elementary} needs a moved real
and does not apply.

\begin{question}[A verified set-level kernel]
Can the coset-slice lemma, finite separating-matrix argument, and saturated
ordered-group extension theorem be formalized as reusable library results,
then connected to a universe-aware model of the surreal construction?
\end{question}
\noindent\emph{Concrete target.} The next section gives declarations and
obligations at the level of mathematical specifications. A verified finite
support computation is not a substitute for these support and class-size
proofs.
\wtag{} \emph{Status in the collection: open.} Its Hahn-field kernel
overlaps \lbl{opa:as:q:formal}; none of the proposed identifiers exists in
the collection's Lean library.

\section{Proof dependencies, verification boundary, and conclusion}
\label{isc:sec:audit}

\subsection{A dependency map}

The proof of the main copy theorem follows the chain
\[
\begin{gathered}
 \text{set-cut property}
 \Longrightarrow\text{one-point avoidance}\\
 \Longrightarrow\text{independent group ranges}
 \Longrightarrow\text{strong monomial lifts}.
\end{gathered}
\]
The upgrade from support intersections to algebraic independence is a
separate chain:
\[
\begin{gathered}
 \text{coset slices}
 \Longrightarrow\text{finite separating matrix}\\
 \Longrightarrow\text{linear disjointness}
 \Longrightarrow\text{polynomial-relation descent}.
\end{gathered}
\]
The arithmetic application adds the common-denominator lemma and the
cofinal gap argument. The compositum application instead adds the explicit
separated-scale construction, iterated Hahn coordinates, and finite
coefficient-field control. Neither application is used to justify the other,
and neither is needed to justify the group construction.

\subsection{A practical formalization decomposition}\label{isc:sub:formal}

A set-level implementation could first express full Hahn series in the
usual increasing-support convention and prove the following specifications.
The names below are proposed identifiers, not claims that declarations with
these names already exist.

\begin{enumerate}[label=\textup{(\arabic*)}]
\item \texttt{cosetSlice\_linear}: restriction of a Hahn series to one
coset of a subgroup, followed by monomial normalization, is linear over the
Hahn field on that subgroup. No convexity assumption is permitted.
\item \texttt{cosetSlice\_separates} and
\texttt{finite\_slice\_matrix}: the family separates points and yields an
invertible finite matrix on a linearly independent tuple.
\item \texttt{hahnSubfields\_linearlyDisjoint}: use the same slice matrix
after enlarging the scalar subgroup from the intersection to the second
subgroup. This is the main reusable field-theoretic lemma.
\item \texttt{orderedEmbedding\_extend\_avoid}: realize a cut of the current
rational span while avoiding a prescribed small subspace, with explicit
cardinality hypotheses for \cref{isc:thm:setversion}.
\item \texttt{boundedSupport\_gap\_transcendental}: formalize the two support
bounds \eqref{isc:eq:upperbound} and \eqref{isc:eq:tailbound}, followed by the
finite root bound. Start with the factorial sequence over $\Q$ before the
transfinite cofinality version.
\item \texttt{iteratedHahn\_coefficientField}: formalize separated ordered
sums, the regrouping isomorphism, and finite coefficient-field control of a
compositum. This is the kernel of the second transcendence proof.
\end{enumerate}

The final proper-class theorem additionally needs a class or universe
framework, a set-like enumeration, and proof that every recursion stage and
every support pullback is a set. A set-level Hahn formalization by itself
does not discharge that final obligation. The repository contains relevant
Hahn and normal-form infrastructure \cite{RepoCatalogue}, but this article
adds no Lean code and makes no machine-verification claim.

\subsection{Claims deliberately not made}\label{isc:sub:notmade}

The following boundaries prevent several plausible but invalid upgrades.
The common field is a \emph{full Hahn field} on a \emph{set-sized} subgroup,
not an arbitrary prescribed real closed field. Group sums are finite sums,
not hidden summable closures. Boolean joins are Hahn joins, not ordinary
composita. Strongness concerns Hahn-summable set families, not convergence
of all real series. The embeddings preserve omnific membership, not an
unproved complete first-order theory of the pair. They are not asserted to
preserve the omega-map, exponentiation, derivation, birthdays, or simplicity.
The proofs establish the stated theorems but do not establish historical
priority or settle a named general factorization conjecture for $\Oz$.
\wtag{} \Cref{isc:sec:nonclaims} collects these boundaries with every other
limitation stated in the source, and adds those that come from the
collection.

\subsection{Conclusion}

Full surreal copies can be arranged with far more control than a bare
universality statement provides. An arbitrary set-sized full Hahn core can
be fixed exactly, an ordinal-indexed family of copies can be made independent
over it, and their arithmetic parts remain genuine copies of the entire
omnific ring. The same construction simultaneously controls the surcomplex
field and Gaussian omnific integers.

The two descent failures are complementary. Passing from intersecting
omnific rings to their fraction field can lose transcendental elements that
belong to both full field copies, even though each copy provides one uniform
denominator for the whole common set field. Passing from two full field
copies to their ordinary compositum can likewise miss transcendental
normal forms allowed by their combined exponent group. Thus support closure,
ordinary field generation, and arithmetic localization are three distinct
operations. The prescribed-intersection construction makes their difference
visible within the actual surreal and surcomplex fields, rather than only
in abstract auxiliary examples.

\section{Consolidated ledger of non-claims}\label{isc:sec:nonclaims}

\wtag{} The source states its limitations where they arise, and each is
kept in place. This ledger collects them with their locations. ``README''
refers to the delivered README, which is not shipped; its content is in the
collection README. Items S1--S24 are the source's. Items W1--W8 were added
when the report joined the collection. \mtag{} The limitations of the second
manuscript (Sections~15--27) are listed separately, as items C1--C23 and
M1--M5 of \cref{isc:cp:sub:nonclaims}; those of the third manuscript
(Sections~28--41) as items J1--J25 and JM1--JM6 of
\cref{isc:hj:sub:nonclaims}.
\begin{enumerate}[label=S\arabic*.,leftmargin=2.6em,itemsep=2pt]
\item The common field is a \emph{full} Hahn field on a \emph{set-sized}
divisible subgroup. It is not an arbitrary prescribed real closed or
set-sized subfield, not the ordinary field generated by a parameter set,
not its real closure, and not a non-Hahn subfield
(after \cref{isc:cor:parameters}; \cref{isc:sub:notmade}; README).
\item Group sums $G_J$ are finite sums, not hidden summable closures
(after \eqref{isc:eq:diagram}; \cref{isc:sub:notmade}).
\item The joins of \cref{isc:thm:boolean} are Hahn joins, not ordinary
composita; \cref{isc:main:join} shows that the difference can be
transcendental (\cref{isc:rem:join}; \cref{isc:sub:notmade}).
\item Strongness concerns Hahn-summable set-indexed families, not
convergence of real series or continuity; ordinary analytic convergence is
not used in any proof (\cref{isc:sec:foundations}; \cref{isc:sub:notmade}).
\item Elementarity is claimed only formula by formula in the languages of
ordered fields and of fields, by quantifier elimination for real closed and
algebraically closed fields. Nothing is asserted in the language with a
predicate for $\Oz$, the omega-map or a derivation, and algebraic
independence is not identified with a model-theoretic independence notion
of a richer theory. Preservation and reflection of $\Oz$ and commutation
with the floor do not give elementarity of the pair
(\cref{isc:sub:model}; research question~9; \cref{isc:sub:notmade}).
\item The embeddings are not asserted to preserve the omega-map, the Gonshor
exponential, the Berarducci--Mantova derivation, birthdays or simplicity.
Monomial preservation is not omega-map commutation: $h\ne\Phi_h$ in
general, and $\Phi_0(\omega^1)\ne\omega^{\Phi_0(1)}$ for the explicit maps,
so they cannot be relabelled as solutions in richer languages
(\cref{isc:warn:omega}; \cref{isc:sec:explicit}; research question~6;
\cref{isc:sub:notmade}).
\item No proper-class base $H$ is treated. The obstacle is a limitation of
the argument, whose first partial domain and forbidden span must be sets,
not a proof that such extensions are impossible (research question~5).
\item \Cref{isc:thm:setversion} is conditional on the existence of the
stated saturated group; it does not assume that every cardinal admits one.
Uncountability is needed for the proof. The lifted Hahn fields need not have
cardinality $\kappa$ (\cref{isc:sec:setversion}).
\item The fraction-field failure does not by itself compute a
normalization, a Picard group or a nontrivial projective module
(research question~4).
\item The bounded-support transcendence mechanism of \cref{isc:sec:gaps}
is not new. It belongs to bounded-support Hahn arithmetic; the
bounded-support report gives stronger independent families, and the
non-fraction obstruction is Proposition~2.4.5 of \cite{LM}. It is reproved
so that \cref{isc:main:fractions} does not depend on an unreviewed
repository result, not claimed as a discovery
(\cref{isc:subsec:provenance}; \cref{isc:sec:gaps}; research question~8;
README).
\item The ingredients are not claimed as new: the set-cut property, normal
forms, Hahn multiplication and inversion, and the closedness criteria. The
group extension is a concrete saturation construction, not a new saturation
principle; the monomial lift is standard Hahn functoriality; the coset-slice
lemma is an elementary reusable lemma, not a certified new theorem of Hahn
field theory (\cref{isc:subsec:provenance}; title page).
\item Priority is not certified. Targeted literature searches do not certify
that the combined formulation is absent from published or unpublished work,
and no claim of having audited every report or every later repository
revision is made (\cref{isc:subsec:provenance}; \cref{isc:sub:notmade};
README).
\item The article is unrefereed and not formalized. It adds no Lean code and
makes no machine-verification claim. The identifiers of
\cref{isc:sub:formal} are proposals, not existing declarations; a set-level
Hahn formalization does not discharge the class-size obligations, and a
verified finite support computation is no substitute for them
(title page; \cref{isc:sub:formal}; research question~10; README).
\item No named longstanding or published conjecture is claimed solved, and
no general factorization conjecture for $\Oz$ is settled (title page;
\cref{isc:sub:notmade}; README).
\item \Cref{isc:main:fractions} is the failure of one descent operation, not
of $\Frac(\Oz)=\No$, which the argument uses (after \cref{isc:main:fractions}).
\item \Cref{isc:lem:finitecoeff} gives a necessary coefficient-field
condition, not membership in $M$ of every series with coefficients in $D$.
The compositum argument is not the invalid inference that every infinite
sum lies outside a compositum; a geometric series already lies in one
(after \cref{isc:lem:finitecoeff}; \cref{isc:rem:infinite}; README).
\item The class framework uses no tensor product indexed by a proper class
and no proper-class-dimensional tensor object, no class truth or
satisfaction predicate, and no recursion with a proper-class state at a
stage. The cut property is never applied to proper-class sides
(\cref{isc:sec:foundations}; \cref{isc:sub:model}).
\item \Cref{isc:thm:intersectionclassification} is a classification inside
the strong, real-fixing, monomial-preserving class only. Strongness alone
should not be assumed to settle research question~1.
\item The maps $h_\alpha$ fix $H$ as a set of exponent values and the
$\Phi_\alpha$ fix $\K{\R}{H}$; the two roles are not identified
(\cref{isc:rem:base}).
\item The positive-exponent truncation plus the ordinary floor of the
constant coefficient is not always the omnific floor
(\cref{isc:rem:floor}). Scaling all exponents by a constant gives an
automorphism, not a proper copy (\cref{isc:rem:scaling}).
\item The saturation proof uses a simplest realizer locally but does not
preserve the simplicity relation; this is not an initial-embedding theorem
(research question~7).
\item The research questions are proposed by these theorems; they are not
presented as previously published open problems or as claims of difficulty.
Research question~8 is not an untouched direction (\cref{isc:sec:questions}).
\item For coefficient fields other than $\R$ and $\C$ the Hahn-field
notation is formal, not a subfield of $\No$ or $\SC$; fraction fields of
class domains are ratios inside the ambient field, not quotients of a
proper class (\cref{isc:sec:foundations}; \cref{isc:sec:fractions}).
\item The encyclopedia page is orientation only; Ehrlich--Kaplan is a
comparison, not a proof input; the repository reports were used for scope
and nonduplication, not as dependencies. The delivered production checks
(compilation, rendering and visual inspection of the PDF) are not formal
verification of the mathematics (bibliography; README).
\end{enumerate}
\begin{enumerate}[label=W\arabic*.,leftmargin=2.6em,itemsep=2pt]
\item No \texttt{isc:} label has a Lean mapping in \rpath{FORMALIZATION.md}.
Independent proof review and formalization are pending. The package shipped
no verification code, data or audit, so the report has no suite to rerun
(\cref{isc:sec:collection}; Appendix~\ref{isc:app:files}). \mtag{} This
holds for source~01; the second manuscript shipped a finite-check suite,
rerun in Appendix~\ref{isc:cp:app:files}, which tests none of source~01's
theorems; so did the third (Appendix~\ref{isc:hj:app:files}), which tests
none of them either.
\item \Cref{isc:thm:lift,isc:prop:arithmeticlift,isc:prop:classification}
and the inner lift of \cref{isc:sec:explicit} are the real-fixing case of
the omnific-preserving report's classification and its inner lift; no
priority is claimed for them relative to that report. \Cref{isc:cor:parameters}
extends \lbl{opa:as:thm:parameters} but does not claim its topological
assertions (\cref{isc:sec:collection}).
\item \Cref{isc:lem:denominator} is \lbl{odg:thm:fractions},
\cref{isc:prop:boundedfrac} is the localization of
\lbl{bst:prop:localdensity}, the coset slice is \lbl{bst:eq:projection},
and \cref{isc:thm:cofinalgap,isc:cor:factorial} duplicate
\lbl{bst:cor:onegap} and the one-series case of \lbl{bst:thm:independent}
(\cref{isc:sec:collection}).
\item Real fixing is a choice. For the explicit copies the bottom gap is
nonzero, so exact-monomial embeddings with the same exponent maps that
preserve $\Oz$ and move reals exist; nothing is claimed about them or about
intersections in that larger class (\cref{isc:sec:collection}).
\item The large-cardinal report's question ``Joint algebraic dependence of
the images'' is not answered: its image $J(\No)$ is not a full Hahn field
on a subgroup (\cref{isc:sec:collection}).
\item None of the omnific-preserving report's questions is answered; in
particular \lbl{opa:as:q:composition}, \lbl{opa:as:q:conjugacy},
\lbl{opa:as:q:elementary} and \lbl{opa:as:q:exp} stay open. None of the
bounded-support report's questions is answered (\cref{isc:sec:collection}).
\item The source's statements about the collection are those of its pin and
are kept as provenance (\cref{isc:subsec:provenance};
Appendix~\ref{isc:app:pinned}).
\item The proofs were read when the report was written
(Appendix~\ref{isc:app:checked}). That reading is not an independent proof
review.
\end{enumerate}

\section{Beyond composita of surreal copies: a later manuscript}
\label{isc:cp:sec:intro}

\mtag{} \Cref{isc:cp:sec:intro,isc:cp:sec:foundations,isc:cp:sec:rank,isc:cp:sec:disjoint,isc:cp:sec:lacunary,isc:cp:sec:generic,isc:cp:sec:cardinal,isc:cp:sec:surreal,isc:cp:sec:resilience,isc:cp:sec:arithmetic,isc:cp:sec:audit,isc:cp:sec:questions,isc:cp:sec:verify}
print a second manuscript, written against this report. In the report's
local numbering (Appendix~\ref{isc:cp:app:source}) the original manuscript
is source~01 and the addition is \emph{source~02}: manuscript~05 of batch~34
of the collection, archive \texttt{Surreal\_Composita\_Research},
\emph{Beyond Composita of Surreal Copies: Exact rank tests, differential
transcendence, and omnific witnesses}, a 29-page PDF dated 23 September
2026 and pinned to repository commit
\nolinkurl{7b256e0792df7718995b816ad12a77bf56731f8f}. At that pin this
report's \texttt{article.tex} was the blob
\nolinkurl{7d5cfe644e7c902f2b594179c337a6c633350fc3}. Since then only
three \wtag{} sentences about the bounded-support report have changed
(\cref{isc:subsec:provenance,isc:sec:collection,isc:sub:bounded}), so the
source's description of \cref{isc:sec:explicit,isc:sec:join,isc:sec:questions}
is accurate. The source takes Research question~3 of
\cref{isc:sec:questions} as its target.

Section~$n$ of the source is Section~$n+14$ here for $n=1,\ldots,12$, and
every numbered statement keeps its position: the source's Theorem~$n.m$ is
Theorem~$(n+14).m$ here, and likewise for lemmas, corollaries, remarks,
definitions and examples. Its Questions~1--12 are Research
questions~11--22. Its Appendices~A and~B form \cref{isc:cp:sec:verify}; its
notation index (Appendix~C) is folded into the table of
\cref{isc:cp:sec:conventions}. Where the source cites this report by section
number, the citation is replaced by a cross-reference to the same section;
the numbering of Sections~\ref{isc:sec:introduction}--\ref{isc:sec:questions} has not
changed since the pin. Where the source reproduces material already printed
here (the explicit copies of \cref{isc:sec:explicit} and the re-expansion of
\cref{isc:sec:join}), it is printed once and the source is credited
(\cref{isc:cp:sub:copies}). Text added in the merge is marked \mtag; symbols
renamed from the source are listed in \cref{isc:cp:sec:conventions}. Every
result, proof, example, question and limitation of the source is printed.

\subsection{Field generation is not Hahn summation}

A Hahn field permits well-ordered, potentially infinite supports. An
ordinary field compositum permits only finitely many field operations on
finitely many elements at a time. These are different closure operations.
Even when two Hahn subfields have a transparent intersection, their
ordinary compositum need not contain all the series supported on the sum
of their exponent groups.

This report constructs explicit strong surreal self-embeddings whose images
$K_0,K_1$ intersect in the real constants (\cref{isc:thm:explicit}), and
\cref{isc:sec:join} proves that the full Hahn join $E$ contains an element
transcendental over $M=K_0K_1$ (\cref{isc:main:join}). Research question~3
asks for an exact description of ordinary composita and the transcendence
degree of the join over the compositum. Source~02 takes that question as
its target. The explicit copies and the one-series coefficient-transcendence
obstruction of \cref{isc:thm:joinobstruction} are this report's; neither is
claimed by source~02 as a new construction.

There are two genuinely different ways for a series to escape the
compositum. Its coefficients can use infinitely much algebraically
independent information. More subtly, its coefficients can all lie in one
field $k(x)$ and nevertheless have complexity in $x$ that defeats every
finite algebraic or differential relation over the compositum. The second
phenomenon shows that the necessary finite coefficient-field condition of
\cref{isc:lem:finitecoeff} is far from sufficient.

\subsection{Main conclusions}

The first conclusion is an exact test. If the \emph{coefficient rank} of a
Hahn series means the dimension over $k$ of the span of its coefficients,
then
\begin{equation}\label{isc:cp:eq:introtest}
 x\in F k((t^\Gamma))
 \quad\Longleftrightarrow\quad
 \exists q\ne0\;\bigl(\crk_k(q)<\infty\ \text{and}\ \crk_k(qx)<\infty\bigr).
\end{equation}
This is an intrinsic test once the coefficient field and Hahn coordinates
have been fixed. It does not assert decidability for arbitrary infinite
input. The proof is elementary tensor algebra and is presented as such.

The second conclusion is a maximal-cardinality theorem. Write
$\cc=2^{\aleph_0}$. For $|k|\le\cc$ and any coefficient extension $F/k$ of
positive transcendence degree,
\begin{equation}\label{isc:cp:eq:introdimension}
 \dtrdeg\bigl(F((t))/F k((t))\bigr)
 =\trdeg\bigl(F((t))/F k((t))\bigr)=|F|^{\aleph_0}.
\end{equation}
Thus even $F=k(x)$, a single transcendental coefficient extension, produces
the full continuum of differential freedom when $k$ is $\Q$, $\R$, or $\C$.
Here the derivation acts only on the outer series variable.

The third conclusion concerns actual omnific integers. For the explicit
separated copies of \cref{isc:sec:explicit}, with the order isomorphism
$\psi(u)=\frac12+\frac{u}{2(1+|u|)}$ of \cref{isc:lem:compression} and
ordinal arithmetic used only in the label $\omega\cdot\alpha+n$, define
\begin{equation}\label{isc:cp:eq:introwitness}
 Y_\alpha
 =\sum_{n\ge1}
 \omega^{\,\omega^{15/4}-n\omega^{7/2}
                    +\omega^{\psi(\omega\cdot\alpha+n)}}
 \qquad(\alpha\in\On).
\end{equation}
Every exponent in this normal form is positive, and the exponents strictly
decrease with $n$. Hence $Y_\alpha$ is a purely infinite omnific integer with
countable support. All its Euler jets, together with the jets of all the
other $Y_\alpha$, are algebraically independent over the compositum in the
precise finite-relation sense. Moreover, given \emph{any set} $P\subseteq E$,
there is an ordinal $\alpha_0$ such that the tail
$\{Y_\alpha:\alpha\ge\alpha_0\}$ remains differentially algebraically
independent over $M\langle P\rangle$.

This last assertion is stronger than producing a fresh witness separately
for each parameter set: a tail of one predefined family works. It also
avoids assigning a set cardinal to a proper-class transcendence degree.
\mtag{} For $\alpha=0$ the label is $\beta_{0,n}=n$, and
$Y_0=\omega^{\omega^{15/4}}S$, where $S$ is the witness
\eqref{isc:eq:explicitS} of \cref{isc:main:join}. Source~02 thus extends this
report's witness to an ordinal family inside $\Oz$, as it says.

\subsection{Scope, sources, and the meaning of new}\label{isc:cp:sub:scope}

Hahn arithmetic, linear disjointness, algebraic extension of derivations in
characteristic zero, and the surreal normal-form construction are
established tools. Standard Hahn derivation frameworks appear in
Kuhlmann--Matusinski \cite{KM}; recurrence descriptions of generalized
series subfields are studied by Krapp--Kuhlmann--Serra \cite{KKS}. The
finite-rank characterization below is a coefficient-space counterpart, not a
replacement for their recurrence theory. Omnific arithmetic has a
substantial factorization literature, including L'Innocente--Mantova
\cite{LM}; no factorization conjecture is addressed here.

The candidate contributions of source~02 are the rapid coefficient-degree
differential-independence theorem, its combination with independent
coefficient trees to give \eqref{isc:cp:eq:introdimension}, and the explicit
omnific tail-resilience theorem and arithmetic consequences. The literature
comparison is targeted, not exhaustive. These claims are proved in the
source but are not certified as previously unpublished. The differential
statements use a specified diagonal Euler derivation, \emph{not} the
exponential-compatible surreal derivation of Berarducci--Mantova \cite{BM}.

Source~02's own source audit (shipped as
\rpath{02-composita-SOURCE\_AUDIT.md}) lists six candidate contributions:
the rapid coefficient-degree theorem over $k((t))(x)$ with arbitrary Laurent
coefficients in the relation; private coefficient-block independence over
the whole finitely-generated-coefficient envelope and its tree
amplification; their combination into the exact degrees
$|F|^{\aleph_0}$; the explicit purely infinite omnific family with support
of order type $\omega$ and the sharpness of its support cardinality; tail
resilience of that predefined family under arbitrary set-sized parameter
adjunction, including differential parameters; and persistent
nonintegrality of the arithmetic join after set adjunction, with the
Gaussian versions. It calls them proposed research contributions with
written proofs; a targeted source search did not find them elsewhere, and
that absence does not establish novelty.

\subsection{What the addition settles in this report}
\label{isc:cp:sec:answers}

\mtag{}
\begin{enumerate}[leftmargin=*,itemsep=3pt]
\item \emph{Research question~3} (an exact description of ordinary
composita) is answered in the \emph{separated} case with trivial common
core, which is the case of \cref{isc:main:join}, and in the stated cardinal
cases. Its first part, the characterization of the compositum, is the
finite coefficient-rank denominator test (\cref{isc:cp:thm:rank}) and, for
the explicit copies, \cref{isc:cp:cor:surrealrank}. Its second part, the
transcendence degree of the join over the compositum, is answered in the
natural set-sized Laurent examples by \cref{isc:cp:thm:cardinal}: both the
algebraic and the Euler-differential transcendence degree of $F((t))$ over
$F k((t))$ are $|F|^{\aleph_0}$ when $|k|\le\cc$ and $\trdeg(F/k)>0$. In
the class-sized explicit join it is answered in the form the question asks
for, by independent families of every set cardinal
(\cref{isc:cp:thm:omnific}), a tail that survives every set of parameters
(\cref{isc:cp:thm:resilience}), and the absence of set-sized generation
(\cref{isc:cp:cor:nosetgen}); no set cardinal is assigned to the class
extension. Open: independent groups that are not separated (neither
infinitesimal relative to the other), copies with a prescribed nonzero
common core $H$, infinite algebraic coefficient extensions, and $|k|>\cc$
with small positive $\trdeg(F/k)$. These are Research
questions~\ref{isc:cp:q:amalgams}, \ref{isc:cp:q:cores},
\ref{isc:cp:q:algebraic} and~\ref{isc:cp:q:large}. \mtag{} Since batch~35
the transcendence degree is also known for explicit non-separated examples
(\cref{isc:hj:sec:answers}); the description of the compositum there is
still open.
\item \Cref{isc:lem:finitecoeff} is shown to be far from sufficient: one
transcendental coefficient $x$ already gives a series with all coefficients
in $k(x)$ that is differentially transcendental over the compositum
(\cref{isc:cp:cor:fginsufficient}).
\item \Cref{isc:thm:joinobstruction} is strengthened for the explicit
witness: $S=\omega^{-b}Y_0$, with $\omega^{-b}\in K_1$ a constant of the
Euler derivation $\Delta$ of \cref{isc:cp:sub:euler}, so $S$ is
differentially transcendental over $\Cfg(K_0/\R;G_1)\supseteq M$ by
\cref{isc:cp:thm:omnific} (see the \mtag{} note after the proof of
\cref{isc:main:join}).
\item No other question of this report is answered. Research question~4
(descent for the arithmetic intersection $\A{k}{H}$) is untouched: the
nonintegrality result, \cref{isc:cp:thm:arithmetic}, concerns the join of two
copies, not their intersection, and computes no normalization.
\end{enumerate}

\subsection{Conventions and renamed symbols}
\label{isc:cp:sec:conventions}

\mtag{} Source~02 uses both exponent conventions of this report: the
decreasing large-$\omega$ convention $\sum c_g\omega^g$ with
reverse-well-ordered support, written $\K{k}{G}$ as here, and the
increasing small-$t$ convention $F((t^\Gamma))$ with well-ordered support,
identified by $t^\gamma=\omega^{-\gamma}$ as in \cref{isc:sub:bounded}. In
\cref{isc:cp:sec:foundations,isc:cp:sec:rank,isc:cp:sec:disjoint,isc:cp:sec:lacunary,isc:cp:sec:generic,isc:cp:sec:cardinal}
the fields are abstract $t$-series fields over arbitrary coefficient fields,
not subfields of $\No$. From \cref{isc:cp:sec:surreal} on, the fields are
the explicit surreal fields of \cref{isc:sec:explicit,isc:sec:join}, with
$k=\R$ (and $\C$ in \cref{isc:cp:thm:complex}).

The following symbols of the source are renamed; everything else is as in
the source.
\begin{center}\small
\begin{longtable}{@{}p{0.27\textwidth}p{0.20\textwidth}p{0.45\textwidth}@{}}
\toprule
Here & Source 02 & Reason\\
\midrule
\endhead
$\Gamma$, $\Gamma_0$, $\Gamma'$: the outer ordered group, a subgroup, a
set subgroup & $H$, $H_0$, $H'$ & $H$ is the prescribed common core of
\cref{isc:main:copies}, and the source itself uses $H$ in that sense in its
Question~4 (Research question~\ref{isc:cp:q:cores})\\
$g>0$ in $\Gamma$ with $\lambda(g)=1$; $g=\omega^{7/2}$ in
\cref{isc:cp:sec:surreal,isc:cp:sec:resilience,isc:cp:sec:arithmetic} & $h$
& $h$, $h_\alpha$ are ordered additive embeddings here; this report already
writes $g=h_1(1)=\omega^{7/2}$ (proof of \cref{isc:main:join}) and uses $g$
in this role in \cref{isc:thm:joinobstruction}\\
$\omega^b$, the shift multiplier & $s$ & $s_n=\omega^{g_n}$ are the
coefficients of \eqref{isc:eq:explicitS}\\
$F[L]$, the ring generated by $F$ and $L$ & $T$ & $T$ is also the source's
inner support set\\
$\Theta$, $G_\Theta$, $C_\Theta$: inner support set, its support group and
localized coefficient field & $T$, $H_T$, $C_T$ & as above; $G_\Theta$ is a
subgroup of $G_0$\\
$R_0[R_1]$, the ring generated by $R_0\cup R_1$; $R_0[R_1][S]$;
$R_0[R_1][i]$ & $B$; $B[S]$; $B_{\C}$ & $B$ is a subgroup in
\cref{isc:lem:iterated}, and $\B{k}{H}$ is the bounded-support ring\\
$j=0,1$, index of the two explicit copies & $i$ & $i=\sqrt{-1}$ in
\cref{isc:cp:sec:arithmetic}; \cref{isc:main:join} writes $K_j$\\
$\mathfrak A$, a family of infinite subsets of $\N_{>0}$ & $\mathcal A$ &
$\A{k}{H}$ is the integer ring of \cref{isc:sec:fractions}\\
$p_{ij}=P_{Y_{ij}}(\mathbf f)$, evaluated partials & $A_{ij}$ & $A_i$ are
members of $\mathfrak A$ in the same proof\\
$W_i(X,x)$, the leading polynomials in \eqref{isc:cp:eq:Hpolynomial} & $H_i(T,x)$ &
$H$ and $T$ as above\\
$w_i$, $\mathbf w$: tails in the proof of \cref{isc:cp:thm:lacunary} & $R_i$,
$\mathbf R$ & $R_j$ are the omnific copies; $R$ is a cut side\\
$\nu$, the degree in $x$ in the same proof & $D$ & $D$ is a finitely
generated coefficient field, as in \cref{isc:lem:finitecoeff}\\
$\gamma_0+\gamma_1$ with $\gamma_j\in G_j$, and $e$: generic
exponents from \cref{isc:cp:sec:surreal} on & $g_0+g_1$, $g$ & $g$ is
$\omega^{7/2}$ above\\
$\crk_k$, coefficient rank & $\operatorname{crk}_k$, macro \texttt{\textbackslash rk} &
same symbol; only the macro name differs\\
\bottomrule
\end{longtable}
\end{center}
No normalization is changed by these renamings. In particular
$\omega^{7/2}=g$ and $\omega^{15/4}=b$ are the same surreals as in the
source, and \eqref{isc:cp:eq:introwitness} is the source's formula verbatim.

\paragraph{Agreements.} $\K{k}{G}$, $E=\K{\R}{G_0+G_1}$, $M=K_0K_1$,
$K_j=\K{\R}{G_j}$, $G_j=h_j(\No)$, $I_j=(3j,3j+1)$, $\psi$, $h_j$, $\Phi_j$,
the convention $t^\gamma=\omega^{-\gamma}$, the finitely generated field $D$
of \cref{isc:lem:finitecoeff}, and a series $S_i=\sum a_{b_i(n)}t^{ng}$ of the
shape of \cref{isc:thm:joinobstruction} mean the same here as in
\cref{isc:sec:explicit,isc:sec:join}. The source's $g_{\alpha,n}$ and
$a_{\alpha,n}$ extend this report's $g_n=\omega^{\psi(n)}$ and
$s_n=\omega^{g_n}$: $g_{0,n}=g_n$ and $a_{0,n}=s_n$.

\paragraph{Local letters and tempting false readings.}
\begin{itemize}[leftmargin=*,itemsep=3pt]
\item In
\cref{isc:cp:sec:foundations,isc:cp:sec:rank,isc:cp:sec:disjoint,isc:cp:sec:lacunary,isc:cp:sec:generic,isc:cp:sec:cardinal},
$E=F((t^\Gamma))$, $L=k((t^\Gamma))$ and $M=FL$. This $L$ is \emph{not} the
common field $L_k=\K{k}{H}$ of \cref{isc:main:copies} and not a cut side;
this $F$ is a coefficient field, not the witness $F$ of
\cref{isc:main:fractions}. In the surreal instance $K_0$ plays the role of
$F$, $K_1\simeq\R((t^{G_1}))$ that of $L$, and $\R$ that of $k$.
\item $\lambda:\Gamma\to k$ is an additive \emph{Euler weight}. It is not a
coset slice $\lambda_{g,H}$ of \eqref{isc:eq:slice}, and not a class
functional.
\item $\partial$ and $\Delta$ are Euler derivations, $d$ and $d_b$
coefficient derivations. The omnific-preserving report reserves $d$ and $D$
for derivations; here $D$ is always a field, and $d$ is also a total degree,
$d_n$ a degree sequence, and $d\in G_j$ a denominator exponent in
\cref{isc:cp:lem:clearing}; none of these is an Euler derivation.
$D_\R=\Z$ and $D_\C=\Z[i]$ keep their meaning from \cref{isc:sec:fractions}.
\item $\jet(y_i:i\in I)$ is a set of jets, not an index set $J$. $I$ is an
index set and $I_j$ an interval, as in \cref{isc:sec:conventions}. The
index set $B$ of a coefficient family $(a_b)_{b\in B}$ and the label maps
$b_i$ are local to \cref{isc:cp:sec:generic}; $b=\omega^{15/4}$ is a
positive element of $G_1$, as the letter $b$ is in \eqref{isc:eq:separated}.
\item $\kappa$ is an infinite cardinal in \cref{isc:cp:cor:tree} and a set
cardinal in \cref{isc:cp:thm:omnific}; it is not $\cf(H)$ of
\cref{isc:sec:gaps}. $\mu=|F|$ in \cref{isc:cp:thm:cardinal}. $\cc$ is
$2^{\aleph_0}$, with no assumption on its value as an aleph.
\item $P$ is a polynomial relation in
\cref{isc:cp:sec:lacunary,isc:cp:sec:generic} and a set of parameters in
\cref{isc:cp:sec:resilience,isc:cp:sec:arithmetic}; $M\langle P\rangle$ is
the differential field generated by $M$ and all jets of $P$. $O=E\cap\Oz$.
$Y_\alpha$ and $Z_\alpha$ are the omnific witnesses (the letter $Y$ is also a
jet variable $Y_{ij}$). $\tau=\omega^{-g}$ is a monomial, not an exponent
map.
\item \emph{Tempting false reading:} ``$\crk_k(x)<\infty$ is necessary for
$x\in M$.'' It is necessary only for $x\in F[L]$; the compositum $M$ is the
fraction field of $F[L]$ and contains series of infinite coefficient rank
(\cref{isc:cp:ex:geometric}).
\item \emph{Tempting false reading:} ``$\Delta$ is a surreal derivation.''
It is a diagonal Euler derivation of the explicit join, defined through the
iterated coordinates, killing $K_0$; it is not the Berarducci--Mantova
derivation, not Hardy type, and not compatible with the exponential.
\item \emph{Strong} means preserving Hahn summability, as in
\cref{isc:sec:conventions}, not convergence in an order topology.
\end{itemize}

\paragraph{Notation index of the source (its Appendix~C), with the names
used here.} $F((t^\Gamma))$: small-$t$ Hahn field, well-ordered support in
$\Gamma$. $\K{k}{G}$: large-$\omega$ Hahn field, reverse-well-ordered
support in $G$. $E,L,M$ in the abstract sections: $F((t^\Gamma))$,
$k((t^\Gamma))$ and the ordinary field compositum $FL$. $F[L]$: the finite
coefficient-rank ring, the image of $F\otimes_kL$. $\crk_k(z)$: dimension
over $k$ of the span of all coefficients of $z$. $\Cfg(F/k;\Gamma)$: union
of $D((t^\Gamma))$ over finitely generated coefficient fields $D/k$ inside
$F$. $\partial$: abstract outer Euler derivation, with weight map
$\lambda:\Gamma\to k$. $h_j,\Phi_j$: inner additive embedding and outer
surreal field embedding. $G_j,K_j,R_j$: exponent-group image, full field
image, and omnific image of copy $j$. $E,M$ from \cref{isc:cp:sec:surreal}
on: the explicit Hahn join $\K{\R}{G_0+G_1}$ and its ordinary compositum
$K_0K_1$. $g,b,\omega^b$: $\omega^{7/2}$, $\omega^{15/4}$, and the monomial
that shifts all witness exponents into the positive cone. $\Delta$: Euler
derivation extracting the inner $\omega^{7/2}$ coefficient from the
dominant exponent coordinate, with the sign in \eqref{isc:cp:eq:Delta}.
$\beta_{\alpha,n}$: the ordinal label $\omega\cdot\alpha+n$, the only place
where ordinal arithmetic is used. $g_{\alpha,n},a_{\alpha,n}$: independent
inner exponents and the corresponding outer monomial coefficients.
$Y_\alpha,Z_\alpha$: real omnific witness and its Gaussian multiple
$(1+i)Y_\alpha$. $\Theta,G_\Theta,C_\Theta$: set of inner support labels,
its full additive support group, and the localized coefficient field.
$R_0[R_1],O$: the ring generated by the two omnific copies, and the full
omnific ring inside their Hahn join. $M\langle P\rangle$: differential field
generated by $M$ and all jets of the set $P$.

\subsection{Place in the collection}\label{isc:cp:sec:collection}

\mtag{} Source~02 cites only this report from the collection. Checked
against the collection at the placement commit \texttt{a7a435f}, the
following reports bear on it; none is a dependency of its proofs.
\begin{enumerate}[leftmargin=*,itemsep=3pt]
\item \emph{This report.} \Cref{isc:cp:lem:coset} is a variant of the
coset-slice result, \cref{isc:thm:sliceddisjoint}: there two full Hahn
fields with the same coefficient field are disjoint over the field of the
intersection group; here $F((t^{\Gamma_0}))$ and $k((t^\Gamma))$, with
different coefficient fields and $\Gamma_0\subseteq\Gamma$, are disjoint
over $k((t^{\Gamma_0}))$, by the same slice argument. The source credits
the coset slice as part of the elementary Hahn toolkit. The first clause
of \cref{isc:cp:lem:clearing}, $\Frac(R_j)=K_j$, is the first clause of
\cref{isc:prop:uniform} for the explicit copies, by the monomial clearing of
\cref{isc:lem:denominator}, as the source says. \Cref{isc:cp:lem:cfg} is
\cref{isc:lem:finitecoeff} in the abstract setting. The localization of a
parameter set in \cref{isc:cp:lem:localizeP} is the support-closure step of
\cref{isc:cor:parameters}, applied one normal-form level deeper.
\item \emph{Three duals of Hahn vector spaces}
(\rpath{surcomplex/three-duals-of-hahn-vector-spaces}, prefix
\texttt{duals:}). Parts~(i)--(ii) of \cref{isc:cp:thm:rank} are the
injectivity paragraph before \lbl{duals:prop:algebraic} and that
proposition, for the $k$-vector space $V=F$: the scalar extension
$K\otimes_kV$, $K=k((t^\Gamma))$, embeds in $V_\Gamma=F((t^\Gamma))$ with
image the vectors of finite total coefficient rank. Source~02 does not cite
that report and presents the tensor fact as standard. Part~(iii), the
fraction-field test, is new but immediate from~(ii).
\item \emph{Tail spans and differential transcendence}
(\rpath{surreal/tail-spans-and-differential-transcendence}, prefix
\texttt{tail:}). \lbl{tail:thm:surreal} uses the same binary-prefix almost
disjoint family as \cref{isc:cp:cor:continuum}, private indices and a
Vandermonde argument to give $\cc$ differentially independent series. The
base ($\Q((t^\Q))$), the derivation (Berarducci--Mantova, via Euler jets)
and the mechanism (the tail-span theorem \lbl{tail:thm:tail} for
multiquadratic coefficients) differ from those of
\cref{isc:cp:thm:lacunary,isc:cp:thm:generic}, whose base is a compositum
or a finite-generation envelope and whose tool is coefficient degree or a
coefficient derivation. Source~02 does not cite it; the two are related,
not duplicates.
\item \emph{Bounded-support transcendence}
(\rpath{surreal/transcendence-over-bounded-support}, prefix \texttt{bst:}).
\lbl{bst:thm:independent} obtains $2^{\cf(G)}$ independent series over the
bounded-support fraction field by coefficient codes realized at cofinally
many indices, a private-index device over a different base, with a
support-gap estimate instead of a derivation. Source~02 does not cite it.
\item \emph{Dilation rigidity}
(\rpath{surcomplex/autonomous-dilation-relations}, prefix \texttt{adr:}).
\lbl{adr:sr:thm:tail} localizes a set of surcomplex parameters in the Hahn
field of the real span of their supports and finds a tail of an explicit
ordinal family whose dilation orbits are independent over that whole field;
\lbl{adr:sr:thm:kappa} gives set-sized families over a prescribed set base.
\Cref{isc:cp:thm:resilience} has the same shape, with Euler jets instead of
dilation orbits and localization at the inner normal-form level.
\item \emph{Omnific-preserving automorphisms} (\texttt{opa:}). The
explicit copies are the maps $M_{1,\iota_{\sigma_j}}$ of that report's
Part~III (\cref{isc:sec:collection}). Research questions~\ref{isc:cp:q:bm}
and~\ref{isc:cp:q:exp} lie inside \lbl{opa:as:q:exp}.
\end{enumerate}
No \texttt{isc:cp:} label has a Lean mapping in \rpath{FORMALIZATION.md}.

\section{Conventions and foundational lemmas of the addition}
\label{isc:cp:sec:foundations}

\subsection{Two support conventions}

For a field $F$ and an ordered abelian group $\Gamma$, the Hahn field
$F((t^\Gamma))$ consists of formal sums
\[
 x=\sum_{\gamma\in\Gamma}x_\gamma t^\gamma
\]
whose nonzero coefficient support is a well-ordered set. All supports are
sets. Addition is coefficientwise. In a product, the coefficient at
$\gamma$ is the finite sum of $x_\alpha y_\beta$ over pairs with
$\alpha+\beta=\gamma$.

Surreal normal forms instead use the large-$\omega$ convention
$x=\sum_g c_g\omega^g$, with reverse-well-ordered support: every nonempty
subset has a greatest element. As in \cref{isc:sec:foundations}, $\K{k}{G}$
is the full field of these normal forms with coefficients in $k$ and
exponents in an additive subgroup $G\subseteq\No$. Whenever the conventions
change, the identification is
\begin{equation}\label{isc:cp:eq:signconvention}
 t^\gamma=\omega^{-\gamma},
\end{equation}
as in \cref{isc:sub:bounded}. The reversal in
\eqref{isc:cp:eq:signconvention} is essential in the final omnific
construction.

For set-sized groups, ordinary Hahn field algebra suffices. For
proper-class exponent groups, each finite calculation occurs in the
subgroup generated by a set of supports. The inverse of a nonzero series
can likewise be constructed in a set-sized Hahn field containing its
support. Thus the formulas define class fields without allowing class-sized
sums. The class assertions are formulated in NBG with global choice, with
ZFC for sets, as in \cref{isc:sec:foundations}. The displayed class families
and derivations are given by formulas; no class-state recursion, class
transcendence basis, or class satisfaction predicate is used.

\subsection{Euler derivations and differential independence}

\begin{definition}\label[definition]{isc:cp:def:euler}
Let $k\subseteq F$ have characteristic zero and let $\lambda:\Gamma\to(k,+)$
be additive. The associated outer Euler derivation is
\begin{equation}\label{isc:cp:eq:euler}
 \partial\Bigl(\sum_\gamma a_\gamma t^\gamma\Bigr)
 =\sum_\gamma \lambda(\gamma)a_\gamma t^\gamma.
\end{equation}
It kills $F$. When $\Gamma=\Z$ and $\lambda(n)=n$, it is $t\,d/dt$.
\end{definition}

\begin{lemma}\label[lemma]{isc:cp:lem:derivation}
Formula \eqref{isc:cp:eq:euler} is a derivation of $F((t^\Gamma))$. It
preserves Hahn-summable families and their sums. If $d:F\to F$ is a
$k$-derivation, its coefficientwise extension to $F((t^\Gamma))$ is also a
derivation and commutes with $\partial$.
\end{lemma}
\begin{proof}
Both maps only delete terms or change coefficients, so the support
condition is preserved. For a product coefficient, finiteness of the
convolution and $\lambda(\alpha+\beta)=\lambda(\alpha)+\lambda(\beta)$ give
the Leibniz identity. The same finite calculation works for $d$. A
Hahn-summable family has a well-ordered union of supports and only finitely
many nonzero contributions at each exponent; neither property is destroyed.
Finally $d(\lambda(\gamma))=0$ because $\lambda(\gamma)\in k$, proving
commutation coefficient by coefficient.
\end{proof}

\begin{definition}\label[definition]{isc:cp:def:diffind}
A family $(y_i)_{i\in I}$ in a differential field extension $E/M$ is
\emph{differentially algebraically independent over $M$} if the family
\[
 \jet(y_i:i\in I)=\{\partial^j y_i:i\in I,\ j\ge0\}
\]
is algebraically independent over $M$. Every relation and every collection
of variables in a polynomial are finite. We write $M\langle P\rangle$ for
the field generated by $M$ and all $\partial^j p$, $p\in P$, $j\ge0$.
\end{definition}

For set fields, differential transcendence degree has its usual meaning.
All lower bounds below are proved by exhibiting independent families of the
stated cardinality. For class fields, only the finite-relation definition
and set-indexed subfamilies are used; no class transcendence basis is
postulated, and no cardinal equal to the size of $\No$.

\subsection{The coefficient-field envelope}

Throughout the set-sized theory, put
\begin{equation}\label{isc:cp:eq:basicfields}
 E=F((t^\Gamma)),\qquad L=k((t^\Gamma)),\qquad M=FL.
\end{equation}
The compositum is taken in $E$ and is a field. In contrast, the ring
$F[L]$ generated by $F$ and $L$ consists of finite sums of products.

\begin{definition}\label[definition]{isc:cp:def:cfg}
Define the finite-generation coefficient envelope by
\begin{equation}\label{isc:cp:eq:cfg}
 \Cfg(F/k;\Gamma)
 =\bigcup_{\substack{k\subseteq D\subseteq F\\D/k\ \mathrm{finitely\ generated}}}
          D((t^\Gamma)).
\end{equation}
Here ``finitely generated'' means as a field, not as a ring.
\end{definition}

\begin{lemma}\label[lemma]{isc:cp:lem:cfg}
The union in \eqref{isc:cp:eq:cfg} is a field containing $M$ and is stable
under every Euler derivation \eqref{isc:cp:eq:euler}. Every finite subset of
$M$ is contained in one $D((t^\Gamma))$ with $D/k$ finitely generated.
\end{lemma}
\begin{proof}
The compositum of two finitely generated coefficient fields is again
finitely generated, so the union is directed. Each constituent is a Hahn
field. A finite field expression in elements of $F$ and $L$ uses only
finitely many elements $a_1,\ldots,a_r$ of $F$, and therefore belongs to
$k(a_1,\ldots,a_r)((t^\Gamma))$. The same argument handles finitely many
expressions simultaneously. Formula \eqref{isc:cp:eq:euler} introduces only
scalars from $k$ into coefficients.
\end{proof}

The last lemma is the finite coefficient-field control of
\cref{isc:lem:finitecoeff}, in the abstract setting. The later constructions
show both strictness of its necessary condition and independence over the
entire envelope.

\section{An exact coefficient-rank test for the compositum}
\label{isc:cp:sec:rank}

\begin{definition}\label[definition]{isc:cp:def:rank}
For $x=\sum x_\gamma t^\gamma\in E$, define
\[
 \crk_k(x)=\dim_k\operatorname{span}_k\{x_\gamma:\gamma\in\Gamma\}.
\]
In particular, $\crk_k(0)=0$. Let $F[L]\subseteq E$ be the ring generated
by $F$ and $L$.
\end{definition}

\begin{theorem}[Finite rank and finite-rank denominators]\label{isc:cp:thm:rank}
For arbitrary fields $k\subseteq F$ and any ordered abelian group $\Gamma$:
\begin{enumerate}[label=\textup{(\roman*)}]
\item The natural multiplication map $F\otimes_k L\to E$ is injective, with
image $F[L]$.
\item $F[L]=\{x\in E:\crk_k(x)<\infty\}$.
\item $M=\Frac(F[L])$, and an element $x\in E$ belongs to $M$ exactly when
there is $q\in E\setminus\{0\}$ for which both $\crk_k(q)$ and
$\crk_k(qx)$ are finite.
\end{enumerate}
\end{theorem}
\begin{proof}
A finite tensor can be written as $\sum_{i=1}^r b_i\otimes\ell_i$ with the
$b_i\in F$ linearly independent over $k$. If its image vanishes, then
comparison at each exponent $\gamma$ gives
\[
 \sum_{i=1}^r b_i(\ell_i)_\gamma=0.
\]
All $(\ell_i)_\gamma$ belong to $k$, so each is zero. Thus every $\ell_i=0$,
proving injectivity. The image is the finite sums of products of elements of
$F$ and $L$, namely $F[L]$.

For $x=\sum_{i=1}^r b_i\ell_i$, every coefficient of $x$ lies in the
finite-dimensional space spanned by the $b_i$. Conversely, suppose the
coefficient span has basis $b_1,\ldots,b_r$. Write uniquely
\[
 x_\gamma=\sum_{i=1}^r b_i c_{i,\gamma},\qquad c_{i,\gamma}\in k.
\]
For every $i$, the support of $\ell_i=\sum_\gamma c_{i,\gamma}t^\gamma$ is
contained in $\supp(x)$, hence is well ordered. Therefore $\ell_i\in L$ and
$x=\sum b_i\ell_i\in F[L]$.

Since $F[L]$ is a subring of a field, its fraction field inside $E$ is
precisely the smallest field containing $F$ and $L$, which is $M$. Finally
$x=p/q$ with $p,q\in F[L]$, $q\ne0$, is equivalent to $q,qx\in F[L]$ for
such a $q$. Part~(ii) gives the asserted test.
\end{proof}

\mtag{} Parts~(i) and~(ii) are, for the $k$-vector space $V=F$, the
injectivity paragraph before \lbl{duals:prop:algebraic} of the three-duals
report and that proposition (\cref{isc:cp:sec:collection}); the source does
not cite it and presents the tensor fact as standard.

\begin{remark}[What the criterion does and does not do]\label[remark]{isc:cp:rem:ranktest}
The theorem describes the subfield exactly in terms of its coefficients,
rather than merely as an unspecified compositum. Nevertheless, the
existential search over $q$ is an infinite-input problem. No algorithm
deciding arbitrary Hahn-series membership follows. The tensor argument is a
standard linear-disjointness mechanism; neither a new tensor-product theorem
nor a new general theory of rational series is claimed. The comparison with
recurrence criteria in \cite{KKS} is conceptual, not a claim that the two
formulations have identical inputs.
\end{remark}

\begin{example}[Infinite rank inside the compositum]\label[example]{isc:cp:ex:geometric}
Let $u\in F$ be transcendental over $k$ and let $g\in\Gamma$ be positive.
The geometric series
\[
 z=\frac1{1-ut^g}=\sum_{n\ge0}u^n t^{ng}
\]
has infinite coefficient rank, because the powers of $u$ are linearly
independent over $k$. Yet $z\in M$: the multiplier $q=1-ut^g$ has rank two
and $qz=1$ has rank one. The support $\{ng:n\ge0\}$ is well ordered in every
ordered abelian group. No cofinality assumption on $g$ is needed.
\end{example}

\begin{corollary}[Finite coefficient extensions]\label[corollary]{isc:cp:cor:finitecoeff}
If $[F:k]<\infty$, then $E=F[L]=M$.
\end{corollary}
\begin{proof}
Every coefficient span is a subspace of the finite-dimensional $k$-space $F$.
Apply \cref{isc:cp:thm:rank}.
\end{proof}

For finite ranks one also has
\[
 \crk_k(x+y)\le\crk_k(x)+\crk_k(y),\qquad
 \crk_k(xy)\le\crk_k(x)\crk_k(y).
\]
Indeed a coefficient of the product lies in the span of pairwise products
of basis elements for the two coefficient spaces. Infinite rank therefore
obstructs membership in $F[L]$, but not in its fraction field.

\section{Linear disjointness under two enlargements}\label{isc:cp:sec:disjoint}

The next lemmas prevent an independence proof over a small coefficient
field or a cyclic exponent group from being lost when the ambient Hahn
field is enlarged. The arguments are included instead of silently treating
base extension as harmless. Linear disjointness has its usual definition
via injectivity of the tensor multiplication map \cite{StacksLD}, as in
\cref{isc:sec:foundations}.

\begin{lemma}[Coefficient extension]\label[lemma]{isc:cp:lem:coeffld}
If $F_0\subseteq F$, then $F$ and $F_0((t^\Gamma))$ are linearly disjoint
over $F_0$ inside $F((t^\Gamma))$.
\end{lemma}
\begin{proof}
A finite relation with coefficients in $F_0((t^\Gamma))$ among elements of
$F$ linearly independent over $F_0$ vanishes coefficientwise only when every
coefficient series is zero. This is the first argument of
\cref{isc:cp:thm:rank} with $k=F_0$.
\end{proof}

\begin{lemma}[Intermediate base change]\label[lemma]{isc:cp:lem:basechange}
Suppose fields $A$ and $B$ are linearly disjoint over $K$ in an ambient
field, and $K\subseteq A_0\subseteq A$. Then $A$ and $A_0B$ are linearly
disjoint over $A_0$.
\end{lemma}
\begin{proof}
Let $a_1,\ldots,a_m\in A$ be linearly independent over $A_0$ and suppose
$\sum_i a_i c_i=0$ with $c_i\in A_0B$. Clear a common nonzero denominator
from the ring generated by $A_0$ and $B$. Each resulting coefficient is a
finite sum of products. Choose a $K$-linearly independent list
$b_1,\ldots,b_r\in B$ spanning all the finitely many $B$-entries used. The
relation becomes
\[
 \sum_{j=1}^r\Bigl(\sum_{i=1}^m a_i a_{ij}\Bigr)b_j=0,
 \qquad a_{ij}\in A_0.
\]
Linear disjointness over $K$ makes the $b_j$ linearly independent over $A$.
Hence every inner sum vanishes. Independence of the $a_i$ over $A_0$ gives
all $a_{ij}=0$, so all the original $c_i$ are zero.
\end{proof}

\begin{lemma}[Coset slices]\label[lemma]{isc:cp:lem:coset}
Let $\Gamma_0\subseteq\Gamma$ be an additive subgroup. Inside
$F((t^\Gamma))$, the fields
\[
 E_0=F((t^{\Gamma_0})),\qquad L=k((t^\Gamma))
\]
are linearly disjoint over $L_0=k((t^{\Gamma_0}))$. No convexity assumption
on $\Gamma_0$ is required.
\end{lemma}
\begin{proof}
Suppose $u_1,\ldots,u_m\in E_0$ are linearly independent over $L_0$ and
$\sum_i u_i v_i=0$, $v_i\in L$. For a coset $C$ of $\Gamma_0$ meeting one of
the supports of the $v_i$, choose a representative $c\in C$. Let $v_{i,C}$
be the sum of the terms of $v_i$ supported on $C$, multiplied by $t^{-c}$.
Its support is a translate of a subset of a well-ordered set, and is
contained in $\Gamma_0$, so $v_{i,C}\in L_0$.

Multiplication by a series supported in $\Gamma_0$ does not change the
coset of an exponent. Extracting the $C$-slice of the original relation and
translating by $-c$ therefore gives
\[
 \sum_i u_i v_{i,C}=0.
\]
Independence over $L_0$ yields $v_{i,C}=0$ for all $i$. Every support
exponent belongs to a coset that was considered; hence every $v_i=0$. Only a
set of cosets is involved, even in the class-supported setting.
\end{proof}

\mtag{} Up to the change of convention \eqref{isc:cp:eq:signconvention},
$v_{i,C}$ is a normalized coset slice as in \eqref{isc:eq:slice}, and the
lemma is a variant of \cref{isc:thm:sliceddisjoint} with different
coefficient fields on the two sides (\cref{isc:cp:sec:collection}).

\begin{corollary}[Independence survives both enlargements]\label[corollary]{isc:cp:cor:transfer}
Let $x\in F$ be transcendental over $k$, let $g>0$ in $\Gamma$, and let
$\lambda:\Gamma\to k$ satisfy $\lambda(g)=1$. If a family in $k(x)((t))$
is differentially algebraically independent over $k((t))(x)$ for
$t\,d/dt$, then replacing $t^n$ by $t^{ng}$ gives a family differentially
algebraically independent over $F k((t^\Gamma))$ in $F((t^\Gamma))$.
\end{corollary}
\begin{proof}
First apply \cref{isc:cp:lem:coeffld} to $k(x)\subseteq F$ and the cyclic
exponent group. Then apply \cref{isc:cp:lem:basechange} with
$A=k(x)((t))$ and $A_0=k((t))(x)$: the field $A$ is linearly disjoint from
$F k((t))$ over $A_0$. Distinct monomials in any finite collection of jets,
which were linearly independent over $A_0$, remain so over $F k((t))$.

Next identify this cyclic Hahn field with $F((t^{\Z g}))$. Apply
\cref{isc:cp:lem:coset}, followed by \cref{isc:cp:lem:basechange} with the
intermediate field $F k((t^{\Z g}))$. This preserves independence after
adjoining $k((t^\Gamma))$. Finally $\partial(t^{ng})=n t^{ng}$, so the
transported elements have exactly the transported jets.
\end{proof}

\section{Differential independence from one coefficient variable}
\label{isc:cp:sec:lacunary}

The next theorem is the key to the small-transcendence-degree case. The
exponents in the outer variable need not be sparse. Instead, the degrees in
a single coefficient variable grow so fast that a late coefficient cannot be
canceled by any fixed differential polynomial.

\begin{definition}\label[definition]{isc:cp:def:private}
A family $\mathfrak A$ of infinite subsets of $\N_{>0}$ has the \emph{finite
private-index property} if, for every finite subfamily $A_1,\ldots,A_r$ of
distinct members and every $i$, the set
\[
 A_i\setminus\bigcup_{j\ne i}A_j
\]
is infinite. In particular, every pairwise almost disjoint family of
infinite subsets has this property.
\end{definition}

\begin{theorem}[Rapid coefficient-degree independence]\label{isc:cp:thm:lacunary}
Let $k$ have characteristic zero, let $x$ be transcendental over $k$, and
let $(d_n)_{n\ge1}$ be a strictly increasing sequence of positive integers
such that
\begin{equation}\label{isc:cp:eq:fastgrowth}
 \frac{d_n}{d_{n-1}}\longrightarrow+\infty.
\end{equation}
For a family $\mathfrak A$ with the finite private-index property, put
\begin{equation}\label{isc:cp:eq:fA}
 f_A=\sum_{n\in A}x^{d_n}t^n\in k(x)((t)).
\end{equation}
Then $(f_A)_{A\in\mathfrak A}$ is differentially algebraically independent
over $k((t))(x)$ for $\partial=t\,d/dt$ and $\partial x=0$.
\end{theorem}
\begin{proof}
Suppose there is a relation in finitely many jets of distinct
$f_{A_1},\ldots,f_{A_r}$. Among all nonzero polynomial relations over
$k((t))(x)$ in finite lists of these jets, choose one of least total jet
degree $d$. Write it as
\[
 P\bigl((\partial^j f_{A_i})_{1\le i\le r,\,0\le j\le m}\bigr)=0,
\]
allowing unused variables. A degree-zero nonzero relation is impossible, so
$d\ge1$.

Clear the denominators in $x$ of the finitely many coefficients of $P$,
then multiply by a power of $t$ large enough to remove their negative outer
exponents. We may therefore assume
\begin{equation}\label{isc:cp:eq:integralP}
 P\in k[[t]][x][Y_{ij}:1\le i\le r,\ 0\le j\le m].
\end{equation}
Let $\nu$ be its degree in $x$. These operations do not change total jet
degree or minimality. We use $\mathbf f$ to denote the tuple of all the
indicated jets.

For each variable set
\[
 p_{ij}=P_{Y_{ij}}(\mathbf f)\in k[x][[t]].
\]
If the formal polynomial $P_{Y_{ij}}$ is nonzero, its total jet degree is
smaller than $d$, so minimality forces $p_{ij}\ne0$. At least one formal
partial is nonzero, because the characteristic is zero and $d\ge1$.
Consequently
\[
 \ell=\min_{p_{ij}\ne0}\ordt(p_{ij})
\]
is a nonnegative integer. Define polynomials in a new indeterminate $X$ by
\begin{equation}\label{isc:cp:eq:Hpolynomial}
 W_i(X,x)=\sum_{j=0}^m [t^\ell]p_{ij}\,X^j\in k[x][X].
\end{equation}
At least one $W_i$ is nonzero: its coefficients occupy distinct powers of
$X$, and at least one $p_{ij}$ attains the valuation $\ell$. Fix such an
index $i_*$.

For an integer $n$, truncate every original series strictly before $t^n$:
\[
 v_i=\sum_{\substack{a\in A_i\\a<n}}x^{d_a}t^a,
 \qquad w_i=f_{A_i}-v_i.
\]
Every jet of $w_i$ has $t$-valuation at least $n$. Let $\mathbf v$ and
$\mathbf w$ be the corresponding tuples of jets. Evaluating any polynomial
from \eqref{isc:cp:eq:integralP} at $\mathbf f$ or $\mathbf v$ changes it only
by a multiple of $t^n$. Therefore, when $n>\ell$, the partial evaluations
$P_{Y_{ij}}(\mathbf v)$ have zero coefficients below $\ell$ and the same
coefficient at $\ell$ as $p_{ij}$.

Choose $n$ in the infinite private-index set
\[
 A_{i_*}\setminus\bigcup_{i\ne i_*}A_i.
\]
We may require simultaneously
\begin{equation}\label{isc:cp:eq:nchoices}
 n>\max\{\ell,1\},\qquad W_{i_*}(n,x)\ne0,\qquad
 d_n>\nu+d\,d_{n-1}.
\end{equation}
The middle condition excludes only finitely many integers, since a nonzero
polynomial in $X$ over the field $k(x)$ has finitely many roots and
characteristic zero gives distinct integer elements. The last condition
holds eventually by \eqref{isc:cp:eq:fastgrowth}.

The polynomial Taylor expansion about $\mathbf v$ is
\begin{equation}\label{isc:cp:eq:taylor}
 P(\mathbf v+\mathbf w)
 =P(\mathbf v)+\sum_{i,j}P_{Y_{ij}}(\mathbf v)\partial^jw_i
       +\text{terms of tail degree at least two}.
\end{equation}
Every omitted term is divisible by $t^{2n}$. Since $n+\ell<2n$, these terms
do not contribute at exponent $n+\ell$. In the linear sum, a tail
coefficient after exponent $n$ could contribute only by multiplying a
partial coefficient before $\ell$, and those coefficients vanish. At
exponent $n$, only $w_{i_*}$ is nonzero, and
\[
 [t^n]\partial^jw_{i_*}=n^j x^{d_n}.
\]
Taking $[t^{n+\ell}]$ in \eqref{isc:cp:eq:taylor}, using $P(\mathbf f)=0$,
gives the exact polynomial identity
\begin{equation}\label{isc:cp:eq:degreecontradiction}
 0=[t^{n+\ell}]P(\mathbf v)+x^{d_n}W_{i_*}(n,x).
\end{equation}

Every coefficient of every truncated jet has degree in $x$ at most
$d_{n-1}$. Thus the first summand in \eqref{isc:cp:eq:degreecontradiction}
has degree at most $\nu+d\,d_{n-1}$. The second is nonzero and has degree
at least $d_n$. The final inequality in \eqref{isc:cp:eq:nchoices} makes
cancellation impossible. This contradiction proves the theorem.
\end{proof}

\begin{remark}[Why arbitrary Laurent coefficients are allowed]\label[remark]{isc:cp:rem:laurentcoeff}
The coefficients of $P$ can be arbitrary elements of $k((t))(x)$, not just
rational functions of $t$ and $x$ over $k$. Clearing denominators makes
their degree in $x$ uniformly finite, while their $t$-series can remain
infinite. The proof uses only this finite $x$-degree and a common lower
bound on outer exponents. It does not impose a recurrence on those
coefficient series.
\end{remark}

\begin{corollary}[A continuum family in one coefficient field]\label[corollary]{isc:cp:cor:continuum}
For every characteristic-zero field $k$ and transcendental $x$, the field
$k(x)((t))$ contains $\cc$ elements differentially algebraically
independent over $k((t))(x)$, all with coefficients in $k(x)$.
\end{corollary}
\begin{proof}
Take $d_n=n!$. Encode the nonempty finite binary strings injectively by
positive integers. For each infinite binary sequence $\eta$, let $A_\eta$
consist of the labels of its finite nonempty prefixes. Distinct branches
have only finitely many common prefixes, so their sets are pairwise almost
disjoint. There are $\cc$ branches. Apply \cref{isc:cp:thm:lacunary}.
\end{proof}

\mtag{} The binary-prefix family is the one of \lbl{tail:eq:Aalpha} in the
tail-spans report, where it serves \lbl{tail:thm:surreal} over a different
base and with a different derivation (\cref{isc:cp:sec:collection}).

\begin{corollary}[Failure of the finite coefficient-field criterion]\label[corollary]{isc:cp:cor:fginsufficient}
Suppose $F/k$ has positive transcendence degree and there are $g>0$ in
$\Gamma$ and additive $\lambda:\Gamma\to k$ with $\lambda(g)=1$. Then
$\Cfg(F/k;\Gamma)$ contains an element differentially transcendental over
$M=F k((t^\Gamma))$. In particular, containment of all coefficients in a
single finitely generated field does not imply membership in $M$, or even
algebraicity over $M$.
\end{corollary}
\begin{proof}
Choose a transcendental $x\in F$ and use
\[
 f=\sum_{n\ge1}x^{n!}t^{ng}.
\]
The singleton family $\{\N_{>0}\}$ satisfies the private-index condition.
Apply \cref{isc:cp:thm:lacunary,isc:cp:cor:transfer}. The series belongs to
$k(x)((t^\Gamma))\subseteq\Cfg(F/k;\Gamma)$.
\end{proof}

\begin{example}[The first terms and the mechanism]\label[example]{isc:cp:ex:mechanism}
The last witness starts as
\[
 xt^g+x^2t^{2g}+x^6t^{3g}+x^{24}t^{4g}+x^{120}t^{5g}+\cdots.
\]
The outer support here is the entire positive cyclic progression, not a
factorially sparse support. A polynomial differential relation of total jet
degree $d$ and coefficient degree $\nu$ sees only degree at most
$\nu+d(n-1)!$ in the truncation before $t^{ng}$. A surviving linear response
to the next coefficient has degree at least $n!$. This is the obstruction.
Merely pointing to gaps in the outer support would not prove this example.
\end{example}

\mtag{} The source's table in \cref{isc:cp:sec:audit} applies
\cref{isc:cp:cor:fginsufficient} to the explicit join of \cref{isc:sec:join},
with $F=K_0$, $k=\R$, $\Gamma=G_1$ and $\lambda$ as in
\cref{isc:cp:sub:euler}; the lemmas of \cref{isc:cp:sec:disjoint} use only
finite relations, so the class-sized $F$ causes no difficulty. The example
is then a series with every coefficient in the one-generator field
$\R(x)$, $x\in K_0$ transcendental over $\R$. \Cref{isc:lem:finitecoeff}
places every finite subset of $M$ in some $\K{D}{G_1}$ with $D$ finitely
generated; the example shows that the converse fails for a single
generator.

\section{Independent coefficient blocks and tree amplification}
\label{isc:cp:sec:generic}

The preceding theorem uses one coefficient variable. The next results
exploit many algebraically independent coefficients and obtain independence
over a larger base than the compositum itself.

\subsection{Coefficient derivations with finite avoidance}

\begin{lemma}[Avoiding a finitely generated coefficient field]\label[lemma]{isc:cp:lem:partial}
Let $F/k$ have characteristic zero, let $(a_b)_{b\in B}\subseteq F$ be
algebraically independent over $k$, and let $D\subseteq F$ be finitely
generated over $k$. There is a finite set $B_0\subseteq B$ such that, for
every $b\in B\setminus B_0$, a $k$-derivation $d_b:F\to F$ exists with
\[
 d_b(a_b)=1,\qquad d_b(a_c)=0\ (c\ne b),\qquad d_b(D)=0.
\]
\end{lemma}
\begin{proof}
Extend the given independent set to a transcendence basis $\mathcal T$ of
the set field $F/k$. Each member of a finite generating list for $D/k$ is
algebraic over $k(\mathcal T_0)$ for some finite
$\mathcal T_0\subseteq\mathcal T$; take one such finite set for the whole
list. Let $B_0$ consist of the $b$ with $a_b\in\mathcal T_0$.

On the rational function field $k(\mathcal T)$ define the coordinate
derivation that sends $a_b$ to $1$ and all other basis elements to $0$. It
extends uniquely over the algebraic extension $F/k(\mathcal T)$, because
this extension is separable. Explicitly, for an algebraic element $z$ with
minimal polynomial $p(X)$, differentiation of $p(z)=0$ forces
\begin{equation}\label{isc:cp:eq:derivextension}
 d_b(z)=-\frac{(d_b p)(z)}{p'(z)}.
\end{equation}
The denominator is nonzero. This formula gives the usual compatible
extension through finite algebraic subextensions and hence through their
union. If $b\notin B_0$, the derivation kills $k(\mathcal T_0)$, so
\eqref{isc:cp:eq:derivextension} makes it kill each of the chosen generators
of $D$. It therefore kills $D$.
\end{proof}

\subsection{The block theorem}

\begin{theorem}[Private coefficient blocks]\label{isc:cp:thm:generic}
Let $F/k$ have characteristic zero. Let $(a_b)_{b\in B}\subseteq F$ be
algebraically independent over $k$, and choose $g>0$ in $\Gamma$ and
additive $\lambda:\Gamma\to k$ with $\lambda(g)=1$. For each $i\in I$, let
\[
 b_i:\N_{>0}\longrightarrow B
\]
be injective. Assume that for every finite set of distinct indices
$i_1,\ldots,i_r$ and every $i=i_q$, there are infinitely many $n$ such that
the label $b_i(n)$ occurs in no other sequence $b_{i_p}$, at any position.
Put
\begin{equation}\label{isc:cp:eq:genericseries}
 S_i=\sum_{n\ge1}a_{b_i(n)}t^{ng}.
\end{equation}
Then $(S_i)_{i\in I}$ is differentially algebraically independent over
$\Cfg(F/k;\Gamma)$ for the Euler derivation associated to $\lambda$.
\end{theorem}
\begin{proof}
Every displayed series has well-ordered support contained in
$\{ng:n\ge1\}$. Suppose a nonzero polynomial relation exists among finitely
many jets of finitely many distinct $S_i$, with coefficients in
$\Cfg(F/k;\Gamma)$. Choose a relation $P$ of least total jet degree. Its
coefficients all lie in one $D((t^\Gamma))$ with $D/k$ finitely generated.
The total degree is positive.

Fix one index $i$ involved in the relation. By the private-label hypothesis,
there are infinitely many positions $n$ whose coefficient label occurs only
in this series among those under consideration. After excluding the finite
exceptional set of \cref{isc:cp:lem:partial}, infinitely many remain. For
such a position, extend $d_{b_i(n)}$ coefficientwise to the Hahn field. By
\cref{isc:cp:lem:derivation}, it kills every coefficient of $P$, commutes
with $\partial$, and satisfies
\[
 d_{b_i(n)}(\partial^jS_i)=n^j t^{ng},\qquad
 d_{b_i(n)}(\partial^jS_{i'})=0\quad(i'\ne i).
\]
Differentiating $P=0$ by this coefficient derivation and dividing by the
nonzero monomial $t^{ng}$ gives
\begin{equation}\label{isc:cp:eq:vandermonde}
 \sum_{j=0}^m n^j P_{Y_{ij}}\bigl((\partial^aS_{i'})_{i',a}\bigr)=0.
\end{equation}
Here $m$ bounds all jet orders used. The coefficients in
\eqref{isc:cp:eq:vandermonde} do not depend on $n$. Since it holds for
infinitely many distinct integers in a characteristic-zero field, the
polynomial in $n$ is identically zero. Thus every evaluated partial
$P_{Y_{ij}}$ is zero. Repeat for each index $i$ in the relation.

A nonconstant polynomial in characteristic zero has some nonzero formal
partial derivative. That partial has smaller total jet degree than $P$, yet
its evaluation vanishes. This contradicts minimality. Therefore no relation
exists.
\end{proof}

\begin{remark}[A finite Vandermonde version]\label[remark]{isc:cp:rem:vandermonde}
For a relation involving jets only through order $m$, it is enough to use
$m+1$ distinct private positions outside the finite exceptional set. The
matrix $(n_a^j)_{0\le a,j\le m}$ has determinant
$\prod_{a<b}(n_b-n_a)\ne0$. The polynomial-root argument in the proof is
simply a convenient way to avoid choosing a particular finite list.
\end{remark}

\begin{corollary}[Tree amplification]\label[corollary]{isc:cp:cor:tree}
If $\trdeg(F/k)\ge\kappa$ for an infinite cardinal $\kappa$, and
$\Gamma,g,\lambda$ satisfy the preceding hypotheses, then $F((t^\Gamma))$
contains a family of cardinality $\kappa^{\aleph_0}$ differentially
algebraically independent over $\Cfg(F/k;\Gamma)$.
\end{corollary}
\begin{proof}
The set of nonempty finite sequences in $\kappa$ has cardinality $\kappa$.
Choose an algebraically independent coefficient $a_s$ for every such finite
sequence $s$. For each branch $\eta\in\kappa^{\N_{>0}}$, use
\[
 S_\eta=\sum_{n\ge1}a_{\eta|n}t^{ng}.
\]
Two distinct branches share only finitely many prefix labels. Relative to
any finite collection of other branches, every sufficiently long prefix is
private. Apply \cref{isc:cp:thm:generic}.
\end{proof}

\subsection{Local use for class coefficient fields}\label{isc:cp:sub:classlocal}

\begin{proposition}[Set-local form]\label[proposition]{isc:cp:prop:local}
The private coefficient-block argument remains valid when $F$ or $\Gamma$
is a class, provided that $k$ is a set field, every series has set support,
and algebraic independence of the coefficient family means independence of
every finite subfamily. The conclusion concerns every finite polynomial
relation. No class transcendence basis is required.
\end{proposition}
\begin{proof}
Fix a proposed relation. It involves finitely many $S_i$, so the union of
their coefficient labels is a set. Its finitely many coefficient series
belong to $D((t^\Gamma))$ for a single set field $D$ finitely generated over
$k$. Form the set field $F'$ generated by $D$, $k$, and all the
coefficients in the finitely many chosen $S_i$. The subgroup $\Gamma'$
generated by $g$ and all supports of the relation coefficients is a set.
Restrict $\lambda$ to $\Gamma'$.

Every expression in the relation lies in $F'((t^{\Gamma'}))$. The
algebraically independent coefficient family used there is a set, so
\cref{isc:cp:lem:partial} applies to $F'/k$ in ordinary set theory. All
differentiations, finite convolution identities, and the contradiction now
occur in this set-sized field. This proves the class statement one finite
relation at a time.
\end{proof}

\begin{corollary}[Three distinct closure levels]\label[corollary]{isc:cp:cor:sandwich}
Assume $\trdeg(F/k)$ is infinite and that $\Gamma$ admits $g>0$ and
$\lambda$ with $\lambda(g)=1$. Then
\begin{equation}\label{isc:cp:eq:sandwich}
 F[L]\subsetneq M\subsetneq\Cfg(F/k;\Gamma)\subsetneq E.
\end{equation}
Moreover, the second strict inclusion has a differential-transcendental
witness over $M$, and the third has one over $\Cfg(F/k;\Gamma)$.
\end{corollary}
\begin{proof}
The first strictness is \cref{isc:cp:ex:geometric}, the second is
\cref{isc:cp:cor:fginsufficient}, and the third follows from a single
sequence of algebraically independent coefficients in
\cref{isc:cp:thm:generic}.
\end{proof}

\section{Exact transcendence degrees for Laurent-series composita}
\label{isc:cp:sec:cardinal}

\begin{theorem}[Maximal Laurent transcendence]\label{isc:cp:thm:cardinal}
Let $k\subseteq F$ be set fields of characteristic zero. Suppose
\[
 |k|\le\cc=2^{\aleph_0},\qquad \trdeg(F/k)>0.
\]
Inside $E=F((t))$, put $M=F k((t))$ and use $\partial=t\,d/dt$, acting
trivially on $F$. Then
\begin{equation}\label{isc:cp:eq:cardinalmain}
 \boxed{\quad\dtrdeg(E/M)=\trdeg(E/M)=|F|^{\aleph_0}.\quad}
\end{equation}
No continuum hypothesis or other cardinal-arithmetic hypothesis is imposed.
\end{theorem}
\begin{proof}
Write $\mu=|F|$, an infinite cardinal. The set of Laurent series has
cardinality $\mu^{\aleph_0}$: a series is specified by an integer lower
bound and a sequence of coefficients, and arbitrary coefficient sequences
already occur in $F[[t]]$. Thus
\[
 \dtrdeg(E/M)\le\trdeg(E/M)\le\mu^{\aleph_0}.
\]
The first inequality also follows directly from the definition: the zeroth
jets of a differentially independent family are algebraically independent.

If $\mu\le\cc$, then
\[
 \cc\le\mu^{\aleph_0}\le\cc^{\aleph_0}=\cc,
\]
so $\mu^{\aleph_0}=\cc$. Choose a transcendental $x\in F$. The continuum
family from \cref{isc:cp:cor:continuum} remains differentially independent
over $M$ by \cref{isc:cp:cor:transfer}, with $\Gamma=\Z$ and $g=1$. This
reaches the upper bound.

If $\mu>\cc$, let $\mathcal T$ be a transcendence basis for $F/k$. A field
algebraic over $k(\mathcal T)$ has cardinality
\[
 \max\{|k|,|\mathcal T|,\aleph_0\}.
\]
This follows by counting polynomials and their finitely many roots. Since
$|k|\le\cc<\mu$, it follows that $|\mathcal T|=\mu$. The tree family from
\cref{isc:cp:cor:tree}, with $\kappa=\mu$, has size $\mu^{\aleph_0}$ and is
differentially independent even over $\Cfg(F/k;\Z)$, hence over $M$. Again
it reaches the upper bound. The two cases prove both equalities.
\end{proof}

\begin{corollary}\label[corollary]{isc:cp:cor:classicalexamples}
For $k=\Q,\R,$ or $\C$ and $x$ transcendental over $k$,
\[
 \dtrdeg\bigl(k(x)((t))/k((t))(x)\bigr)
 =\trdeg\bigl(k(x)((t))/k((t))(x)\bigr)=\cc.
\]
More generally, \cref{isc:cp:thm:cardinal} applies to every
positive-transcendence-degree extension of any of these three fields,
irrespective of the cardinality of $F$.
\end{corollary}

\subsection{Boundary cases that must not be suppressed}\label{isc:cp:sub:boundary}

If $F/k$ is finite, \cref{isc:cp:cor:finitecoeff} gives $E=M$, so both
transcendence degrees are zero. An infinite algebraic coefficient extension
is not covered by either the finite-extension conclusion or
\cref{isc:cp:thm:cardinal}; the formula is not silently extended to that
case.

If $\Gamma=0$, there is no outer series variable and $E=M=F$. If $\Gamma$
is nonzero but no specified map $\lambda:\Gamma\to k$ takes a chosen
positive element to $1$, the particular differential statements here have
not been formulated for it. The purely algebraic finite-rank test still
applies. For divisible groups and characteristic-zero coefficient fields,
suitable maps can often be chosen by linear algebra over $\Q$, but the
theorems keep the needed map explicit.

Characteristic zero is substantive. In characteristic $p>0$, the standard
Euler derivation on $F((t))$ satisfies
\[
 \partial^p=\partial,
\]
because $n^p=n$ in the prime field. Every element then satisfies the
differential polynomial $Y^{(p)}-Y^{(1)}=0$. The differential-independence
conclusion is therefore impossible for that derivation in positive
characteristic. This does not settle the corresponding ordinary
algebraic-transcendence problem.

The hypothesis $|k|\le\cc$ enters only in the cardinal upper-bound matching
when the coefficient transcendence degree is small. Both independence
constructions themselves work over every characteristic-zero base field.
For larger $k$, they still give their explicit lower bounds, but it is not
claimed that the bounds always equal $|F|^{\aleph_0}$.

\section{Separated surreal copies and explicit omnific witnesses}
\label{isc:cp:sec:surreal}

The set-local results now apply to actual surreal normal forms. The
source reproduces the explicit two-copy construction of
\cref{isc:sec:explicit} to make its later proofs independent of any
unreviewed general class-embedding theorem, in particular of
\cref{isc:thm:group}. Only normal-form algebra and the elementary order
comparison are needed.

\subsection{The inner and outer levels, printed once}\label{isc:cp:sub:copies}

\mtag{} The source's Subsections~8.1 and~8.2 re-derive, with the same
proofs, material printed in \cref{isc:sec:explicit,isc:sec:join}; it is
printed once, there, and the source is credited as an independent
re-derivation. The correspondence is: the order isomorphism $\psi$ and its
surjectivity proof are \cref{isc:lem:compression}; the intervals
$I_j=(3j,3j+1)$, the inner maps $h_j$ of \eqref{isc:eq:inner}, their
non-unitality \eqref{isc:eq:innerone} and the description of
$G_j=h_j(\No)$ as all normal forms supported in $I_j$ are in
\cref{isc:sec:explicit} (with $\alpha=j\in\{0,1\}$); the outer lifts
$\Phi_j$ of \eqref{isc:eq:explicitouter} with images $K_j=\K{\R}{G_j}$ are
\cref{isc:thm:lift,isc:thm:explicit}; the scale separation is
\eqref{isc:eq:separated}; and the re-expansion is \cref{isc:lem:iterated},
whose proof the source repeats. The source adds the following, each
immediate from those statements.

Conway's omnific integers are
\begin{equation}\label{isc:cp:eq:oz}
 \Oz=\Bigl\{\sum_{e>0}r_e\omega^e+n:
           r_e\in\R,\ n\in\Z,\ \supp\text{ reverse well ordered}\Bigr\}.
\end{equation}
A \emph{purely infinite} omnific integer has constant term zero; the term
does not require every coefficient to be positive. Since $h_j$ preserves
the sign and zero of an exponent, $\Phi_j$ preserves and reflects the
condition in \eqref{isc:cp:eq:oz} (\cref{isc:prop:arithmeticlift}), so
\begin{equation}\label{isc:cp:eq:Ri}
 R_j:=\Phi_j(\Oz)=K_j\cap\Oz,\qquad R_0\cap R_1=\Z,
\end{equation}
as in \cref{isc:thm:explicit}, and $\Phi_j$ preserves and reflects purely
infinite elements. None of this says that $\Phi_j$ commutes with the unary
omega-map: $\Phi_0(\omega)=\omega^{\omega^{1/2}}\ne\omega^{\Phi_0(1)}=\omega$
(\cref{isc:warn:omega}). By \eqref{isc:eq:separated}, $G_0+G_1$ is a direct
sum as a group, ordered lexicographically with $G_1$ dominant. With
\[
 E=\K{\R}{G_0+G_1},\qquad M=K_0K_1,
\]
as in \cref{isc:main:join}, regrouping by the $G_1$-coordinate
(\cref{isc:lem:iterated}) and using $t^\gamma=\omega^{-\gamma}$ identifies
\begin{equation}\label{isc:cp:eq:iterated}
 E\simeq K_0((t^{G_1})),\qquad
 K_1\simeq\R((t^{G_1})).
\end{equation}
Here the coefficient field $K_0$ is a class, but each individual outer
series uses only a set of coefficients.

\subsection{The exact test for the explicit copies}

\begin{corollary}[The exact test for these copies]\label[corollary]{isc:cp:cor:surrealrank}
Expand $z\in E$ as a Hahn series in $G_1$ with coefficients in $K_0$. Then
$z\in M$ if and only if some nonzero $q\in E$ has both its coefficient space
and the coefficient space of $qz$ finite-dimensional over $\R$.
\end{corollary}
\begin{proof}
Apply the coefficientwise proof of \cref{isc:cp:thm:rank} in the coordinates
\eqref{isc:cp:eq:iterated}. Every tensor, finite-dimensional span, and
fraction used in that proof is set-sized even though the ambient fields are
classes.
\end{proof}

\subsection{A concrete Euler derivation}\label{isc:cp:sub:euler}

Define
\begin{equation}\label{isc:cp:eq:hb}
 g=\omega^{7/2}\in G_1,\qquad b=\omega^{15/4}\in G_1,
 \qquad \tau=\omega^{-g}.
\end{equation}
Here $\tau$ is the single cyclic monomial indexed by $g$ in the full
small-$t$ notation $t^\gamma=\omega^{-\gamma}$, which is retained only for
the full iterated Hahn field. The element $g$ is $h_1(1)$, the same $g$ as
in the proof of \cref{isc:main:join}.

Let $\lambda:G_1\to\R$ extract the coefficient of the inner monomial
$\omega^{7/2}$ from an element of $G_1$. Coefficient extraction is
additive, with
\[
 \lambda(g)=1,\qquad\lambda(b)=0.
\]
The associated derivation $\Delta$ on $E$ kills $K_0$ and, in Conway
notation, satisfies
\begin{equation}\label{isc:cp:eq:Delta}
 \Delta(\omega^{\gamma_0+\gamma_1})
  =-\lambda(\gamma_1)\omega^{\gamma_0+\gamma_1}
 \quad(\gamma_j\in G_j).
\end{equation}
In particular $\Delta\tau=\tau$ and $\Delta\omega^b=0$. It preserves both
$K_1$ and $M$. By \cref{isc:cp:lem:derivation} it is a strong derivation
for Hahn summation. It does not alter support exponents, and therefore maps
$E\cap\Oz$ into its purely infinite ideal: any constant term is killed.

This is an explicitly specified Euler derivation adapted to the separated
coordinates. It is not asserted to be a Hardy-type derivation, the
Berarducci--Mantova derivation, or compatible with surreal exponentiation
\cite{KM,BM}.

\subsection{The countably supported family}

For every ordinal $\alpha$ and integer $n\ge1$, let
\begin{equation}\label{isc:cp:eq:labels}
 \beta_{\alpha,n}=\omega\cdot\alpha+n,
 \qquad g_{\alpha,n}=\omega^{\psi(\beta_{\alpha,n})},
 \qquad a_{\alpha,n}=\omega^{g_{\alpha,n}}.
\end{equation}
The product and sum in the label $\beta_{\alpha,n}$ are ordinary
\emph{ordinal} operations. All arithmetic in $g_{\alpha,n}$, $a_{\alpha,n}$
and subsequent exponent formulas is surreal-field arithmetic. Distinct pairs
$(\alpha,n)$ give distinct ordinal labels: the finite tail after
$\omega\cdot\alpha$ lies strictly before $\omega\cdot(\alpha+1)$.

Each $g_{\alpha,n}$ lies in $G_0$. The family of all these $g_{\alpha,n}$ is
linearly independent over $\Q$, because their inner normal forms are
distinct monomials. Hence the family of all $a_{\alpha,n}$ is algebraically
independent over $\R$: distinct monomials in any finite polynomial produce
distinct sums of the $g_{\alpha,n}$ as outer exponents, and cannot cancel.
\mtag{} For $\alpha=0$ these are the $g_n=g_{0,n}$ and $s_n=a_{0,n}$ of the
proof of \cref{isc:main:join}, with the same independence argument.

\begin{theorem}[Explicit omnific differential freedom]\label{isc:cp:thm:omnific}
Define
\begin{align}
 Y_\alpha
 &=\omega^b\sum_{n\ge1}a_{\alpha,n}\omega^{-ng}\label{isc:cp:eq:Yfactored}\\
 &=\sum_{n\ge1}\omega^{b-ng+g_{\alpha,n}}.
 \label{isc:cp:eq:Ynormal}
\end{align}
Then:
\begin{enumerate}[label=\textup{(\roman*)}]
\item $Y_\alpha$ is a purely infinite omnific integer in $E$ whose support,
read in decreasing order, has order type $\omega$.
\item Its jets are
\begin{equation}\label{isc:cp:eq:Yjets}
 \Delta^jY_\alpha=\sum_{n\ge1}n^j\omega^{b-ng+g_{\alpha,n}}
 \qquad(j\ge0).
\end{equation}
\item The family $(Y_\alpha)_{\alpha\in\On}$ is differentially algebraically
independent over $\Cfg(K_0/\R;G_1)$, hence over $M$, in the finite-relation
sense.
\item For every set cardinal $\kappa$, the subfamily indexed by
$\alpha<\kappa$ has cardinality $\kappa$ and has the same independence
property.
\end{enumerate}
\end{theorem}
\begin{proof}
The inner leading monomial of $b$ has exponent $15/4$, that of $g$ has
exponent $7/2$, and all the $g_{\alpha,n}$ have inner leading exponent below
$1$. Consequently
\[
 b-ng+g_{\alpha,n}>0
\]
for every finite $n$. Moreover
\[
 (b-ng+g_{\alpha,n})-(b-(n+1)g+g_{\alpha,n+1})
   =g+g_{\alpha,n}-g_{\alpha,n+1}>0
\]
by \eqref{isc:eq:separated}. Thus \eqref{isc:cp:eq:Ynormal} is a
reverse-well-ordered normal form with distinct positive exponents indexed in
decreasing order by the positive integers. Its coefficients are all $1$,
and it has no constant term. This proves~(i).

Formula \eqref{isc:cp:eq:Delta} gives eigenvalue $n$ on the $n$th summand,
since $\lambda(b)=0$, $\lambda(g)=1$, and $\Delta$ kills the $G_0$-part.
Strongness gives \eqref{isc:cp:eq:Yjets}.

The coefficient sequences $n\mapsto a_{\alpha,n}$ are disjoint blocks in an
algebraically independent family. Apply
\cref{isc:cp:thm:generic,isc:cp:prop:local} to the unscaled sums in
\eqref{isc:cp:eq:Yfactored}, using $k=\R$, $F=K_0$, and $\Gamma=G_1$. The
multiplier $\omega^b$ is a nonzero member of $M$ with $\Delta\omega^b=0$.
Multiplying every jet by $\omega^b$ is an invertible substitution over the
base, so the scaled family remains independent. This proves~(iii).
Part~(iv) follows by restriction; independence in particular implies
distinctness.
\end{proof}

\mtag{} Since $Y_0=\omega^bS$ with $S$ the witness \eqref{isc:eq:explicitS}
and $\omega^{-b}\in K_1$ is a $\Delta$-constant, part~(iii) contains the
statement that $S$ is differentially transcendental over
$\Cfg(K_0/\R;G_1)$, a strengthening of \cref{isc:thm:joinobstruction} for this
$S$ in the real case. \Cref{isc:thm:joinobstruction} itself needs neither
characteristic-zero derivations nor the explicit coefficients: it holds for
any algebraically independent $s_n$ and also over $\C$.

\begin{corollary}[Countable support is the sharp cardinal threshold]\label[corollary]{isc:cp:cor:threshold}
Every element of $E$ with finite Conway normal-form support belongs to $M$.
Nevertheless, $E\cap\Oz$ contains countably supported purely infinite
elements differentially transcendental over $M$. Thus $\aleph_0$ is the
least possible support cardinality for this phenomenon in the explicit
join.
\end{corollary}
\begin{proof}
A finite-support element is a finite sum of terms
\[
 c\,\omega^{\gamma_0+\gamma_1}=(c\,\omega^{\gamma_0})\omega^{\gamma_1}\in K_0K_1
 \qquad(\gamma_j\in G_j).
\]
Apply \cref{isc:cp:thm:omnific} for the countable examples. The threshold is
about support cardinality, not birthday or cofinality in the value group.
\end{proof}

\begin{remark}[Algebraic enlargements of the base]\label[remark]{isc:cp:rem:algebraicbase}
Every finite family of jets in \cref{isc:cp:thm:omnific} remains
algebraically independent after any algebraic base extension within a
common ambient extension field. Indeed, for a finite tuple $\mathbf y$, if
$N/M$ is algebraic, then $N(\mathbf y)/M(\mathbf y)$ is algebraic;
transcendence degree in the two towers gives
$\trdeg(N(\mathbf y)/N)=\trdeg(M(\mathbf y)/M)$. Only finitely many
algebraic coefficients are involved in a proposed relation. In particular,
the witnesses remain differentially independent over the relative algebraic
closure of $M$ in $E$, with the induced derivation: differentiating a
minimal polynomial shows that this relative algebraic closure is stable
under the derivation. This observation is not a computation of that
closure.
\end{remark}

\section{A fixed tail survives every set of parameters}\label{isc:cp:sec:resilience}

The class-sized strength of the preceding construction is not just that its
indices run through the ordinals. Set-supported normal forms allow every
set of parameters to be localized to a bounded collection of inner
supports. All sufficiently late coefficient labels escape that collection.

\begin{lemma}[Coefficient localization for a parameter set]\label[lemma]{isc:cp:lem:localizeP}
For every set $P\subseteq E$, there is a set $\Theta\subseteq(0,1)$ such
that, if $G_\Theta$ is the additive group of all inner surreal normal forms
supported in $\Theta$ and
\[
 C_\Theta=\K{\R}{G_\Theta}\subseteq K_0,
\]
then $G_\Theta$ and $C_\Theta$ are sets and
\begin{equation}\label{isc:cp:eq:paramcontain}
 P\subseteq C_\Theta((t^{G_1})).
\end{equation}
\end{lemma}
\begin{proof}
Expand each element of $P$ in the iterated coordinates
\eqref{isc:cp:eq:iterated}. Collect all its outer coefficients in $K_0$.
This is a set, since $P$ and all its supports are sets. Expand each of these
coefficients in its Conway normal form, whose exponents lie in $G_0$, and
collect the resulting set of exponents. Expand each such exponent once more
in its inner normal form. Its support is a subset of $(0,1)$. The union
$\Theta$ of all these supports is again a set.

Every collected exponent belongs to $G_\Theta$, so every original outer
coefficient belongs to $C_\Theta$. This gives \eqref{isc:cp:eq:paramcontain}.
The class of inner normal forms supported in a fixed set $\Theta$ injects
into the set of functions $\Theta\to\R$, and is closed under addition and
negatives; it is therefore a set-sized additive group. Likewise the full
Hahn field on the set group $G_\Theta$ injects into the set of functions
$G_\Theta\to\R$ and is a set. All these statements use ordinary replacement
and union, not a class-sized support operation.
\end{proof}

\mtag{} This is the support-closure step of \cref{isc:cor:parameters},
applied to the outer coefficients and then to their exponents: there the
parameters are placed in $\K{\R}{H}$ with $H$ the rational span of their
supports; here their outer coefficients are placed in $\K{\R}{G_\Theta}$.

\begin{lemma}[A tail is independent over the localized coefficient field]\label[lemma]{isc:cp:lem:tailcoeff}
For $\Theta$ and $C_\Theta$ as above, there is an ordinal $\alpha_0$ such
that the family
\[
 \{a_{\alpha,n}:\alpha\ge\alpha_0,\ n\ge1\}
\]
is algebraically independent over $C_\Theta$.
\end{lemma}
\begin{proof}
Because $\psi$ is injective, the ordinals $\beta$ satisfying
$\psi(\beta)\in\Theta$ form a set: each member of $\Theta$ has at most one
inverse image, and replacement collects the ordinal inverse images. Choose
$\alpha_0$ greater than every such ordinal. If $\alpha\ge\alpha_0$, then
$\beta_{\alpha,n}=\omega\cdot\alpha+n\ge\alpha$ and hence
$\psi(\beta_{\alpha,n})\notin\Theta$.

The corresponding $g_{\alpha,n}$ are distinct inner monomials supported
outside $\Theta$. Every nonzero rational linear combination of finitely many
of them still has a nonzero coefficient outside $\Theta$, and thus does not
belong to $G_\Theta$. They are therefore linearly independent over $\Q$
modulo $G_\Theta$.

In a nonzero polynomial in finitely many $a_{\alpha,n}$ with coefficients in
$C_\Theta$, the support of each coefficient is contained in $G_\Theta$.
Multiplication by a monomial in the $a_{\alpha,n}$ translates that support
by the corresponding integer linear combination of the $g_{\alpha,n}$.
Distinct polynomial monomials produce distinct cosets modulo $G_\Theta$.
Their supports are consequently disjoint, so their nonzero terms cannot
cancel. This proves algebraic independence over $C_\Theta$.
\end{proof}

\begin{theorem}[Set-parameter tail resilience]\label{isc:cp:thm:resilience}
For every set $P\subseteq E$, there is an ordinal $\alpha_0$ such that the
predefined tail
\[
 (Y_\alpha)_{\alpha\ge\alpha_0}
\]
is differentially algebraically independent over $M\langle P\rangle$. In
fact it is independent over the larger field
\begin{equation}\label{isc:cp:eq:largebase}
 \Cfg(K_0/C_\Theta;G_1)
\end{equation}
for a set coefficient field $C_\Theta$ supplied by
\cref{isc:cp:lem:localizeP}.
\end{theorem}
\begin{proof}
Choose $\Theta$, $C_\Theta$, and $\alpha_0$ from the preceding lemmas. Apply
the private coefficient-block theorem with base field $C_\Theta$,
coefficient field $K_0$, outer group $G_1$, and the tail coefficients
$a_{\alpha,n}$. Its hypotheses hold by \cref{isc:cp:lem:tailcoeff}. The
Euler weight map takes its values in $\R\subseteq C_\Theta$, so it is valid
over this enlarged base. The set-local formulation \cref{isc:cp:prop:local}
applies because $C_\Theta$ is a set. This gives differential independence of
the unscaled tail over \eqref{isc:cp:eq:largebase}, and then of the
$Y_\alpha$ after multiplication by $\omega^b$.

It remains to verify inclusion of the desired parameter field. The envelope
\eqref{isc:cp:eq:largebase} contains every element of $K_0$ as an outer
constant, since an individual element generates a finitely generated
extension of $C_\Theta$. It contains $K_1$ because its coefficients lie in
$\R\subseteq C_\Theta$. It contains $P$ by \eqref{isc:cp:eq:paramcontain}.
Finally it is stable under $\Delta$, by the proof of \cref{isc:cp:lem:cfg}.
Therefore it contains $M\langle P\rangle$, and independence descends to that
smaller base.
\end{proof}

\begin{corollary}[No set-sized algebraic or differential generation]\label[corollary]{isc:cp:cor:nosetgen}
For no set $P\subseteq E$ is $E$ algebraic over $M(P)$. For no such $P$ is
$E$ differential-algebraic over $M\langle P\rangle$. In particular, no set of
ordinary or differential generators, even followed by the respective
algebraic closure operation within $E$, generates $E$ over $M$.
\end{corollary}
\begin{proof}
A member of the tail supplied by \cref{isc:cp:thm:resilience} is
differentially transcendental over $M\langle P\rangle$ and therefore
algebraically transcendental over its subfield $M(P)$. This contradicts
either asserted algebraicity.
\end{proof}

\begin{remark}[Why the set qualification is sharp for the method]\label[remark]{isc:cp:rem:setsharp}
A proper class of parameters could use every inner support label
$\psi(\beta)$. The localization set $\Theta$ would then no longer be a set,
and the ordinal bound in \cref{isc:cp:lem:tailcoeff} would not be available.
The theorem is a genuine set-parameter result, not a statement about
arbitrary proper-class parameter collections. It also does not assert
independence over a differential-algebraic closure of $M\langle P\rangle$;
that stronger base-change statement is unnecessary for
\cref{isc:cp:cor:nosetgen}.
\end{remark}

\section{Arithmetic joins and the surcomplex extension}\label{isc:cp:sec:arithmetic}

\subsection{Fraction fields of the three integer rings}

Let $R_j=K_j\cap\Oz$ as in \eqref{isc:cp:eq:Ri}. Let
\begin{equation}\label{isc:cp:eq:arithrings}
 R_0[R_1]=\text{the subring generated by }R_0\cup R_1,
 \qquad O=E\cap\Oz.
\end{equation}
Thus $R_0[R_1]$ is formed by finite ring operations, whereas $O$ permits
every set-supported omnific normal form whose exponents lie in $G_0+G_1$.
Both are class domains. Fractions and integrality are interpreted by their
ordinary finite formulas, as in \cref{isc:sec:fractions}.

\begin{lemma}[Monomial denominator clearing]\label[lemma]{isc:cp:lem:clearing}
One has
\begin{equation}\label{isc:cp:eq:fractions}
 \Frac(R_j)=K_j,\qquad \Frac(R_0[R_1])=M,\qquad \Frac(O)=E.
\end{equation}
\end{lemma}
\begin{proof}
For $x\in K_j$, its support is a set in $G_j$. Pull it back through the
ordered-group isomorphism $h_j:\No\to G_j$. Every set of surreal numbers has
an upper bound, so choose $d\in G_j$, $d>0$, with $d+e>0$ for every $e$ in
the support of $x$. Both $\omega^d$ and $\omega^d x$ have only positive
exponents, and hence belong to $R_j$. Thus $x=(\omega^d x)/\omega^d$ is a
fraction of elements of $R_j$. The reverse inclusion is immediate.

The fraction field of $R_0[R_1]$ contains both $K_j$ and is contained in
their compositum, giving $\Frac(R_0[R_1])=M$.

Finally take $x\in E$. Project its set of support exponents to a set in
$G_1$. Choose $d\in G_1$, $d>0$, strictly greater than the negatives of all
these projections. If $\gamma_0+\gamma_1$ occurs in $x$, then $d+\gamma_1$
is positive in $G_1$ and dominates $\gamma_0$ by \eqref{isc:eq:separated}.
Hence $d+\gamma_0+\gamma_1>0$. The same denominator $\omega^d$ clears $x$
into $O$, proving $\Frac(O)=E$.
\end{proof}

The first clearing mechanism is already part of this report's omnific
fraction theory (\cref{isc:lem:denominator}, and $\Frac(R_\alpha)=K_\alpha$
in \cref{isc:prop:uniform}). It is included here because the exact fraction
field of the \emph{two-copy generated ring} is needed in the next result.

\begin{theorem}[Arithmetic failure cannot be repaired by a set]\label{isc:cp:thm:arithmetic}
With $R_0[R_1]\subseteq O$ as in \eqref{isc:cp:eq:arithrings}:
\begin{enumerate}[label=\textup{(\roman*)}]
\item For every set $J\subseteq\On$, substitution $X_\alpha\mapsto Y_\alpha$
embeds the polynomial $R_0[R_1]$-algebra in the variables
$(X_\alpha)_{\alpha\in J}$ into $O$.
\item For every set $S\subseteq O$, the extension $O/R_0[R_1][S]$ is not
integral. In fact $E$ is not algebraic over $\Frac(R_0[R_1][S])=M(S)$.
\item The obstructing element in~(ii) can be chosen from a tail of the same
family \eqref{isc:cp:eq:Ynormal}; it has purely positive, countable Conway
support.
\end{enumerate}
\end{theorem}
\begin{proof}
Since $R_0[R_1]\subseteq M$, algebraic independence over $M$ implies the
injectivity in~(i); the image lies in $O$ by \cref{isc:cp:thm:omnific}. For
(ii), apply \cref{isc:cp:thm:resilience} with $P=S$. Any tail member is
transcendental over $M(S)$, while \cref{isc:cp:lem:clearing} gives
$\Frac(R_0[R_1][S])=M(S)$. If it were integral over $R_0[R_1][S]$, a monic
polynomial over that ring would make it algebraic over $M(S)$, a
contradiction. The same element proves~(iii).
\end{proof}

\begin{remark}[Not a normalization calculation]\label[remark]{isc:cp:rem:normalization}
Nonintegrality of $O/R_0[R_1][S]$ does not compute the integral closure of
$R_0[R_1]$ in $M$ or in $E$. Nor does it compute a Picard group, a class
group, or a projective module. The theorem isolates a much larger
transcendental obstruction that no set of arithmetic parameters removes.
Questions about normalization inside the fixed fraction field remain
distinct.
\end{remark}

\subsection{Surcomplex and Gaussian omnific versions}

Put
\[
 E_{\C}=E[i],\qquad K_{j,\C}=K_j[i],\qquad M_{\C}=M[i].
\]
Because $E$ is ordered, adjoining $i$ gives a genuine quadratic field
extension. Separating real and imaginary coefficients identifies these
fields with the corresponding full complex-coefficient Hahn fields on the
same groups. The union of two set supports is still a set, so this
identification preserves the Hahn conventions.

Extend $\Delta$ by $\Delta i=0$. Define
\[
 O_{\C}=E_{\C}\cap\Oz[i]=O[i].
\]
The Gaussian integer images are $R_j[i]$, and the ring they generate is
$R_0[R_1][i]$.

\begin{theorem}[Surcomplex independence and Gaussian witnesses]\label{isc:cp:thm:complex}
The following statements hold.
\begin{enumerate}[label=\textup{(\roman*)}]
\item The real family $(Y_\alpha)$ remains differentially algebraically
independent over $M_{\C}$ in $E_{\C}$.
\item The genuinely nonreal elements
\begin{equation}\label{isc:cp:eq:Zfamily}
 Z_\alpha=(1+i)Y_\alpha
\end{equation}
belong to $\Oz[i]$, have countable support, and are differentially
algebraically independent over $M_{\C}$.
\item For every set $P\subseteq E_{\C}$, a tail of each of these families is
differentially algebraically independent over $M_{\C}\langle P\rangle$.
\item $\Frac(R_0[R_1][i])=M_{\C}$ and $\Frac(O_{\C})=E_{\C}$. For every set
$S\subseteq O_{\C}$, the ring $O_{\C}$ is not integral over
$R_0[R_1][i][S]$.
\end{enumerate}
\end{theorem}
\begin{proof}
For a polynomial relation over $M[i]$ among real jets, write each
coefficient uniquely as $a+ib$, $a,b\in M$. Its real and imaginary parts
give polynomial relations over $M$. At least one is a nonzero polynomial if
the original was nonzero. This contradicts \cref{isc:cp:thm:omnific},
proving~(i). Multiplication by the nonzero constant $1+i$ is an invertible
substitution on all jets, proving~(ii).

For~(iii), take the set $P_{\R}\subseteq E$ consisting of the real and
imaginary parts of all members of $P$. Apply \cref{isc:cp:thm:resilience} to
$P_{\R}$. Splitting coefficients as above gives independence over
$M\langle P_{\R}\rangle[i]$, which contains $M_{\C}\langle P\rangle$. Again
rescale to obtain the $Z_\alpha$ version.

For~(iv), the displayed fraction fields follow from
\cref{isc:cp:lem:clearing} after adjoining $i$, or by clearing real and
imaginary parts with a common monomial. A tail element supplied by~(iii) is
transcendental over $\Frac(R_0[R_1][i][S])=M_{\C}(S)$, so cannot be integral
over $R_0[R_1][i][S]$.
\end{proof}

No statement here depends on extending analytic functions to all
surcomplex arguments, choosing a complex logarithm, or assuming an
exponential-compatible derivation. The passage from real to complex is
finite algebraic extension and coefficient separation.

\section{What the addition resolves, and what it does not}\label{isc:cp:sec:audit}

\subsection{Relation to Research question 3}

Research question~3 has two parts: an exact description of the ordinary
compositum and a dimension calculation for the Hahn join over it. The
results of the addition give the following answers in specified domains.

\begin{center}\small
\begin{tabular}{@{}p{0.29\textwidth}p{0.63\textwidth}@{}}
\toprule
Problem component & Conclusion proved by source~02 \\
\midrule
Membership in separated coordinates & The finite coefficient-rank
denominator criterion, \cref{isc:cp:thm:rank,isc:cp:cor:surrealrank}. \\
Finite coefficient-field condition & Not sufficient even for algebraicity;
a witness uses only $k(x)$, \cref{isc:cp:cor:fginsufficient}. \\
Natural set-sized dimensions & Exact algebraic and differential dimensions
for $F((t))/F k((t))$ under the hypotheses of \cref{isc:cp:thm:cardinal}. \\
Explicit class-sized join & Independent families of every set cardinal,
with countable purely infinite supports, \cref{isc:cp:thm:omnific}. \\
Dependence on set parameters & A fixed tail remains independent after any
set adjunction, \cref{isc:cp:thm:resilience}. \\
Arithmetic and complex forms & Nonintegrality persists after set
adjunction; the same phenomena occur in $\Oz[i]$,
\cref{isc:cp:thm:arithmetic,isc:cp:thm:complex}. \\
\bottomrule
\end{tabular}
\end{center}

These answers do not solve the complete problem for arbitrary non-separated
ordered amalgams of exponent groups, arbitrary coefficient-field
cardinalities, and arbitrary prescribed common Hahn cores. The rank
criterion applies whenever an iterated coefficient interpretation has been
established. The cardinal theorem does not assign a set cardinal to the
class extension. The question addressed is a question posed in this report,
not a claimed resolution of a separately established, long-standing
published conjecture.

\subsection{Provenance and proof dependencies}

The source compared itself with this report at commit
\nolinkurl{7b256e0792df7718995b816ad12a77bf56731f8f}, where the report's
source had blob identifier
\nolinkurl{7d5cfe644e7c902f2b594179c337a6c633350fc3}. The key portions
consulted were the explicit-copy and re-expansion material of
\cref{isc:sec:explicit,isc:sec:join} and the question list of
\cref{isc:sec:questions}. The report's README and the collection's main
catalogue were also consulted. The source's audit (shipped as
\rpath{02-composita-SOURCE\_AUDIT.md}) records the paths and distinguishes
sources read for substance from sources consulted only for orientation.

The explicit maps \eqref{isc:eq:inner} and \eqref{isc:eq:explicitouter}, the
scale separation, and the finite coefficient-field observation are credited
to this report. The coset-slice argument belongs to the same elementary Hahn
toolkit and is not presented as a priority claim. The proofs of the new
class conclusions do not import this report's general transfinite
independent-copy recursion (\cref{isc:thm:group}), nor, in the source's
words, the repository's proposed automatic-strongness results,
curve-rigidity results, or unreviewed quotient theorems.

The main dependency chains are
\begin{align*}
 \text{rapid degrees}
 &\Longrightarrow \text{continuum family}
 \Longrightarrow \text{small-cardinality case},\\
 \text{coefficient derivations}
 &\Longrightarrow \text{private blocks}
 \Longrightarrow \text{tree amplification},\\
 \text{tree amplification}
 &\Longrightarrow \text{large-cardinality case},\\
 \text{separated copies and private blocks}
 &\Longrightarrow \text{omnific witnesses},\\
 \text{inner-support localization}
 &\Longrightarrow \text{tail resilience},\\
 \text{tail resilience}
 &\Longrightarrow \text{arithmetic and complex consequences}.
\end{align*}
No numerical experiment replaces any arrow in these chains.

\subsection{Critical proof checks}

The rapid-degree argument depends on more than comparing leading terms
informally. The polynomial relation is made integral in the outer variable
before valuations are compared. Minimal total jet degree guarantees a
nonzero evaluated partial derivative. The selected coefficient lies at
$n+\ell<2n$, so all quadratic tail terms vanish there. Partial coefficients
below $\ell$ vanish, so later tail coefficients cannot interfere. Finally the
coefficient degree $\nu$ is fixed even though the outer coefficient series
are arbitrary Laurent series.

The private-block argument depends on a different finiteness principle. A
finitely generated field $D$ is algebraic over a field generated by finitely
many elements of a transcendence basis. Coordinate derivations outside that
finite set kill $D$ after separable extension. This is why the argument
controls whole coefficient series in $D((t^\Gamma))$, not merely their
leading coefficients. Characteristic zero is used both for separability and
for distinct integer eigenvalues.

The class argument uses two levels of support localization. An outer
coefficient is a surreal series in exponents from $G_0$; those exponents
themselves have inner supports in $(0,1)$. Taking the union of the inner
supports for a \emph{set} of parameters produces a set $\Theta$. The
predefined ordinal coefficient labels outside $\Theta$ remain independent
modulo $G_\Theta$. Replacing this argument with an unproved claim that ``a
class is larger than a set'' would not establish the required independence
over the parameter field.

The source's proof audit (shipped as \rpath{02-composita-PROOF\_AUDIT.md})
adds three points. The auxiliary coordinate derivations of
\cref{isc:cp:thm:generic} need not preserve any \emph{inner} surreal
summation; they are extended coefficientwise only in the explicitly used
\emph{outer} Hahn representation. The rapid-degree argument does not require
sparse outer support, but it does require a uniform finite degree in $x$
after denominators are cleared. In \cref{isc:cp:thm:omnific} the
$g_{\alpha,n}$ are inner monomials in $(0,1)$ and the $a_{\alpha,n}$ are
independent outer coefficient monomials; the Euler derivation has weight
$-\lambda(\gamma_1)$ in Conway notation, so the $n$th witness term has
eigenvalue $n$, and the shift multiplier has derivative zero.

\subsection{Limitations and nonclaims, as stated by the source}

The proofs are mathematical arguments, not Lean-checked results. The finite
program supplied with the source tests identities and finite models only.
Independent review should prioritize
\cref{isc:cp:thm:lacunary,isc:cp:thm:generic,isc:cp:thm:resilience}. The
source search did not certify novelty or exhaustive priority. The adjective
``strong'' refers to Hahn summability, not convergence in an order topology.

No claim is made about Gonshor-exponential compatibility, preservation of
simplicity, initiality of the two images, birthday optimality, or the
Berarducci--Mantova derivation. The full algebraic closure of the
compositum, its normalization, and its projective modules are not
calculated. No general algorithm deciding compositum membership is proved.
These distinctions are mathematically relevant, not merely editorial
qualifications.

\subsection{What the addition does not claim}\label{isc:cp:sub:nonclaims}

\mtag{} This list collects every limitation stated by source~02, in its
article, delivered README, source and proof audits, recorded run and build
audit, with its location here, followed by those added in the merge.
\begin{enumerate}[label=C\arabic*.,leftmargin=2.6em,itemsep=2pt]
\item The rank test is not a decision procedure: the existential search
over $q$ is an infinite-input problem, and no algorithm deciding arbitrary
Hahn-series membership follows (\cref{isc:cp:rem:ranktest};
Research question~\ref{isc:cp:q:effective}).
\item The tensor fact is not new: it is a standard linear-disjointness
mechanism. No new tensor-product theorem and no general theory of rational
series is claimed (\cref{isc:cp:rem:ranktest}; source audit).
\item The comparison with Krapp--Kuhlmann--Serra \cite{KKS} is conceptual,
not a claim that the formulations have identical inputs; that paper was
used for the literature comparison, not as a proof of the compositum
theorems (\cref{isc:cp:rem:ranktest}; bibliography).
\item Infinite coefficient rank obstructs membership in $F[L]$, not in $M$;
a geometric series of infinite rank lies in $M$
(\cref{isc:cp:ex:geometric,isc:cp:cor:finitecoeff}).
\item Base extension is not treated as harmless; coefficient extension uses
coefficientwise independence of constants, intermediate base change clears
denominators, and exponent-group extension uses coset slices, not an
unjustified claim that arbitrary base extension preserves algebraic
independence. No convexity is required for the coset slice, but the
surreal iterated presentation requires the explicit scale separation
(\cref{isc:cp:sec:disjoint}; proof audit).
\item The rapid-degree theorem does not rest on gaps in the outer support;
merely pointing to such gaps would not prove \cref{isc:cp:ex:mechanism}. It
requires characteristic zero, a transcendental coefficient variable,
unbounded ratios $d_n/d_{n-1}$, infinitely many private indices relative to
each finite subfamily, and a uniform finite $x$-degree after clearing
denominators (\cref{isc:cp:thm:lacunary,isc:cp:rem:laurentcoeff}; proof
audit).
\item The formula $|F|^{\aleph_0}$ is not extended to infinite algebraic
coefficient extensions, to $\Gamma=0$, to groups without a specified
$\lambda$ with $\lambda(g)=1$ (only the algebraic rank test applies there),
or to $|k|>\cc$ with small positive transcendence degree, where only the
explicit lower bounds are claimed; the rapid-degree family need not reach
the cardinality of the Laurent field for large $k$
(\cref{isc:cp:sub:boundary}; Research questions~\ref{isc:cp:q:large}
and~\ref{isc:cp:q:algebraic}).
\item Characteristic zero is essential for the differential statements: in
characteristic $p$ the Euler derivation satisfies $\partial^p=\partial$. This
does not settle the ordinary algebraic-transcendence problem in
characteristic $p$ (\cref{isc:cp:sub:boundary}).
\item The derivation $\Delta$ is an explicitly specified diagonal Euler
derivation. It is not asserted to be Hardy type, the Berarducci--Mantova
derivation, or compatible with exponentiation. It is proved to be a
derivation directly and is not identified with a Hardy-type derivation of
Kuhlmann--Matusinski \cite{KM}, whose Hardy-type theory the source did not
review in full (\cref{isc:cp:sub:euler,isc:cp:sub:scope}; source audit;
Research question~\ref{isc:cp:q:bm}).
\item The explicit copies do not commute with the omega-map, and strong
Hahn summability alone does not provide compatibility with it or with the
exponential (\cref{isc:cp:sub:copies}; Research question~\ref{isc:cp:q:exp}).
\item \Cref{isc:cp:rem:algebraicbase} is not a computation of the relative
algebraic closure of $M$ in $E$.
\item The class results use only finite relations, set-indexed families and
set-parameter resilience. No class transcendence basis is chosen, no set
cardinal is assigned to the proper-class transcendence degree, and no
cardinal equal to the size of $\No$ is postulated. No class-state recursion
or class satisfaction predicate is used (\cref{isc:cp:sec:foundations};
\cref{isc:cp:prop:local}; delivered README; proof audit).
\item Tail resilience is for \emph{set} parameters only; proper-class
parameter collections are not covered. Independence is not asserted over
the differential-algebraic closure of $M\langle P\rangle$
(\cref{isc:cp:rem:setsharp}).
\item Nonintegrality of $O/R_0[R_1][S]$ does not compute the integral
closure of $R_0[R_1]$ in $M$ or $E$, a normalization, a Picard group, a
class group or a projective module (\cref{isc:cp:rem:normalization};
Research question~\ref{isc:cp:q:normalization}).
\item The complex versions use only finite algebraic extension and
coefficient separation: no analytic continuation to surcomplex arguments,
no complex logarithm, no exponential-compatible derivation
(\cref{isc:cp:sec:arithmetic}).
\item Research question~3 is not solved for non-separated ordered amalgams,
arbitrary coefficient-field cardinalities, or arbitrary prescribed common
Hahn cores; the cardinal theorem assigns no set cardinal to the class
extension; the question is this report's own, not a long-standing
published conjecture (\cref{isc:cp:sec:audit}).
\item The explicit maps, scale separation and finite coefficient-field
observation are this report's; the coset slice is elementary. The proofs do
not import this report's transfinite recursion, nor automatic-strongness,
curve-rigidity or quotient theorems of the repository. No numerical
experiment replaces a proof step (\cref{isc:cp:sec:audit}).
\item Nothing is claimed on Gonshor-exponential compatibility, simplicity,
initiality, birthday optimality, the full algebraic closure of the
compositum, normalization or projective modules; no general membership
algorithm is claimed (\cref{isc:cp:sec:audit}; Research
question~\ref{isc:cp:q:birthday}).
\item Status: the proofs are not Lean-checked, not refereed and not
independently reviewed; the source's proof audit is author-side, not a
referee report; review should prioritize
\cref{isc:cp:thm:lacunary,isc:cp:thm:generic,isc:cp:thm:resilience}.
Priority is not certified: the literature comparison was targeted, a
search that finds nothing does not establish novelty, and not all related
literature was read (\cref{isc:cp:sub:scope,isc:cp:sec:audit}; title page of
the source; source and proof audits; build audit).
\item The source's comparison with this report was targeted: it read two
line ranges of the pinned source for substance, not every proof of the
report. The Stacks Project was used for terminology; the encyclopedia page
on surreal numbers for orientation only; L'Innocente--Mantova for context,
with no factorization result claimed (source audit; bibliography).
\item The research questions are proposed by the proved statements. Except
for Research question~3 of this report, they are not represented as
previously published open problems (\cref{isc:cp:sec:questions}).
\item The finite checks are narrower than the theorems. They test
polynomial identities, finite coefficient ranks, finite combinatorial
examples and finite ordered-group models; they do not implement $\No$ or
infinite Hahn summation, do not represent countably supported surreals, and
prove no independence, cardinal or proper-class assertion. The
lexicographic model is a model of the comparisons, not an approximation to
a surreal number. Neither typesetting nor finite computation is proof
verification (\cref{isc:cp:sec:verify}; recorded run).
\item The formalization decomposition names implementation targets, not
declarations that exist (\cref{isc:cp:sub:formal}).
\end{enumerate}
\begin{enumerate}[label=M\arabic*.,leftmargin=2.6em,itemsep=2pt]
\item Parts~(i)--(ii) of \cref{isc:cp:thm:rank} are
\lbl{duals:prop:algebraic} of the three-duals report and its preceding
injectivity paragraph; part~(iii) is new but immediate. The source does not
cite that report. \Cref{isc:cp:lem:coset} is a variant of
\cref{isc:thm:sliceddisjoint} (\cref{isc:cp:sec:collection}).
\item The private-index and tree mechanisms are related to
\lbl{tail:thm:surreal} and \lbl{bst:thm:independent}, which use different
bases, derivations and tools, and the tail resilience to
\lbl{adr:sr:thm:tail}; the source cites none of them
(\cref{isc:cp:sec:collection}).
\item Research question~3 is answered only in the separated case with
trivial core and in the cardinal cases of \cref{isc:cp:thm:cardinal}; its
remaining part is Research questions~\ref{isc:cp:q:amalgams}
and~\ref{isc:cp:q:cores}. \mtag{} Source~03 later answers the degree part
for explicit non-separated examples (\cref{isc:hj:sec:answers}). No other question of this report is answered, and
no question of another report (\cref{isc:cp:sec:answers}).
\item No \texttt{isc:cp:} label has a Lean mapping in
\rpath{FORMALIZATION.md}. The proofs were read when the addition was
merged (Appendix~\ref{isc:cp:app:source}); that reading is not an
independent review.
\item The recorded run was repeated in the merge on a copy, with identical
counts; the shipped \rpath{code/02-composita-build.sh} does not run
unchanged under the shipped file names (Appendix~\ref{isc:cp:app:files}).
\end{enumerate}

\section{Research questions of the addition}\label{isc:cp:sec:questions}

The following questions are proposed by the proved statements of
source~02. Except for Research question~3, which the source takes from this
report, they are not represented as previously published open problems.
Each asks for a specific extension, obstruction, or effective refinement.
\mtag{} The source's Question~$n$ is Research question~$10+n$ here.

\begin{question}[Non-separated ordered amalgams]\label{isc:cp:q:amalgams}
Let $G_0,G_1$ be independent ordered subgroups but assume neither is
infinitesimal relative to the other. Determine the image of
$\K{k}{G_0}\K{k}{G_1}$ inside $\K{k}{G_0+G_1}$ and its transcendence defect.
Can a finite-rank denominator criterion be formulated without an iterated
Hahn coefficient field?
\end{question}
The obstacle is not the finite tensor argument itself, but choosing slices
with support properties strong enough to characterize the entire join. This
is the main part of Research question~3 left untreated.
\mtag{} \emph{Status (batch 35): the transcendence defect is known for
explicit examples; the image and a denominator criterion are open.} For the
interleaved groups of \cref{isc:hj:main:cardinal} the defect is $2^\kappa$,
for the dense Archimedean groups of \cref{isc:hj:main:prime} it is $\cc$,
and for the class groups of \cref{isc:hj:main:class} there is no set-sized
transcendence basis (\cref{isc:hj:sec:answers}). These groups are explicit
subgroups, not the exponent groups of copies of $\No$ (the question allows
any independent groups); the criterion \cref{isc:hj:main:mixed} applies to
any pair of full Hahn fields on groups meeting its hypotheses. The smallest open dense case is Research
question~\ref{isc:hj:q:rankone}; the image is Research
question~\ref{isc:hj:q:characterization}.

\begin{question}[Very large base fields]\label{isc:cp:q:large}
For $|k|>\cc$ and $0<\trdeg(F/k)<|k|$, must
\[
 \trdeg\bigl(F((t))/F k((t))\bigr)=|F|^{\aleph_0}
\]
still hold? Does the same equality hold for differential transcendence
degree under the outer Euler derivation?
\end{question}
The rapid-degree theorem gives a continuum family over any $k$, but that
family need not reach the cardinality of the entire Laurent field. A
stronger combinatorial coefficient construction or a counterexample would
settle the gap in the cardinal argument.

\begin{question}[Infinite algebraic coefficient extensions]\label{isc:cp:q:algebraic}
For an infinite algebraic extension $F/k$ of characteristic zero,
characterize $F k((t))$ inside $F((t))$ and determine the algebraic and
differential transcendence degrees of the latter over the former.
\end{question}
Finite extensions give equality of fields, while the
transcendental-coordinate derivations no longer apply. Growth in the
algebraic degrees needed to contain successive coefficients is a natural
replacement invariant, but no theorem about it is claimed.

\begin{question}[Prescribed nonzero common Hahn cores]\label{isc:cp:q:cores}
Can the explicit omnific tail-resilience theorem be extended to independent
surreal copies with a prescribed set-sized common field $\K{\R}{H}$? Which
separation conditions on the ordered amalgam over $H$ suffice?
\end{question}
This report already constructs copies with prescribed cores
(\cref{isc:main:copies}). What is needed is a compatible local coefficient
decomposition and an independent-label family modulo the shared support
group, not just an abstract pair of embeddings. \mtag{} Here $H$ is the
common core of \cref{isc:main:copies}, as in the source's question.
\mtag{} \emph{Status (batch 35): continued, not answered.} The
mixed-support criterion \cref{isc:hj:main:mixed} allows a nonzero core
algebraically, but no example with $H\ne0$ is constructed and no
resilience is proved; the ordered construction is Research
question~\ref{isc:hj:q:cores}, and full copies over a core are Research
question~\ref{isc:hj:q:copies}.

\begin{question}[Presentation-independent coefficient rank]\label{isc:cp:q:presentation}
Which part of the finite-rank denominator criterion can be recovered from
the pair of embedded fields alone, without naming a monomial cross-section
or a chosen iterated Hahn presentation?
\end{question}
The compositum is intrinsic, but coefficient rank is a
presentation-dependent certificate. An invariant description of the
smallest possible certificate spaces could distinguish coordinate artifacts
from structure of the pair.

\begin{question}[Effective membership on controlled inputs]\label{isc:cp:q:effective}
For rational, recurrence-defined, automatic, or otherwise effectively
presented supports and coefficients, when is the finite-rank denominator
test decidable? Can one bound the denominator rank in a useful subclass?
\end{question}
The generalized recurrence framework of \cite{KKS} is a relevant starting
point. A useful algorithm must specify its representation of infinite
series and should not treat a formal coefficient oracle as an effective
finite input.

\begin{question}[Natural surreal derivations]\label{isc:cp:q:bm}
Do comparable compositum-independence and set-parameter resilience
statements hold for a Hardy-type surreal derivation, particularly the
Berarducci--Mantova derivation, on subfields stable under it?
\end{question}
The proofs exploit a diagonal derivation that kills the entire lower-scale
coefficient field. A natural surreal derivation generally differentiates
those coefficients and may not preserve the explicit copy images.
Establishing an appropriate stable pair is part of the problem, not an
automatic corollary of \cref{isc:cp:thm:omnific}. \mtag{} Research
question~6 asks the analogous question for the prescribed-intersection
theorem.

\begin{question}[Exponential-compatible copies]\label{isc:cp:q:exp}
Can one obtain a similarly large gap for a pair of proper surreal
self-embeddings commuting with the Gonshor exponential, or with the
omega-map? Which common cores must then be closed under the additional
operation?
\end{question}
The displayed copies fail to commute with the omega-map. Strong Hahn
summability alone does not provide the desired compatibility. This question
connects the compositum programme to the richer embedding questions already
distinguished in this report (Research question~6).

\begin{question}[Normalization within the fraction field]\label{isc:cp:q:normalization}
For $R_0[R_1]$ from \eqref{isc:cp:eq:arithrings}, compute or constrain its
integral closure inside $M$. How does that ring compare with $M\cap\Oz$, and
what can be said about finitely generated projective modules over it?
\end{question}
The nonintegrality of $O/R_0[R_1][S]$ comes from transcendence outside $M$
and therefore does not answer this within-fraction-field problem. A
normalization result would complement rather than restate
\cref{isc:cp:thm:arithmetic}. \mtag{} Research question~4 asks the
corresponding questions for the intersection ring $\A{k}{H}$.

\begin{question}[Birthday complexity of minimal-support witnesses]\label{isc:cp:q:birthday}
Among countably supported omnific elements transcendental over the explicit
compositum, what birthday bounds are possible? Can one lower the nested
normal-form complexity while retaining set-parameter tail resilience?
\end{question}
The support cardinality $\aleph_0$ is already sharp by
\cref{isc:cp:cor:threshold}. Birthday optimality is a different invariant
and requires genuine estimates on surreal construction lengths.

\begin{question}[Several scales and commuting derivations]\label{isc:cp:q:scales}
For three or more separated copies, construct several commuting Euler
derivations and determine algebraic independence of mixed jets over the
corresponding ordinary compositum. What is the correct tree-label condition
for joint differential independence?
\end{question}
Coefficient derivations still offer finite avoidance, but one must separate
multivariable eigenvalue interpolation from the ordered-group support
problem. A multidimensional Vandermonde argument alone does not settle the
Hahn support conditions. \mtag{} The algebraic counterpart for $r$ factor
fields, without separation, is Research question~\ref{isc:hj:q:factors}.

\begin{question}[Formalization of the core mechanisms]\label{isc:cp:q:formal}
Can the coefficient-rank theorem, the coefficient-extension and coset-slice
lemmas, the rapid-degree Taylor argument, and the private-block theorem be
formalized as reusable Hahn-field results before specializing to surreal
numbers?
\end{question}
A promising order is to formalize the finite polynomial and
coefficient-space identities first, then the set-sized Hahn statements, and
only afterwards the class-local interface. This would test the most
important new proofs without requiring a monolithic formalization of the
entire report. \mtag{} It overlaps Research question~10 in the coset-slice
kernel, and Research question~\ref{isc:hj:q:formal} in the coset and
Euler-derivation kernel.

\section{Finite verification and a formalization decomposition of the addition}
\label{isc:cp:sec:verify}

\subsection{Finite verification and reproducibility}

The source's program, shipped as \rpath{code/02-composita-verify.py},
performs exact finite checks. Its recorded output, shipped as
\rpath{data/02-composita-verification.json}, records the actual checks and
their counts: 2,530 checks in 21 categories. These checks are deliberately
narrower than the theorems. They do not represent countably supported
surreal numbers and do not prove an independence or cardinality statement.

The checked mechanisms include finite geometric-series convolution and
coefficient ranks, Euler and coefficient-derivation identities on
polynomial examples, Vandermonde determinants at distinct integer
eigenvalues, finite prefix-tree private-label behaviour, factorial growth
inequalities used in \eqref{isc:cp:eq:nchoices}, and a multivariable Taylor
coefficient calculation with quadratic tails excluded at the relevant
order. A finite lexicographic exponent model tests the sign and strict
descent used in \eqref{isc:cp:eq:Ynormal}. It is a model of the comparisons,
not a numerical approximation to an actual surreal number.

A characteristic-$p$ check illustrates the explicit boundary
$\partial^p=\partial$ on monomials. Finite complex substitution checks
confirm that multiplying variables by $1+i$ is invertible; the proof for
arbitrary polynomials remains the elementary field argument above.

The program uses the standard library and SymPy. It writes a file only when
given \texttt{--output}. The source's build script also supports repeated
\texttt{pdflatex} passes when \texttt{latexmk} is unavailable, and its build
audit records the compiler result, page count, reference checks, and layout
checks for the delivered PDF. Neither successful typesetting nor finite
computation constitutes proof verification. \mtag{} The shipped file names,
the rerun and its comparison with the recorded output are in
Appendix~\ref{isc:cp:app:files}.

\subsection{A proposed formalization decomposition}\label{isc:cp:sub:formal}

The following are implementation targets, not declarations claimed to exist
in the collection.

\paragraph{Coefficient-space layer.}
Represent a finite-dimensional coefficient span and prove that its
coordinate functions define Hahn series with supports contained in the
original support. Prove tensor injectivity by coefficient comparison.
Derive the rank-denominator criterion using the ordinary fraction field
universal property.

\paragraph{Support and base-change layer.}
Define restriction to an exponent coset, prove its compatibility with
multiplication by series on the subgroup, and establish the two
linear-disjointness transfer lemmas. Make the chosen ambient field explicit
to avoid conflating abstract tensor products with a specified compositum.

\paragraph{Polynomial perturbation layer.}
For a polynomial in finitely many jet variables over $k[[t]][x]$, prove a
first-order Taylor congruence modulo $t^{2n}$ after a perturbation divisible
by $t^n$. Prove the coefficient-degree bound $\nu+d\,d_{n-1}$. These finite
statements are the delicate reusable kernel of \cref{isc:cp:thm:lacunary}.

\paragraph{Derivation and interpolation layer.}
Formalize coefficientwise Hahn derivations, their commutation with Euler
derivations, extension of a derivation over separable algebraic fields,
and the finite Vandermonde elimination of evaluated partials. This yields
the set-sized private-block theorem.

\paragraph{Cardinal and surreal layer.}
Add the prefix-tree construction and cardinal counting. For the explicit
surreals, prove the two-level normal-form relabeling, lexicographic
comparison, and set-support localization. Each proposed relation must be
transported into a set-sized field before invoking a transcendence basis.
No unrestricted class basis should be introduced just to shorten the proof.

\mtag{} The coset-slice and iterated-Hahn items overlap items (1)--(3)
and~(6) of \cref{isc:sub:formal}.

\section{Maximal transcendence in Hahn joins: a third manuscript}
\label{isc:hj:sec:intro}

\mtag{} \Cref{isc:hj:sec:intro,isc:hj:sec:foundations,isc:hj:sec:lemmas,isc:hj:sec:mixed,isc:hj:sec:combinatorics,isc:hj:sec:cardinal,isc:hj:sec:archimedean,isc:hj:sec:prime,isc:hj:sec:classes,isc:hj:sec:strong,isc:hj:sec:scope,isc:hj:sec:questions,isc:hj:sec:conclusion,isc:hj:sec:verify}
print a third manuscript, written against this report. In the report's
local numbering (Appendix~\ref{isc:hj:app:source}) it is \emph{source~03}:
manuscript~01 of batch~35 of the collection, archive
\texttt{Surreal\_Hahn\_Joins\_Research}, \emph{Maximal Transcendence in Hahn
Joins: Mixed-support independence and prime omnific integers}, a 27-page
PDF dated 23 September 2026 and pinned to repository commit
\nolinkurl{f388f950f44664f21f50f04b8ac74dab56c8ff2f}. At that pin this
report's \texttt{article.tex} was the blob
\nolinkurl{7d5cfe644e7c902f2b594179c337a6c633350fc3}, the same snapshot
against which source~02 was written. Source~03 therefore saw
\cref{isc:sec:introduction}--\cref{isc:sec:nonclaims} (source~01) but not
\cref{isc:cp:sec:intro}--\cref{isc:cp:sec:verify} (source~02), which were
merged afterwards; the two later sources are independent of each other. The
source takes Research question~3 of \cref{isc:sec:questions} as its target.

Section~$n$ of the source is Section~$n+27$ here for $n=1,\ldots,13$, and
every numbered statement keeps its position: the source's Lemma~$n.m$ is
Lemma~$(n+27).m$ here, and likewise for theorems, propositions,
corollaries, definitions and remarks, and for equations. The source's
Theorems~A, B, C and~D share the lettered counter of this report's
Theorems~A--C and are therefore printed as \cref{isc:hj:main:mixed},
\cref{isc:hj:main:cardinal}, \cref{isc:hj:main:prime} and
\cref{isc:hj:main:class}: its Theorem~A is Theorem~D here, B is~E, C is~F
and D is~G. Its Questions~1--10 are Research questions~23--32. Its
Appendices~A--C form \cref{isc:hj:sec:verify}; its Appendix~D (sources and
provenance) is folded into Appendix~\ref{isc:hj:app:source}. Where the
source cites this report, the citation is replaced by a cross-reference to
the same place. Text added in the merge is marked \mtag; symbols renamed from
the source are listed in \cref{isc:hj:sec:conventions}. Every result, proof,
example, question and limitation of the source is printed.

\subsection{Finite field operations versus Hahn summation}

Suppose two additive subgroups $A,B\subseteq\No$ have been chosen. Write,
as in \cref{isc:sec:introduction},
\[
 \K{k}{A}=\left\{\sum_{a\in S}c_a\omega^a:
 c_a\in k,\ S\subseteq A\text{ reverse well ordered and set-sized}\right\},
 \qquad k=\R,\C.
\]
For $k=\C$ this notation denotes the corresponding subfield of $\No(i)$.
Every individual monomial $\omega^{a+b}$ is already a product of an
$A$-monomial and a $B$-monomial. It does not follow that an arbitrary
allowable sum of these products belongs to the field generated by the two
Hahn fields. There are two distinct objects:
\begin{equation}\label{isc:hj:eq:twoobjects}
 M=\K{k}{A}\K{k}{B}
 \quad\subseteq\quad
 E=\K{k}{A+B}.
\end{equation}
The left-hand side is the \emph{ordinary compositum}: each of its elements
uses only finitely many elements of the two fields and finitely many field
operations. The right-hand side is the \emph{full Hahn join}: it allows
every set-sized reverse-well-ordered support in $A+B$.

The distinction is not a simple ``finite sum versus infinite sum''
observation. Each factor field already contains unrestricted Hahn sums on
its own exponent group. Moreover, a geometric series in a mixed monomial
is the inverse of a binomial and therefore belongs to $M$. The problem is
to distinguish genuinely new mixed sums from sums disguised as rational
expressions in arbitrary elements of the factor fields.

This report, in Research question~3 (headed \emph{An exact description of
ordinary composita}), asks both for a characterization of $M$ inside $E$
and for the transcendence degree of $E/M$ in natural examples. Its
compositum-gap argument (\cref{isc:main:join,isc:thm:joinobstruction}) uses
separated scales and infinitely many independent coefficients. Here the
main mechanism is different: a finite differential annihilator works
directly on the exponent group and does not require an iterated-series
presentation. The degree question is solved in explicit families, including
one with Archimedean, interleaved exponent groups.

\subsection{Four principal statements}

\begin{mainthm}[Mixed-support independence]\label{isc:hj:main:mixed}
Let $k$ be a field of characteristic zero, $\Gamma_0,\Gamma_1$ divisible
ordered abelian subgroups of a common ordered abelian group, and
$H=\Gamma_0\cap\Gamma_1$. Suppose a set-indexed family
$(a_i,b_i)_{i\in I}$ satisfies:
\begin{enumerate}[label=\textup{(\roman*)}]
\item $a_i\in\Gamma_0$, and the classes $a_i+H$ are
$\Q$-linearly independent;
\item $b_i\in\Gamma_1$, and $b_i+H$ are pairwise distinct;
\item $\{a_i+b_i:i\in I\}$ is well ordered.
\end{enumerate}
Let $\mathcal U$ be a family of subsets of $I$ such that, for every finite
list of distinct $U_1,\ldots,U_s\in\mathcal U$, each
$U_j\setminus\bigcup_{\ell\ne j}U_\ell$ is infinite. Then the series
\[
 F_U=\sum_{i\in U}t^{a_i+b_i}\qquad(U\in\mathcal U)
\]
are algebraically independent over
$k((t^{\Gamma_0}))k((t^{\Gamma_1}))$ inside
$k((t^{\Gamma_0+\Gamma_1}))$.
\end{mainthm}

The change from $\omega$ to $t$ in this theorem is intentional:
$t^\gamma=\omega^{-\gamma}$, so well-ordered $t$-supports correspond to
reverse-well-ordered Conway supports. The theorem is proved in
\cref{isc:hj:sec:mixed}. Its hypotheses concern the exponents of the
proposed witnesses, not the supports of the coefficients of a putative
relation.

\begin{mainthm}[The largest possible set-sized gap]\label{isc:hj:main:cardinal}
For every infinite cardinal $\kappa$ there are explicit divisible
subgroups $A_\kappa,B_\kappa\subseteq\No$, each of cardinality $\kappa$,
with $A_\kappa\cap B_\kappa=0$, such that, for $k=\R,\C$,
\[
 \trdeg\!\left(
 \K{k}{A_\kappa+B_\kappa}
 \big/
 \K{k}{A_\kappa}\K{k}{B_\kappa}
 \right)=2^\kappa.
\]
The two factor fields are linearly disjoint over $k$. For $k=\R$ they
are real closed; for $k=\C$ they are algebraically closed. There is a
witnessing independent family of $2^\kappa$ elements of $\Oz$, valid
simultaneously for both coefficient fields.
\end{mainthm}

All three Hahn fields in this statement have cardinality $2^\kappa$.
Consequently the lower bound is not a cardinality comparison between a
small base and a larger ambient field. It asserts maximal transcendence
above a base that already has the same cardinality.

\begin{mainthm}[An interleaved family of omnific primes]\label{isc:hj:main:prime}
Let $(p_n)_{n\geq0}$ enumerate the ordinary positive primes increasingly,
and put
\[
 \lambda_n=p_{2n}^{-3/2},\qquad
 \mu_n=p_{2n+1}^{-3/2},\qquad \rho_n=\lambda_n+\mu_n,
 \qquad A=\Span\{\lambda_n\},\quad B=\Span\{\mu_n\}.
\]
There is an explicit family $(\mathfrak B_\eta)_{\eta\in2^{\N}}$ of infinite,
pairwise almost-disjoint subsets of $\N$ for which
\begin{equation}\label{isc:hj:eq:primefamily-intro}
 P_\eta=1+\sum_{n\in\mathfrak B_\eta}\omega^{\rho_n}
\end{equation}
are algebraically independent over $\K{\C}{A}\K{\C}{B}$.
Each $P_\eta$ is prime in both $\Oz$ and $\Oz[i]$. Moreover,
\[
 \trdeg\bigl(\K{k}{A+B}/\K{k}{A}\K{k}{B}\bigr)=\cc
 \quad(k=\R,\C),\qquad \cc=2^{\aleph_0}.
\]
Here $A+B\subseteq\R$ is Archimedean and $A\cap B=0$.
\end{mainthm}

The primality input is classical, from Pitteloud and the transfer work of
Biljakovic--Kochetov--Kuhlmann; it is also available through the broader
criteria of L'Innocente--Mantova \cite{Pitteloud,BKK,LM}. The proposed
contribution in \cref{isc:hj:main:prime} is the simultaneous maximal
algebraic independence over two entire Hahn subfields, together with the
explicit prime realization, not a new one-row primality theorem.

\begin{mainthm}[One fixed class-sized example]\label{isc:hj:main:class}
There are explicit proper-class divisible subgroups $A_*,B_*\subseteq\No$
with $A_*\cap B_*=0$ such that the full Hahn join
$\K{k}{A_*+B_*}$ contains an ordinal-indexed, algebraically independent
family of omnific integers over $\K{k}{A_*}\K{k}{B_*}$, for
$k=\R,\C$. The extension therefore has no set-sized transcendence basis.
For every infinite cardinal $\kappa$ it contains a set of $2^\kappa$
algebraically independent omnific integers over that same class compositum.
\end{mainthm}

This last statement concerns explicitly defined class-sized support
subfields. It does \emph{not} identify either factor with a full copy of
$\No$, and does not use this report's construction of such copies
(\cref{isc:main:copies}), which the source treats as unreviewed.
Algebraic independence of a class always means that every finite
subfamily is algebraically independent.

\subsection{How to read the source}

The essential proof is \cref{isc:hj:sec:lemmas,isc:hj:sec:mixed}. Readers
mainly interested in concrete surreal numbers may then go directly to
\cref{isc:hj:sec:archimedean,isc:hj:sec:prime}. Cardinal amplification and
the class example are in
\cref{isc:hj:sec:combinatorics,isc:hj:sec:cardinal,isc:hj:sec:classes}.
The exact scope of the advance over this report, failure cases, and
further questions are recorded in
\cref{isc:hj:sec:scope,isc:hj:sec:questions}. \Cref{isc:hj:sec:verify}
gives a proof-dependency audit, finite exact verification, and a
formalization roadmap.

\subsection{What the addition settles in this report}
\label{isc:hj:sec:answers}

\mtag{}
\begin{enumerate}[leftmargin=*,itemsep=3pt]
\item \emph{Research question~\ref{isc:cp:q:amalgams}} (non-separated
ordered amalgams) asks for the image of $\K{k}{G_0}\K{k}{G_1}$ in
$\K{k}{G_0+G_1}$, for its transcendence defect, and for a finite-rank
denominator criterion without an iterated Hahn coefficient field. Its
\emph{transcendence-defect part} is answered for explicit non-separated
examples: for the interleaved groups $A_\kappa,B_\kappa$ of
\cref{isc:hj:main:cardinal} (neither is infinitesimal relative to the
other: $f_0\in B_\kappa$ lies between $e_0$ and $e_1$ of $A_\kappa$, and so on)
the defect is $2^\kappa$ at every infinite cardinal $\kappa$; for the dense
Archimedean groups of \cref{isc:hj:main:prime} it is $\cc$; for the class
groups of \cref{isc:hj:main:class} there is no set-sized transcendence
basis. The image of the compositum and a denominator criterion are
\emph{not} obtained; they remain open, and the source's own
Research question~\ref{isc:hj:q:characterization} asks for them.
\item \emph{Research question~3}, second part (the transcendence degree of
the join over the compositum in natural examples), was answered in
batch~34 in the separated case and the Laurent cardinal cases. Source~03
adds the explicit non-separated examples just listed, with omnific and
prime witnesses. The first part (an exact description of the compositum)
remains answered only in the separated case (\cref{isc:cp:thm:rank}).
\item \emph{Research question~\ref{isc:cp:q:cores}} (prescribed nonzero
common cores) is continued, not answered. \Cref{isc:hj:main:mixed}
allows a nonzero intersection $H=\Gamma_0\cap\Gamma_1$ algebraically
(hypotheses (i) and~(ii) are modulo $H$), but every example constructed has
$H=0$, and no tail resilience is proved; the ordered construction for
$H\ne0$ is the source's Research question~\ref{isc:hj:q:cores}.
\item \Cref{isc:main:join} receives a second route: the witness $S$ of
\eqref{isc:eq:explicitS} is an instance of \cref{isc:hj:main:mixed} in the
class-local form of \cref{isc:hj:lem:classlocal}
(\cref{isc:hj:rem:secondroute}).
\item No other question of this report is answered. Research
questions~\ref{isc:cp:q:scales}, \ref{isc:cp:q:formal} and~10 are continued
by Research questions~\ref{isc:hj:q:factors} and~\ref{isc:hj:q:formal};
Research question~3 and~\ref{isc:cp:q:cores} by
Research questions~\ref{isc:hj:q:cores} and~\ref{isc:hj:q:copies}.
\end{enumerate}

\subsection{Conventions and renamed symbols}
\label{isc:hj:sec:conventions}

\mtag{} Source~03 uses both exponent conventions of this report: the
decreasing large-$\omega$ convention $\sum c_g\omega^g$, written $\K{k}{G}$
as here, and the increasing small-$t$ convention $k((t^\Gamma))$ with
well-ordered support, identified by $t^\gamma=\omega^{-\gamma}$ as in
\cref{isc:sub:bounded,isc:cp:sec:foundations}. In
\cref{isc:hj:sec:lemmas,isc:hj:sec:mixed} and in \cref{isc:hj:main:mixed}
the fields are abstract $t$-series fields over an arbitrary
characteristic-zero field $k$; from \cref{isc:hj:sec:cardinal} on they are
subfields of $\No$ or $\SC$, with $k=\R$ or $\C$.

The following symbols of the source are renamed; everything else is as in
the source. No normalization is changed.
\begin{center}\small
\begin{longtable}{@{}>{\raggedright\arraybackslash}p{0.29\textwidth}>{\raggedright\arraybackslash}p{0.19\textwidth}>{\raggedright\arraybackslash}p{0.44\textwidth}@{}}
\toprule
Here & Source 03 & Reason\\
\midrule
\endhead
$M=K_0K_1$, the ordinary compositum; $M'$, $M^{\mathrm{alg}}$, $M_*$ & $L$;
$L'$, $L^{\mathrm{alg}}$, $L_*$ & this report writes $M$ for the ordinary
compositum (\cref{isc:main:join}, source~02); $L$ is the common core
$L_k=\K{k}{H}$ of \cref{isc:main:copies} and, in
\cref{isc:cp:sec:foundations}, the field $k((t^\Gamma))$\\
$\Gamma_0,\Gamma_1$; $K_0=k((t^{\Gamma_0}))$, $K_1=k((t^{\Gamma_1}))$;
$K_{0,*},K_{1,*}$ & $\Gamma_1,\Gamma_2$; $K_1,K_2$; $K_{1,*},K_{2,*}$ &
this report indexes the two factors of a join by $0,1$ ($G_0,G_1$,
$K_0,K_1$ in \cref{isc:sec:join}); the first factor, over which the
annihilator is solved, plays the role of the coefficient field $K_0$ there\\
$\K{k}{G}$ & $\mathcal H_k(G)$ & the same full Hahn field; this report's
macro\\
$\partial_\chi$, $\partial_i$ (Euler derivations);
$\partial_{\mathbf c}=\sum c_i\partial_i$ (the annihilator) & $D_\chi$,
$D_i$; $D$ & $D$ is always a field in \cref{isc:cp:sec:conventions}, and
$D_\R=\Z$, $D_\C=\Z[i]$ are rings of integers; $\partial_\chi$ is the
Euler derivation \eqref{isc:cp:eq:euler} with weight $\lambda=\chi$\\
$\delta$, an ordinary relative derivation (\cref{isc:hj:sec:strong}) & $D$
& as above\\
$\Gamma'$, a subgroup in \cref{isc:hj:lem:cosets,isc:hj:lem:euler} &
$\Delta$ & $\Delta$ is the Euler derivation of \cref{isc:cp:sub:euler}\\
$\PiG{k}{G}$, purely infinite series; $\A{k}{G}=D_k\oplus\PiG{k}{G}$ &
$I_k(G)$; $R_k(G)$ & the same objects as in \cref{isc:sec:fractions}; $I$
is an index set and $I_\alpha$ an interval here, and $R_j$ are omnific
copies\\
$\Ad{D}{k}{G}=D\oplus\PiG{k}{G}$ for a unital ring $D\subseteq k$, in
particular $\Ad{k}{k}{G}$ & $S$; $T$ (\cref{isc:hj:lem:pullback}) & $S$ is
a support set and the witness \eqref{isc:eq:explicitS}; $T$ is kept for
the local set in \cref{isc:hj:lem:independentfamily}\\
$\A{k}{A_\kappa}$, $\A{k}{B_\kappa}$, $\A{k}{A_\kappa+B_\kappa}$;
$\A{k}{A_\kappa}[\A{k}{B_\kappa}]$ & $R_{A,k}$, $R_{B,k}$, $R_{E,k}$; $S_k$
& as above; the ring generated by two rings is written as
$R_0[R_1]$ in \cref{isc:cp:sec:arithmetic}\\
$\mathcal T$, a transcendence basis; $\varphi:\mathcal T\to E$; $dy$,
$y\in\mathcal T$ & $T$; $h$; $dt$, $t\in T$ & as source~02 does
(\cref{isc:cp:lem:partial}); $h$, $h_\alpha$ are embeddings here; $t$ is
the Hahn variable\\
$\mathcal I_\kappa$, the index set of \cref{isc:hj:lem:independentfamily}
& $I_\kappa$ & $I_\alpha$ is an interval of \cref{isc:sec:explicit}\\
$\mathfrak B_\eta$, the branch sets & $B_\eta$ & $B$ is the group
$\Span\{\mu_n\}$ in the same theorem\\
$\Xi_U$, the omnific witnesses of \cref{isc:hj:main:cardinal} & $Z_U$ &
$Z_\alpha=(1+i)Y_\alpha$ are source~02's Gaussian witnesses
(\cref{isc:cp:thm:complex})\\
$\Lambda_\alpha$, the ordinal blocks of \cref{isc:hj:sec:classes} &
$C_\alpha$ & $C_\Theta$ is a localized coefficient field
(\cref{isc:cp:lem:localizeP})\\
Theorems D, E, F, G & Theorems A, B, C, D & the lettered counter
continues this report's Theorems A--C\\
Research questions 23--32 & Questions 1--10 & the question counter
continues\\
\bottomrule
\end{longtable}
\end{center}

\paragraph{Symbols kept, and tempting false readings.}
\begin{itemize}[leftmargin=*,itemsep=3pt]
\item $H=\Gamma_0\cap\Gamma_1$ is the common core of the two exponent groups,
in the same role as the core $H$ of \cref{isc:main:copies}; here it is an
arbitrary subgroup of an abstract ordered group, and all examples have
$H=0$. $\Gamma=\Gamma_0+\Gamma_1$ is the exponent group of the join $E$,
as $\Gamma$ is the outer group of $E=F((t^\Gamma))$ in source~02.
\item $A,B$; $A_\kappa,B_\kappa$; $A_*,B_*$ are exponent \emph{subgroups},
as $A,B$ are in \cref{isc:lem:iterated}. They are not the rings
$\A{k}{H}$, $\B{k}{H}$ (calligraphic), not a parameter set $A$
(\cref{isc:cor:parameters}), and not the index set $B$ or the labels
$b_i$ of \cref{isc:cp:sec:generic}.
\item $\chi$, $\chi_i$ are additive \emph{characters} $\Gamma\to k$: the
Euler weight $\lambda$ of \cref{isc:cp:def:euler} under another name. They
are not the characters $\chi$ of the monomial maps $M_{\chi,\tau}$ of the
omnific-preserving report.
\item $\lambda_n,\mu_n,\rho_n$ (\cref{isc:hj:main:prime}) are positive
\emph{real numbers} used as exponents. $\lambda_n$ is not the Euler weight
$\lambda$ of \cref{isc:cp:def:euler} or a coset slice $\lambda_{g,H}$ of
\eqref{isc:eq:slice}; $\mu_n$ is not $\mu=|F|$ of
\cref{isc:cp:thm:cardinal}; $\rho_n$ is not the embedding $\rho$ of the
omnific-preserving classification (\cref{isc:sec:collection}). $p_n$ is
the $n$th prime, not an evaluated partial $p_{ij}$ of
\cref{isc:cp:thm:lacunary}.
\item $e_\alpha=\omega^{-\iota_0(\alpha)}$, $f_\alpha=\omega^{-\iota_1(\alpha)}$
and $r_\alpha=e_\alpha+f_\alpha$ are infinitesimal surreal \emph{exponents}
(for $\alpha\ge1$; $e_0=1$). Here $\iota_0(\alpha)=2\cdot_{\Ord}\alpha$ is an
\emph{ordinal} product. \emph{Tempting false reading:}
``$e_\alpha=\omega^{-2\alpha}$ with the surreal product $2\alpha$.'' For
infinite $\alpha$ the two differ: $2\cdot_{\Ord}\omega=\omega$, whereas the
surreal $2\omega=\omega+\omega$.
\emph{Tempting false reading:} ``$r_\alpha$ is a real coefficient, as
$r_g$ in $\sum r_g\omega^g$.'' It is not; $r_\alpha$ is itself a Conway
exponent of the witnesses $\Xi_U$ and $W_\alpha$. Nor is $e$ the generic
exponent of \cref{isc:cp:sec:surreal}.
\item $\iota_0(\alpha)=2\cdot_{\Ord}\alpha$ and $\iota_1(\alpha)$ are
ordinal-index maps, not the inner lifts $\iota_{\sigma_\alpha}$ of the
omnific-preserving report.
\item $F_U$, $F_\eta$ are mixed sums indexed by an incidence set or a
branch; they are not the witness $F$ of \cref{isc:main:fractions}, the
coefficient field $F$ of source~02, or a polynomial. $P_\eta$ is a prime
witness indexed by a binary branch; $P$ without index is a polynomial
relation, and $P_\alpha$ are the partial sums of \cref{isc:sec:gaps}.
$W_\alpha$ (\cref{isc:hj:eq:classwitness}) is an omnific witness, not a
leading polynomial $W_i(X,x)$ of \cref{isc:cp:thm:lacunary}.
\item $\eta,\theta\in2^{\N}$ are binary branches, as $\eta$ is in
\cref{isc:cp:cor:continuum}; they are not the $\eta_\kappa$-orders of
\cref{isc:sec:setversion}, and $\theta$ is not the inner support set
$\Theta$. $s$ is a finite subset of $\kappa$ or a binary word, and in
\cref{isc:hj:thm:classicalprime} $s_n$ are real $t$-exponents, not the
coefficients $s_n=\omega^{g_n}$ of \eqref{isc:eq:explicitS}. $\delta$ is a
derivation here, not the Taylor parameter of the omnific-preserving
report; $\delta_{ih}$ is the Kronecker symbol. $\kappa$ is an arbitrary
infinite cardinal, not $\cf(H)$, and $\lambda$ without index is a cardinal
in the proofs of \cref{isc:hj:main:class,isc:hj:thm:automorphisms}.
\item $D$ is a unital subring of $k$ in \cref{isc:hj:sec:prime}, in
particular $D_\R=\Z$ or $D_\C=\Z[i]$ as in \cref{isc:sec:fractions}; it is
never a derivation (source~02's rule that $D$ is a field applies to its own
sections). $q$ is an index of the generators $u_q$ in
\cref{isc:hj:lem:annihilator} and a prime series in
\cref{isc:hj:sec:prime}; $g$ is a generic exponent in the proof of
\cref{isc:hj:prop:strongrigidity}, and $\omega^{7/2}$ only in
\cref{isc:hj:rem:secondroute}.
\item \emph{Tempting false reading:} ``the annihilator
$\partial_{\mathbf c}$ is a derivation of $E$ killing the compositum $M$.''
It kills only a finitely generated subfield $K_1(u_1,\ldots,u_m)$ chosen
after a hypothetical relation is fixed; a nonzero strong derivation
killing all of $M$ does not exist (\cref{isc:hj:prop:strongrigidity}).
\item \emph{Tempting false reading:} ``the automorphisms of
\cref{isc:hj:thm:automorphisms} are automorphisms of $\No(i)$ or of
$\Oz$.'' They are automorphisms of one set-sized complex Hahn join $E$
fixing $M$, and nonidentity ones are not strong.
\item \emph{Strong} means preserving Hahn summability of set-indexed
families, as in \cref{isc:sec:conventions}; ordinal operations are written
$\cdot_{\Ord}$, $+_{\Ord}$ as in the source.
\end{itemize}

\subsection{Place in the collection}\label{isc:hj:sec:collection}

\mtag{} Source~03 cites only this report from the collection (and the
catalogue, at its pin). Checked against the collection at the placement
commit \texttt{dec8d56}, the following reports bear on it; none is a
dependency of its proofs.
\begin{enumerate}[leftmargin=*,itemsep=3pt]
\item \emph{This report, source~01.} \Cref{isc:hj:prop:ld} is the case
$H=0$ of \cref{isc:thm:sliceddisjoint} (the intersection is $k$ and the
fields are linearly disjoint over it), by the same coefficientwise
argument; \cref{isc:hj:lem:cosets} is the elementary coset fact behind the
slices \eqref{isc:eq:slice}. The definitions of $\PiG{k}{G}$ and
$\A{k}{G}$ in \cref{isc:hj:sec:foundations} are those of
\cref{isc:sec:fractions}. The source's comparison with source~01 is
accurate: \cref{isc:thm:joinobstruction} uses separated scales and an
infinite independent coefficient list, and the source does not present
that mechanism as new.
\item \emph{This report, source~02} (\cref{isc:hj:sub:compare}). Written
from the same snapshot, source~02 answered Research question~3 in the
separated case; the two sources overlap in no theorem.
\item \emph{Tail spans and differential transcendence}
(\rpath{surreal/tail-spans-and-differential-transcendence}, prefix
\texttt{tail:}). The binary-branch sets \eqref{isc:hj:eq:branchsets} are the
almost disjoint family \lbl{tail:eq:Aalpha} with the coding
\lbl{tail:eq:binarycode} shifted by one: the source's code of a word $s$ of
length $m$ is $2^m-1+\sum_js_j2^{m-1-j}$, that report's is
$N(s)=2^m+\sum_js_j2^{m-j}$ with $1$-based indices, so the source's code is
$N(s)-1$. The same family serves \cref{isc:cp:cor:continuum}. The
linear independence of square roots of distinct primes, used for the
exponents \eqref{isc:hj:eq:realexponents}, is also the multiquadratic input
of \lbl{tail:thm:surreal}, where the square roots are coefficients, not
exponents. Source~03 cites neither.
\item \emph{Omnific Diophantine geometry}
(\rpath{surreal/omnific-diophantine-geometry}, prefix \texttt{odg:}). After
\lbl{odg:prop:irreduciblect} that report cites the primality of
$1+\sum_{k\ge1}\omega^{1/k}$ in $\Oz$ from the same sources (Pitteloud;
Biljakovic--Kochetov--Kuhlmann; as reviewed in \cite[Section~1.3]{LM}).
\Cref{isc:hj:prop:fullprime} is the same one-row input with the
constant-term pullback written out, and the $P_\eta$ are a new explicit
continuum subfamily of such primes. A search of the collection found no
other nonconstant prime element of $\Oz[i]$ (the set-sized-quotients report
treats the ideals generated by ordinary Gaussian primes such as $1+i$), so
primality of the $P_\eta$ in $\Oz[i]$ appears here first in the collection;
this is not a priority claim. Every $P_\eta$ has constant term $1$,
as \lbl{odg:prop:irreduciblect} requires of a nonconstant irreducible.
\item \emph{Set-sized quotients of omnific integers}
(\rpath{surreal/set-sized-quotients-of-omnific-integers}, prefix
\texttt{osq:}). \lbl{osq:cor:primeexample} uses the prime
$\omega^{\sqrt2}+\omega+1$ (L'Innocente--Mantova's Theorem~B) to obtain a
proper-class domain $\Oz/(q)$ with no nonzero set-sized image. Its proof
uses only primality and $\ct(q)=1$, so each $P_\eta$ gives such a quotient;
this is an observation of the merge, not a statement of either source.
\item \emph{Bounded-support transcendence}
(\rpath{surreal/transcendence-over-bounded-support}, prefix
\texttt{bst:}). \lbl{bst:cor:exactcard} and \lbl{bst:cor:realgroup} compute
transcendence degrees $2^\kappa$ and $\cc$ of one full Hahn field over its
bounded-support fraction field. \Cref{isc:hj:main:cardinal,isc:hj:main:prime}
compute the same cardinals for a join over the compositum of two full
Hahn subfields. The bases differ, and the results are consistent: in each
case the degree equals the cardinality of the ambient field. The
Archimedean join $\K{k}{A+B}$ of \cref{isc:hj:main:prime} has
$\cc$ as its degree over both bases.
\item \emph{Omnific-preserving automorphisms} (\texttt{opa:}).
\lbl{opa:as:thm:main} shows that every ring automorphism of $\Oz$ extends
to a strong automorphism of $\No$, and \lbl{opa:as:cor:setpairs} that every
automorphism of a full real Hahn field on a nonzero divisible set group
stabilizing its omnific ring is strong. The non-strong automorphisms of
\cref{isc:hj:thm:automorphisms} are automorphisms of a complex Hahn join,
not of $\Oz$, and are not claimed to preserve any omnific ring, so there is
no conflict; \cref{isc:hj:rem:omnificaut} combines the two results.
\end{enumerate}
No \texttt{isc:hj:} label has a Lean mapping in \rpath{FORMALIZATION.md}.

\subsection{Comparison with source 02}\label{isc:hj:sub:compare}

\begin{remark}[Two answers to the degree question]\label[remark]{isc:hj:rem:compare}
\mtag{} Sources~02 and~03 were written from the same snapshot of this report
and do not cite each other.
\begin{enumerate}[label=\textup{(\alph*)},leftmargin=*,itemsep=2pt]
\item \emph{Settings.} \Cref{isc:cp:thm:cardinal} extends coefficients: it
computes $\dtrdeg=\trdeg=|F|^{\aleph_0}$ for $F((t))$ over $F k((t))$, with
exponents in $\Z$ and the gap coming from the coefficient field $F$.
\Cref{isc:hj:main:cardinal} mixes exponents: both factors are full Hahn
fields over the same $k$ on interleaved groups, and the degree is
$2^\kappa$. Both degrees are maximal, equal to the cardinality of the
ambient field; they agree in value for $\kappa=\aleph_0$ and $|F|=\cc$.
\item \emph{Mechanisms.} Source~03's finite-exclusive infinitude
(\cref{isc:hj:def:exclusive}) is word for word the finite private-index
property \cref{isc:cp:def:private}, and the private-label hypothesis of
\cref{isc:cp:thm:generic} is its label form. Source~02 kills a finitely
generated coefficient field with coordinate \emph{coefficient} derivations
$d_b$ (\cref{isc:cp:lem:partial}); source~03 kills finitely many elements
of the factor $K_0$ with a $K_0$-linear combination of exponent-coordinate
Euler derivations $\partial_i$ (\cref{isc:hj:lem:annihilator}), which plays
the same role. In the second route of \cref{isc:hj:rem:secondroute} the
$\partial_i$ act on the outer coefficients $s_n=\omega^{g_n}$ of
\eqref{isc:eq:explicitS} as $s_n\,\partial/\partial s_n$, a logarithmic form
of the coefficient derivations $d_b$.
\item \emph{Scope.} Source~02 proves differential independence, tail
resilience over every set of parameters and arithmetic nonintegrality, in
the \emph{separated} explicit join. Source~03 proves algebraic independence
only, but without convex separation, with an exact set-sized degree,
prime witnesses and a count of relative derivations and automorphisms.
\item \emph{Class size.} \Cref{isc:cp:thm:omnific,isc:cp:cor:nosetgen} work
in the join of two actual copies of $\No$ (\cref{isc:sec:explicit});
\cref{isc:hj:main:class} works in explicit class \emph{support subfields}
$\K{k}{A_*}$, $\K{k}{B_*}$, which are not copies of $\No$, as the source
says.
\item \emph{Coverage.} \Cref{isc:hj:main:mixed} as stated does not cover
source~02's family $(Y_\alpha)$: in the labels $(\alpha,n)$ its
$G_1$-exponents $b-ng$ do not depend on $\alpha$, so hypothesis~(ii) fails
for the family, although each single $Y_\alpha$ satisfies it.
\end{enumerate}
\end{remark}

\section{Conventions and foundational inputs of the third source}
\label{isc:hj:sec:foundations}

\subsection{Hahn fields and summable families}

Throughout the abstract argument, $k$ is a set-sized field of
characteristic zero and $\Gamma$ is a set-sized divisible ordered abelian
group. Divisibility makes $\Gamma$ a $\Q$-vector space, uniquely, because
an ordered abelian group is torsion free. Its Hahn field consists of
formal sums
\[
 k((t^\Gamma))=
 \left\{x=\sum_{\gamma\in\Gamma}x_\gamma t^\gamma:
 x_\gamma\in k,\ \supp(x)=\{\gamma:x_\gamma\ne0\}
 \text{ is well ordered}\right\}.
\]
Addition is coefficientwise and multiplication is convolution. The
classical Hahn support theorem ensures that the sum of two well-ordered
supports is well ordered and that each exponent has only finitely many
representations from the two supports. It follows that every convolution
coefficient is an ordinary finite sum. Inversion makes this a field.
These are standard foundations, not results proposed as new here
\cite{Hahn,Gonshor}.

A set-indexed family $(x_j)$ is \emph{Hahn summable} if its union of
supports is well ordered and every exponent belongs to only finitely many
of the supports. Its Hahn sum is then defined coefficientwise. No
analytic convergence or passage to a topological limit is intended.
In particular, the families of distinct monomials in \cref{isc:hj:main:mixed}
are summable because their total support is well ordered.

The natural valuation is $v(x)=\min\supp(x)$ for $x\ne0$. It is useful for
orientation but does not enter the annihilator argument. In particular, no
assumption about the cofinality of the valuation group is needed.

\subsection{The Conway realization and the two coefficient roles}

Conway normal form identifies $\No$ with real-coefficient, set-supported
reverse-well-ordered sums $\sum r_g\omega^g$. Sending $t^\gamma$ to
$\omega^{-\gamma}$ identifies a Hahn field on a subgroup of $\No$ with
a subfield of $\No$. Complex coefficients are handled through
$\No(i)=\No\oplus i\No$. The union of the real and imaginary supports
is again reverse well ordered, so this agrees with the formal complex
Hahn presentation. The usual real-closedness theorem gives that
$\K{\R}{G}$ is real closed when $G$ is divisible, and its
complexification $\K{\C}{G}$ is algebraically closed
\cite{Gonshor,Kaplansky}.

An exponent that happens to be a real number is not a coefficient.
For example, $\lambda_n$ in \cref{isc:hj:main:prime} is fixed when it is
regarded as an element of the coefficient field, while the monomial
$t^{\lambda_n}$ can have a nonzero Euler derivative. This distinction is
essential. The derivations below are not assumed to respect a surreal
exponential function, and are not the Berarducci--Mantova surreal
derivation. They are additive-character derivations of a Hahn field.

\subsection{Omnific and Gaussian omnific rings}

For a subgroup $G\subseteq\No$ the additive group of purely infinite series
is, as in \cref{isc:sec:fractions},
\[
 \PiG{k}{G}=\left\{\sum_{g\in S}c_g\omega^g\in\K{k}{G}:
 S\subseteq G^{>0}\right\}.
\]
The product of two such series again has strictly positive exponents.
Thus $k\oplus\PiG{k}{G}$ is a ring and its constant-term map to $k$ is a
ring homomorphism. For $k=\R$ and $k=\C$ respectively,
\[
 \A{\R}{G}=\Z\oplus\PiG{\R}{G}=\Oz\cap\K{\R}{G},
 \qquad
 \A{\C}{G}=\Z[i]\oplus\PiG{\C}{G}=\Oz[i]\cap\K{\C}{G}.
\]
These descriptions are exact. A real surreal is omnific precisely when
it has no negative Conway exponent and its constant coefficient is an
ordinary integer \cite{Conway,Gonshor}.

We will not use the false blanket assertion
$\Frac(\A{k}{G})=\K{k}{G}$ for arbitrary set-sized $G$.
Even though every surreal is a quotient of full omnific integers, this
identity need not remain true after imposing a fixed exponent group.
The arithmetic corollary in \cref{isc:hj:sec:arithmetic} only uses
inclusions that follow directly from the definitions.
\mtag{} This report proves the failure: for $H\ne0$ the field
$\Frac\A{k}{H}$ is $\Frac\B{k}{H}$ (\cref{isc:prop:boundedfrac}), over which
$\K{k}{H}$ is transcendental (\cref{isc:thm:cofinalgap}).

\subsection{Set theory and class statements}

\Cref{isc:hj:sec:foundations} through \cref{isc:hj:sec:prime} use
ordinary set mathematics with choice. The optional class results can be
read in a class theory such as NBG, the framework of
\cref{isc:sec:foundations}. Every Hahn support is still a set. The class
exponent groups used there have explicit finite-support coordinate
descriptions, so no proper-class Hamel basis is selected. Finite linear
algebra and finite polynomial identities over a class field have their
usual meanings. Where a published theorem is stated for set groups, the
application to a class ring will be localized to a set group containing all
relevant supports.

\section{Three algebraic lemmas}\label{isc:hj:sec:lemmas}

Whenever the data of \cref{isc:hj:main:mixed} are used in this section, write
$\Gamma=\Gamma_0+\Gamma_1$, $K_j=k((t^{\Gamma_j}))$,
$E=k((t^\Gamma))$, and $H=\Gamma_0\cap\Gamma_1$.

\subsection{Distinct exponent cosets cannot cancel}

\begin{lemma}[Coset independence]\label[lemma]{isc:hj:lem:cosets}
Let $\Gamma'\subseteq\Gamma$ be an additive subgroup. If
$\beta_1,\ldots,\beta_r$ represent distinct cosets modulo $\Gamma'$, then
$t^{\beta_1},\ldots,t^{\beta_r}$ are linearly independent over
$k((t^{\Gamma'}))$ inside $k((t^\Gamma))$.
\end{lemma}
\begin{proof}
If $c_j\in k((t^{\Gamma'}))$, then
$\supp(c_jt^{\beta_j})\subseteq\Gamma'+\beta_j$.
These cosets are pairwise disjoint. For a nonzero $c_j$, any exponent of
$c_jt^{\beta_j}$ has a nonzero coefficient that cannot receive a
contribution from another summand. Hence
$\sum_jc_jt^{\beta_j}=0$ forces every $c_j=0$.
\end{proof}

No convexity assumption appears in this argument. Even very closely
interleaved cosets remain disjoint as sets of exponents. This is the
replacement for a separated-scale coefficient expansion.

\begin{proposition}[Intersection and linear disjointness]\label[proposition]{isc:hj:prop:ld}
Suppose $\Gamma_0\cap\Gamma_1=0$, and put
$K_j=k((t^{\Gamma_j}))$. Then $K_0\cap K_1=k$, and $K_0,K_1$ are
linearly disjoint over $k$ in $k((t^{\Gamma_0+\Gamma_1}))$.
\end{proposition}
\begin{proof}
An element of the intersection has support contained in both groups,
hence in $\{0\}$, which proves the first assertion. To prove the second,
let $y_1,\ldots,y_r\in K_1$ be linearly independent over $k$, and assume
$\sum_j x_jy_j=0$ with $x_j\in K_0$.
Every exponent in $\Gamma_0+\Gamma_1$ has a unique expression $a+b$ with
$a\in\Gamma_0$, $b\in\Gamma_1$. For a fixed $a$, the coefficients in
the coset $a+\Gamma_1$ of the asserted relation therefore give
\[
 \sum_{j=1}^r (x_j)_a\,y_j=0.
\]
The independence of the $y_j$ implies $(x_j)_a=0$ for every $j$. As $a$
was arbitrary, all $x_j$ vanish. This is the linear-disjointness criterion.
\end{proof}

\mtag{} This is the case $H=0$ of \cref{isc:thm:sliceddisjoint}
(\cref{isc:hj:sec:collection}).

\subsection{Coordinate Euler derivations}

\begin{lemma}[An additive character gives a derivation]\label[lemma]{isc:hj:lem:euler}
Let $\chi:\Gamma\to k$ be an additive map. Then
\begin{equation}\label{isc:hj:eq:euler}
 \partial_\chi\left(\sum_\gamma x_\gamma t^\gamma\right)
 =\sum_\gamma\chi(\gamma)x_\gamma t^\gamma
\end{equation}
defines a $k$-linear derivation of $k((t^\Gamma))$. It preserves Hahn
summability and commutes with Hahn sums. If $\chi$ vanishes on a subgroup
$\Gamma'$, it annihilates $k((t^{\Gamma'}))$.
\end{lemma}
\begin{proof}
The support of the output is contained in that of the input, so the
series is well defined. Additivity and $k$-linearity are immediate.
For a product, its coefficient at $\gamma$ is a finite convolution sum.
Using $\chi(\alpha+\beta)=\chi(\alpha)+\chi(\beta)$ in that finite sum
shows
\[
 \partial_\chi(xy)=\partial_\chi(x)y+x\partial_\chi(y).
\]
For a summable family, its union of supports cannot increase under
$\partial_\chi$, and coefficientwise local finiteness cannot fail. The
equality between differentiating the sum and summing the derivatives
reduces to linearity in a finite coefficient sum. The last assertion
follows from \eqref{isc:hj:eq:euler}.
\end{proof}

\mtag{} This is \cref{isc:cp:lem:derivation} for $F=k$, with the Euler
weight $\lambda=\chi$.

\begin{lemma}[Coordinate characters modulo the other group]\label[lemma]{isc:hj:lem:characters}
Under hypotheses \textup{(i)} and \textup{(ii)} of \cref{isc:hj:main:mixed},
for each $i\in I$ there is a $\Q$-linear map
$\chi_i:\Gamma_0+\Gamma_1\to\Q\subseteq k$ such that
\[
 \chi_i\vert_{\Gamma_1}=0,\qquad \chi_i(a_h)=\delta_{ih}
 \quad(h\in I).
\]
Writing $\partial_i=\partial_{\chi_i}$, one has $\partial_i(K_0)\subseteq K_0$,
$\partial_i(K_1)=0$, and
\begin{equation}\label{isc:hj:eq:coordinateaction}
 \partial_i(F_U)=\begin{cases}t^{a_i+b_i}&i\in U,\\0&i\notin U.\end{cases}
\end{equation}
\end{lemma}
\begin{proof}
The classes of the $a_h$ are linearly independent in
$(\Gamma_0+\Gamma_1)/\Gamma_1$: a finite rational combination lying in
$\Gamma_1$ also lies in $\Gamma_0$, hence in $H$, so all coefficients
vanish. Extend these classes to a $\Q$-basis of the quotient, assign value
$1$ to the class of $a_i$ and $0$ to the other basis vectors, and compose
with the quotient map. Apply \cref{isc:hj:lem:euler}. The value on
$a_h+b_h$ is $\delta_{ih}$, which gives \eqref{isc:hj:eq:coordinateaction}.
\end{proof}

Only hypothesis (i) is needed to construct these characters; hypothesis
(ii) will ensure that a linear combination of the resulting nonzero
monomial derivatives does not cancel.

\subsection{Annihilating finitely many arbitrary Hahn series}

\begin{lemma}[Finite differential annihilator]\label[lemma]{isc:hj:lem:annihilator}
Let $u_1,\ldots,u_m\in K_0$, and let $J\subseteq I$ consist of exactly $m+1$
distinct indices. There are coefficients $(c_i)_{i\in J}\in K_0^J$, not
all zero, such that the derivation
\begin{equation}\label{isc:hj:eq:annihilator}
 \partial_{\mathbf c}=\sum_{i\in J}c_i\partial_i
\end{equation}
annihilates $K_1(u_1,\ldots,u_m)$. In addition,
\begin{equation}\label{isc:hj:eq:noncancel}
 \sum_{i\in J}c_i t^{a_i+b_i}\ne0.
\end{equation}
\end{lemma}
\begin{proof}
Consider the homogeneous equations
\[
 \sum_{i\in J}c_i\partial_i(u_q)=0\qquad(1\leq q\leq m).
\]
There are $m$ equations in $m+1$ unknowns over the field $K_0$, because
$\partial_i(u_q)\in K_0$. Gaussian elimination gives a nonzero solution.
A finite field-linear combination of derivations is again a derivation,
so \eqref{isc:hj:eq:annihilator} is well defined. It annihilates every
$u_q$ and all of $K_1$. The quotient rule then implies that it annihilates
the field they generate.

For \eqref{isc:hj:eq:noncancel}, write the left side as
$\sum_{i\in J}(c_it^{a_i})t^{b_i}$.
The cosets $b_i+\Gamma_0$ are pairwise distinct: a difference in
$\Gamma_0$ would be in $H$, contrary to hypothesis (ii).
The coefficients $c_it^{a_i}$ lie in $K_0$, and some are nonzero.
Now apply \cref{isc:hj:lem:cosets}.
\end{proof}

\begin{remark}[Why arbitrary coefficient supports are harmless]\label[remark]{isc:hj:rem:finiteness}
The $u_q$ may each have infinitely many, or uncountably many, support
exponents. We do not choose coordinates outside all their supports. We
choose one more differential direction than there are generators, and
solve a finite homogeneous system over $K_0$ itself. This is the principal
reason the proof applies to a compositum of \emph{full} Hahn fields.
\end{remark}

\section{The mixed-support independence theorem}\label{isc:hj:sec:mixed}

\begin{definition}\label[definition]{isc:hj:def:exclusive}
A family $\mathcal U$ of subsets of $I$ has the \emph{finite-exclusive
infinitude property} if, for every finite list of distinct members
$U_1,\ldots,U_s$ and each $j$, the set
\[
 U_j\setminus\bigcup_{\ell\ne j}U_\ell
\]
is infinite. This is a property of the incidence sets only; it says
nothing about the order on the exponent group.
\end{definition}

\mtag{} This is the finite private-index property of
\cref{isc:cp:def:private}, for families of subsets of an arbitrary index
set.

\begin{proof}[Proof of \cref{isc:hj:main:mixed}]
Write $K_j=k((t^{\Gamma_j}))$, $M=K_0K_1$, and
$E=k((t^{\Gamma_0+\Gamma_1}))$.
The mixed exponents $a_i+b_i$ are pairwise distinct. Indeed, equality for
two indices would give a nontrivial relation between the corresponding
$a_i+H$. Their assumed well-ordering therefore makes every $F_U$ a
well-defined Hahn sum.

Suppose there is an algebraic relation among finitely many distinct
$F_{U_1},\ldots,F_{U_s}$. Choose a nonzero polynomial
\[
 P\in M[X_1,\ldots,X_s],\qquad
 P(F_{U_1},\ldots,F_{U_s})=0,
\]
of least possible total degree among nonzero relations on this tuple.
It is not a constant polynomial. In characteristic zero, some partial
$P_j=\partial P/\partial X_j$ is a nonzero polynomial. It has smaller
total degree, so minimality gives
\begin{equation}\label{isc:hj:eq:partialnonzero}
 P_j(F_{U_1},\ldots,F_{U_s})\ne0.
\end{equation}

There are only finitely many coefficients of $P$. Each is a rational
expression in finitely many elements of $K_0$ and $K_1$. Taking the union
of the finite lists of $K_0$-elements used in these expressions, there
are $u_1,\ldots,u_m\in K_0$ such that all coefficients of $P$ lie in
$K_1(u_1,\ldots,u_m)$.

Choose $m+1$ distinct indices in
$U_j\setminus\bigcup_{\ell\ne j}U_\ell$, and call the resulting set $J$.
Apply \cref{isc:hj:lem:annihilator} to obtain
$\partial_{\mathbf c}=\sum_{i\in J}c_i\partial_i$.
This derivation annihilates the coefficients of $P$. By
\eqref{isc:hj:eq:coordinateaction}, it also annihilates each $F_{U_\ell}$
with $\ell\ne j$, whereas
\[
 \partial_{\mathbf c}(F_{U_j})=\sum_{i\in J}c_it^{a_i+b_i}\ne0.
\]
Differentiating the polynomial identity, using the ordinary finite
polynomial chain rule, now yields
\[
 0=\partial_{\mathbf c}\bigl(P(F_{U_1},\ldots,F_{U_s})\bigr)
  =P_j(F_{U_1},\ldots,F_{U_s})\,\partial_{\mathbf c}(F_{U_j}).
\]
Both factors on the right are nonzero by \eqref{isc:hj:eq:partialnonzero}
and \cref{isc:hj:lem:annihilator}, contradicting that $E$ is a field.
Thus no algebraic relation exists.
\end{proof}

\subsection{A finite-generator formulation}

The proof has a useful quantitative form even when the exclusive sets
are not known to be infinite.

\begin{proposition}[Generator-sensitive obstruction]\label[proposition]{isc:hj:prop:quantitative}
Fix $u_1,\ldots,u_m\in K_0$ and finite subsets
$U_1,\ldots,U_s$ or arbitrary summable subsets of $I$. If each exclusive
part $U_j\setminus\bigcup_{\ell\ne j}U_\ell$ has at least $m+1$
elements, then the corresponding $F_{U_j}$ are algebraically independent
over $K_1(u_1,\ldots,u_m)$.
\end{proposition}
\begin{proof}
Choose a relation of minimal total degree over this fixed coefficient
field. Its nonzero partial selects one $j$. The $m+1$ available exclusive
indices give the same annihilator and contradiction as above.
\end{proof}

This statement is consistent with finite mixed sums lying in $M$:
representing enough of them in a field $K_1(u_1,\ldots,u_m)$ requires
enough $K_0$-generators. It does not give a bound on the expression length
of an arbitrary Hahn series, because the choice of generators is itself
unrestricted.

\subsection{Stable variants}

\begin{corollary}[Weights and finite modifications]\label[corollary]{isc:hj:cor:weights}
In \cref{isc:hj:main:mixed}, replace $F_U$ by
\[
 \widetilde F_U=d_U+
 \sum_{i\in U}r_{U,i}t^{a_i+b_i},\qquad d_U\in M,\quad
 r_{U,i}\in k^\times.
\]
Then the resulting family is still algebraically independent over $M$.
Finite changes of the supports or coefficients also preserve the
conclusion, provided the remaining nonzero terms satisfy the same
exclusive infinitude property.
\end{corollary}
\begin{proof}
First take $d_U=0$. On an exclusive set the derivative has the form
$\sum c_i r_{U,i}t^{a_i+b_i}$; its coefficients in the distinct cosets
remain nonzero whenever $c_i\ne0$. The proof is unchanged. Translation
by elements of the base field is an invertible polynomial change of
variables. Finally, a finite sum of mixed monomials belongs to $M$, so
finite modifications can also be absorbed into the translations.
\end{proof}

\begin{corollary}[Algebraic enlargement of the base]\label[corollary]{isc:hj:cor:algbase}
Let $M'$ be an algebraic extension of $M$ contained in $E$. Every family
proved independent over $M$ above remains independent over $M'$.
In particular this applies to the relative algebraic closure of $M$ in
$E$.
\end{corollary}
\begin{proof}
For a finite tuple $\mathbf F$, the extension
$M'(\mathbf F)/M(\mathbf F)$ is algebraic. The transcendence-degree
formula, or the elementary exchange lemma for algebraic dependence,
gives
$\trdeg(M'(\mathbf F)/M')=\trdeg(M(\mathbf F)/M)$.
Thus the degree equals the length of every finite tuple of distinct
members, as required.
\end{proof}

When $E$ is real closed, the relative algebraic closure of $M$ in $E$
is real closed: positive square roots and roots of odd-degree polynomials
in $E$ are algebraic over $M$. When $E$ is algebraically closed, the
relative algebraic closure is algebraically closed for the same reason.
Consequently passing from the compositum to its appropriate algebraic
closure does not eliminate the transcendental gap.
\mtag{} \Cref{isc:cp:rem:algebraicbase} is the same observation for the
Euler jets of source~02.

\section{Families of supports of the required size}\label{isc:hj:sec:combinatorics}

The differential theorem separates the algebra from the combinatorics.
This section supplies the support families. The constructions themselves
are elementary standard independent-family and almost-disjoint-family
arguments; no novelty is claimed for them.

\subsection{Every infinite cardinal}

\begin{lemma}[A maximal independent incidence family]\label[lemma]{isc:hj:lem:independentfamily}
For every infinite cardinal $\kappa$ there is a set $\mathcal I_\kappa$ of
cardinality $\kappa$ and a family $(U_X)_{X\subseteq\kappa}$ of $2^\kappa$
subsets of $\mathcal I_\kappa$ such that every finite Boolean pattern on
distinct members is realized at $\kappa$ indices. In particular the family
has finite-exclusive infinitude.
\end{lemma}
\begin{proof}
Take
\[
 \mathcal I_\kappa=\{(s,T,\xi):s\in[\kappa]^{<\omega},\
                         T\subseteq\mathcal P(s),\ \xi<\kappa\}.
\]
There are $\kappa$ finite subsets $s$, only finitely many choices of $T$
for each fixed finite $s$, and $\kappa$ choices of $\xi$.
Thus $|\mathcal I_\kappa|=\kappa$. For $X\subseteq\kappa$ define
\[
 U_X=\{(s,T,\xi):X\cap s\in T\}.
\]
Given distinct $X_1,\ldots,X_r$, choose, for each pair, an element of
their symmetric difference. The finite set $s$ of the chosen elements
makes all traces $X_j\cap s$ distinct. For any prescribed pattern
$\varepsilon_j\in\{0,1\}$, choose $T$ to contain exactly those of these
traces for which $\varepsilon_j=1$. Then every $(s,T,\xi)$, $\xi<\kappa$,
realizes the desired pattern. This gives at least $\kappa$ witnesses,
and there cannot be more because $|\mathcal I_\kappa|=\kappa$.
Taking a pattern that distinguishes a pair also shows that all $U_X$
are different.
\end{proof}

Fixing any bijection between $\mathcal I_\kappa$ and $\kappa$ transfers the
family to subsets of an ordinal-indexed support. The exponent groups in the
next section are explicit without this choice; the formula for the
witnesses is explicit relative to the chosen bijection.

\subsection{A concrete countable support}

For a finite binary word $s=(s_0,\ldots,s_{m-1})$, set
\[
 \operatorname{code}(s)=2^m-1+
            \sum_{j=0}^{m-1}s_j2^{m-1-j}.
\]
Words of length $m$ have codes in $[2^m-1,2^{m+1}-2]$, and words of
different lengths have disjoint code intervals. For an infinite binary
sequence $\eta\in2^{\N}$ let
\begin{equation}\label{isc:hj:eq:branchsets}
 \mathfrak B_\eta=\{\operatorname{code}(\eta\vert_m):m\geq1\}.
\end{equation}
These are infinite subsets of $\N$.

\mtag{} With $1$-based indices this is $N(s)-1$ for the coding
\lbl{tail:eq:binarycode} of the tail-spans report, and
$\mathfrak B_\eta$ is its set \lbl{tail:eq:Aalpha} shifted down by one; the
same kind of family serves \cref{isc:cp:cor:continuum}.

\begin{lemma}[Binary branches give exclusive tails]\label[lemma]{isc:hj:lem:branches}
Distinct $\mathfrak B_\eta,\mathfrak B_\theta$ have finite intersection.
Hence the continuum sized family $(\mathfrak B_\eta)_{\eta\in2^{\N}}$ has
finite-exclusive infinitude.
\end{lemma}
\begin{proof}
Once two branches first differ, their prefixes of every larger length
differ, and therefore so do their codes. They can share only the
finitely many earlier prefixes. For finitely many other branches, remove
the finite union of these finite intersections from a fixed
$\mathfrak B_\eta$; an infinite tail remains.
\end{proof}

For example, the constantly zero and constantly one branches give
\[
 \mathfrak B_{000\cdots}=\{1,3,7,15,\ldots\},\qquad
 \mathfrak B_{111\cdots}=\{2,6,14,30,\ldots\}.
\]
This fully specified family is used for the prime witnesses.

\section{Maximal transcendence at every set cardinal}\label{isc:hj:sec:cardinal}

\subsection{Two explicit independent exponent groups}

It is important to distinguish ordinal operations used as indices from
surreal arithmetic used inside the field. Define ordinal indices
\[
 \iota_0(\alpha)=2\cdot_{\Ord}\alpha,\qquad
 \iota_1(\alpha)=(2\cdot_{\Ord}\alpha)+_{\Ord}1.
\]
Here $\cdot_{\Ord}$ and $+_{\Ord}$ are ordinary ordinal multiplication
and addition, not Hessenberg operations. They satisfy
\[
 \iota_0(\alpha)<\iota_1(\alpha)<\iota_0(\beta)
 \qquad(\alpha<\beta).
\]
Regard these ordinals as surreal numbers and put
\begin{equation}\label{isc:hj:eq:innerexponents}
 e_\alpha=\omega^{-\iota_0(\alpha)},\qquad
 f_\alpha=\omega^{-\iota_1(\alpha)},\qquad
 r_\alpha=e_\alpha+f_\alpha.
\end{equation}
Thus $e_0=1$, $f_0=\omega^{-1}$,
$e_1=\omega^{-2}$, $f_1=\omega^{-3}$, and so on.
The symbols $e_\alpha,f_\alpha$ are themselves surreal numbers that will
be used as exponents in a second, outer Conway normal form.

\begin{lemma}\label[lemma]{isc:hj:lem:groups}
For any infinite cardinal $\kappa$ define
\[
 A_\kappa=\Span\{e_\alpha:\alpha<\kappa\},\qquad
 B_\kappa=\Span\{f_\alpha:\alpha<\kappa\}.
\]
Then $A_\kappa\cap B_\kappa=0$, both groups have cardinality $\kappa$,
and $(r_\alpha)_{\alpha<\kappa}$ is strictly decreasing and positive.
\end{lemma}
\begin{proof}
All the $e_\alpha$ and $f_\alpha$ are distinct Conway monomials.
Uniqueness of finite Conway normal form implies that they are linearly
independent over $\R$, hence over $\Q$. This proves the intersection
claim. A finite rational combination of $\kappa$ generators has only
$\kappa$ possibilities, and there are $\kappa$ distinct generators, so
the cardinalities are exactly $\kappa$.

If $\alpha<\beta$, the largest Conway exponent occurring in
$r_\alpha-r_\beta$ is $-\iota_0(\alpha)$, with coefficient $1$.
Therefore the difference is positive. Each $r_\alpha$ is also positive.
\end{proof}

\subsection{Building the independent omnific integers}

Transfer the family of \cref{isc:hj:lem:independentfamily} to a family
$\mathcal U_\kappa\subseteq\mathcal P(\kappa)$ by a fixed bijection. Define
\begin{equation}\label{isc:hj:eq:cardwitness}
 \Xi_U=\sum_{\alpha\in U}\omega^{r_\alpha}
 \qquad(U\in\mathcal U_\kappa).
\end{equation}
The decreasing support ensures that this is a Conway normal form. Its
exponents are positive and its constant coefficient is zero, so it is an
omnific integer. There are no analytic convergence conditions to check,
even when $\kappa$ is uncountable.

\begin{proof}[Proof of \cref{isc:hj:main:cardinal}]
In \cref{isc:hj:main:mixed}, use
\[
 \Gamma_0=A_\kappa,\quad\Gamma_1=B_\kappa,\quad
 a_\alpha=-e_\alpha,\quad b_\alpha=-f_\alpha.
\]
The classes of the $a_\alpha$ are independent modulo the zero
intersection. The $b_\alpha$ are distinct. The mixed exponents
$-r_\alpha$ are strictly increasing, with order type $\kappa$, so are
well ordered. Since $t^{-r_\alpha}=\omega^{r_\alpha}$, the theorem
proves that the $2^\kappa$ elements in \eqref{isc:hj:eq:cardwitness} are
algebraically independent over the ordinary compositum, for $k=\R$ and
$k=\C$.

Let $\Gamma=A_\kappa+B_\kappa$. Its cardinality is $\kappa$.
A Hahn series is determined by a function from $\Gamma$ to $k$, so
\[
 |\K{k}{\Gamma}|
 \leq |k|^\kappa
 =(2^{\aleph_0})^\kappa=2^\kappa.
\]
The independent family supplies the opposite bound and proves that the
transcendence degree is exactly $2^\kappa$.
The factors are linearly disjoint by \cref{isc:hj:prop:ld}; their real or
algebraic closedness follows from divisibility and the classical Hahn
closedness theorem. The witnesses have real coefficients, so exactly the
same family works inside the complexified fields.
\end{proof}

Each factor field itself has cardinality $2^\kappa$: arbitrary subsets of
the decreasing support $(e_\alpha)_{\alpha<\kappa}$, or of the analogous
$f$-support, give $2^\kappa$ different series, and the preceding upper
bound applies to each factor. The ordinary compositum also has cardinality
$2^\kappa$, since it contains a factor and is contained in the join.

\begin{corollary}\label[corollary]{isc:hj:cor:realclosuregap}
The relative real algebraic closure of
$\K{\R}{A_\kappa}\K{\R}{B_\kappa}$ in the real Hahn join remains
of transcendence codimension $2^\kappa$. The relative algebraic closure
of the complex compositum has the same property in the complex Hahn join.
\end{corollary}
\begin{proof}
The independent family survives by \cref{isc:hj:cor:algbase}; the upper
bound is unchanged.
\end{proof}

\subsection{A fraction-field consequence that does not assume descent}
\label{isc:hj:sec:arithmetic}

Let $\A{k}{A_\kappa}$, $\A{k}{B_\kappa}$ and $\A{k}{A_\kappa+B_\kappa}$ be
as defined in \cref{isc:hj:sec:foundations}. Denote by
$\A{k}{A_\kappa}[\A{k}{B_\kappa}]$ the subring of the join generated by
$\A{k}{A_\kappa}$ and $\A{k}{B_\kappa}$.

\begin{corollary}[Maximal arithmetic-join gap]\label[corollary]{isc:hj:cor:arithmetic}
For $k=\R,\C$ one has
\[
 \trdeg\bigl(\Frac\A{k}{A_\kappa+B_\kappa}\big/
 \Frac\bigl(\A{k}{A_\kappa}[\A{k}{B_\kappa}]\bigr)\bigr)=2^\kappa.
\]
The family $(\Xi_U)_{U\in\mathcal U_\kappa}$ witnesses this equality.
\end{corollary}
\begin{proof}
The ring $\A{k}{A_\kappa+B_\kappa}$ contains both factor rings and every
$\Xi_U$. Its fraction field is a subfield of the Hahn join. Also,
\[
 \Frac\bigl(\A{k}{A_\kappa}[\A{k}{B_\kappa}]\bigr)\subseteq
 \K{k}{A_\kappa}\K{k}{B_\kappa}.
\]
Independence over the larger field implies independence over
$\Frac(\A{k}{A_\kappa}[\A{k}{B_\kappa}])$. This gives the lower bound
$2^\kappa$. The ambient Hahn join has cardinality $2^\kappa$, which
supplies the upper bound.
\end{proof}

This is a failure of reconstructing the arithmetic Hahn join by finite
ring operations and fractions from its two pieces. It is distinct from
this report's failure of fraction-field descent for \emph{intersections}
(\cref{isc:main:fractions}). Neither argument authorizes replacing
fixed-support omnific fraction fields by the corresponding full Hahn fields
without proof.
\mtag{} It is the set-sized, non-separated counterpart of
\cref{isc:cp:thm:arithmetic}, which obtains nonintegrality of the omnific
join over the ring generated by two explicit separated copies, after any
set adjunction.

\section{A fully Archimedean, interleaved example}\label{isc:hj:sec:archimedean}

The nested monomial example has many distinct growth scales. We now
remove that feature. All exponent groups in this section are subgroups
of the ordinary real numbers.

\subsection{Independent real exponents approaching zero}

We use the elementary multiquadratic fact that the square roots of
distinct positive primes are linearly independent over $\Q$. One way
to see it is to work in the finite multiquadratic extension generated by
any finite list of those roots. Unique prime factorization makes their
square classes independent in $\Q^\times/(\Q^\times)^2$, so adjoining
the roots gives independent sign-changing automorphisms. Applying the
automorphism that changes just one root's sign to a rational linear
relation, and subtracting, forces that root's coefficient to vanish.
Doing this for every root proves the assertion. This is ordinary
finite Galois theory, not an assumption about infinite algebraic
extensions.

Let $p_0=2,p_1=3,p_2=5,\ldots$ and define
\begin{equation}\label{isc:hj:eq:realexponents}
 \lambda_n=\frac{\sqrt{p_{2n}}}{p_{2n}^{2}},\qquad
 \mu_n=\frac{\sqrt{p_{2n+1}}}{p_{2n+1}^{2}},\qquad
 \rho_n=\lambda_n+\mu_n.
\end{equation}
Since $p^{-3/2}=\sqrt p/p^2$, these are exactly the quantities in
\cref{isc:hj:main:prime}. Set
\[
 A=\Span\{\lambda_n:n\geq0\},\qquad
 B=\Span\{\mu_n:n\geq0\},\qquad \Gamma=A+B.
\]

\begin{lemma}\label[lemma]{isc:hj:lem:realgroups}
The combined family of all $\lambda_n,\mu_n$ is $\Q$-linearly
independent. The groups $A,B,\Gamma$ are countable and divisible,
$A\cap B=0$, and $\rho_n$ strictly decreases to zero in $\R$.
The group $\Gamma$ is Archimedean and has no nonzero proper convex subgroup.
\end{lemma}
\begin{proof}
Each number in the combined family is a nonzero rational multiple of
the square root of a different prime, so the preceding independence
fact applies. The claims about cardinalities, divisibility and
intersection follow.
As the primes strictly increase without bound, both $(\lambda_n)$ and
$(\mu_n)$ strictly decrease to zero; their sum does too.
Every ordered subgroup of $\R$ is Archimedean. If a convex subgroup
contains $g>0$, then for any $h>0$ in $\Gamma$ there is a natural number
$m$ with $h<mg$. Convexity then puts $h$ in that subgroup. Hence a
nonzero convex subgroup is all of $\Gamma$.
\end{proof}

Both $A$ and $B$ are dense in $\R$: each contains the rational multiples
of a nonzero real number. Thus this is not a construction in which all
nonzero exponents from one factor occupy a convex scale separated from
the other. In particular, a proof depending on a nontrivial convex
quotient of $\Gamma$ cannot be used here.
\mtag{} In particular the iterated presentation of \cref{isc:lem:iterated},
and with it the rank test \cref{isc:cp:cor:surrealrank}, is not available for
this pair.

\subsection{Continuum many mixed sums}

For the branch sets from \eqref{isc:hj:eq:branchsets}, define
\[
 F_\eta=\sum_{n\in\mathfrak B_\eta}\omega^{\rho_n},\qquad P_\eta=1+F_\eta.
\]
Every $\mathfrak B_\eta$ is an unbounded subset of $\N$, so its increasing
enumeration selects a subsequence of $(\rho_n)$ that still decreases to
zero. These sums therefore have reverse-well-ordered support of type
$\omega$; adjoining the constant term gives reverse support order type
$\omega+1$.

\begin{proposition}\label[proposition]{isc:hj:prop:archindependence}
The family $(P_\eta)_{\eta\in2^{\N}}$ is algebraically independent over
$\K{k}{A}\K{k}{B}$, for both $k=\R$ and $k=\C$. In either case
\[
 \trdeg\bigl(\K{k}{\Gamma}/\K{k}{A}\K{k}{B}\bigr)=\cc.
\]
\end{proposition}
\begin{proof}
Use $a_n=-\lambda_n$, $b_n=-\mu_n$ in \cref{isc:hj:main:mixed}.
The mixed $t$-exponents $-\rho_n$ strictly increase to zero, hence form
a well-ordered set of order type $\omega$. Apply \cref{isc:hj:lem:branches}
and then \cref{isc:hj:cor:weights} to add the constant terms.
The lower bound is $|2^{\N}|=\cc$. For the upper bound, $\Gamma$
is countable and $|k|=\cc$, so the entire Hahn join has cardinality
at most $\cc^{\aleph_0}=\cc$.
\end{proof}

For illustration, two members begin as
\[
 \begin{aligned}
 P_{000\cdots}&=1+\omega^{\rho_1}+\omega^{\rho_3}
                         +\omega^{\rho_7}+\omega^{\rho_{15}}+\cdots,\\
 P_{111\cdots}&=1+\omega^{\rho_2}+\omega^{\rho_6}
                         +\omega^{\rho_{14}}+\omega^{\rho_{30}}+\cdots.
 \end{aligned}
\]
Each displayed monomial already lies in the ordinary compositum because
\[
 \omega^{\rho_n}=\omega^{\lambda_n}\omega^{\mu_n}.
\]
Nevertheless the full family of Hahn sums is algebraically independent
over that compositum. This is the sense in which the obstruction is a
mixed-summation obstruction rather than a missing-monomial obstruction.

\section{Prime omnific and Gaussian omnific witnesses}\label{isc:hj:sec:prime}

\subsection{The isolated external theorem}

The following is the only specialized factorization theorem imported
into the source. It is the one-row result of Pitteloud together
with the transfer of Biljakovic--Kochetov--Kuhlmann, as stated in
\cite[\S1.3]{LM}; the latter paper also gives broader primality and
transfer criteria in Theorem~F and Proposition~8.2.9.

\begin{theorem}[Classical one-row prime theorem]\label{isc:hj:thm:classicalprime}
Let $k$ be any characteristic-zero field and let $G$ be a set-sized
ordered abelian group containing the ordered additive group $\R$.
If $s_0<s_1<\cdots<0$ are real numbers with supremum zero and all
$c_n\in k$ are nonzero, then
\[
 q=1+\sum_{n\geq0}c_nt^{s_n}
\]
is prime in the ring of Hahn series with coefficients in $k$ and
exponents in $G^{\leq0}$.
\end{theorem}

No proof of \cref{isc:hj:thm:classicalprime} is claimed here; the references
\cite{Pitteloud,BKK,LM} supply this classical input. In particular,
the independence argument does not reprove or replace the factorization
theory behind it. Everything needed to pass from this input to the
specific omnific and Gaussian statements is proved next.
\mtag{} Checked in the merge against the arXiv text of \cite{LM}
(version~5, Section~1.3): Pitteloud proved that $1+b$, for $b$ a sum over
an $\omega$-sequence of real exponents increasing to $0$, is prime in
$K((\R^{\le0}))$ for a field $K$ of characteristic zero, and
Biljakovic--Kochetov--Kuhlmann observed that it remains prime in
$K((G^{\le0}))$ for every $G\supseteq\R$. The statement above agrees with that
account; the original papers were not read in the merge.

\subsection{Pulling primality through the constant term}

\begin{lemma}[Constant-term-one pullback]\label[lemma]{isc:hj:lem:pullback}
Let $k$ be a field, $D\subseteq k$ a unital subring, and $G$ an ordered
abelian group. In the $\omega$ convention consider the rings
\[
 \Ad{k}{k}{G}=k\oplus\PiG{k}{G},\qquad \Ad{D}{k}{G}=D\oplus\PiG{k}{G}.
\]
Suppose $q\in\Ad{D}{k}{G}$ is prime in $\Ad{k}{k}{G}$ and $\ct(q)=1$. Then
$q$ is prime in $\Ad{D}{k}{G}$. More precisely, divisibility by $q$ in
$\Ad{k}{k}{G}$ of an element of $\Ad{D}{k}{G}$ is already divisibility in
$\Ad{D}{k}{G}$.
\end{lemma}
\begin{proof}
If $x=qu$ with $x\in\Ad{D}{k}{G}$, $u\in\Ad{k}{k}{G}$, apply the
constant-term homomorphism. It gives
$\ct(u)=\ct(x)/\ct(q)=\ct(x)\in D$, so $u\in\Ad{D}{k}{G}$.
Now suppose $q$ divides $xy$ in $\Ad{D}{k}{G}$. It divides the same product
in $\Ad{k}{k}{G}$, so it divides $x$ or $y$ in $\Ad{k}{k}{G}$. The preceding
observation descends that divisibility to $\Ad{D}{k}{G}$. Also $q$ cannot
be a unit in $\Ad{D}{k}{G}$, since that would make it a unit in
$\Ad{k}{k}{G}$. It is nonzero, as required for a prime element.
\end{proof}

The constant-term condition is not cosmetic. Without it a quotient
that exists after enlarging the constant ring need not lie in the
smaller ring. The lemma uses exactly the feature supplied by the
constant $1$ in every $P_\eta$.
\mtag{} For $D=D_k$ one has $\Ad{D_k}{k}{G}=\A{k}{G}$; the ring
$\Ad{\R}{\R}{G}$ is the omnific-preserving report's
$\mathcal A_\Gamma=\R\oplus\Pi_\Gamma$ (\cref{isc:sec:conventions}).

\subsection{Localization handles the full proper-class rings}

\begin{proposition}\label[proposition]{isc:hj:prop:fullprime}
Let $r_0>r_1>\cdots>0$ be ordinary real numbers decreasing to zero.
Then
\[
 q=1+\sum_{n\geq0}c_n\omega^{r_n}
\]
is prime in $\Oz$ when $c_n\in\R^\times$, and is prime in
$\Oz[i]$ when $c_n\in\C^\times$.
\end{proposition}
\begin{proof}
We give the complex case, with $D=\Z[i]$; the real case is identical
with $k=\R$ and $D=\Z$. Suppose $xy=qz$ in $\Oz[i]$.
The union of the Conway supports of $x,y,z,q$ is a set. Let $G$ be its
$\Q$-linear span together with the whole set $\R$ inside $\No$.
This is a set-sized ordered divisible group containing $\R$.
All four elements belong to $\Ad{D}{k}{G}$.

Under $t^g=\omega^{-g}$, the series $q$ has real exponents
$-r_n$ increasing to zero. By \cref{isc:hj:thm:classicalprime}, it is prime
in $\Ad{k}{k}{G}$. By \cref{isc:hj:lem:pullback} it is prime in
$\Ad{D}{k}{G}$. Thus $q$ divides $x$ or $y$ in that set-sized ring, and
hence in the full Gaussian omnific ring. The series is nonzero and is not a
unit: its largest Conway exponent is positive, and the largest exponent of
a product of two nonzero elements of $\Ad{D}{k}{\No}$ is the sum of their
nonnegative largest exponents. Such a product cannot be $1$ if one
factor has positive largest exponent. This proves primality in the
full class ring without applying a set theorem directly to a proper class.
\end{proof}

\begin{proof}[Completion of the proof of \cref{isc:hj:main:prime}]
The exponent and field assertions, and algebraic independence, are
\cref{isc:hj:lem:realgroups,isc:hj:prop:archindependence}.
For each branch $\eta$, enumerate $\mathfrak B_\eta$ increasingly.
The resulting subsequence of $\rho_n$ is strictly decreasing to zero.
Apply \cref{isc:hj:prop:fullprime} with every coefficient equal to $1$.
This gives both primality assertions simultaneously.
\end{proof}

\begin{corollary}[Pairwise nonassociate primes]\label[corollary]{isc:hj:cor:nonassociate}
The $P_\eta$ form continuum many pairwise nonassociate primes in
$\Oz$, and also in $\Oz[i]$, while remaining algebraically independent
over the respective full Hahn composita.
\end{corollary}
\begin{proof}
The only units of $\Oz$ are $\pm1$, and those of $\Oz[i]$ are
$\pm1,\pm i$. Indeed, the leading-exponent argument in
\cref{isc:hj:prop:fullprime} forces any unit and its inverse to have only
constant terms, and the constants must be units of $\Z$ or $\Z[i]$.
Since every $P_\eta$ has constant term $1$, two associates would have
to be equal. Algebraic independence implies that distinct members are
not equal.
\end{proof}

The prime conclusion is arithmetic, whereas the independence conclusion
is relative to a much larger field than $\Q$ or $\Q(i)$.
One should not confuse the two: primality alone supplies no such
relative transcendence statement. Conversely the mixed-support proof
works for many series with constant term zero, which are not prime in
$\Oz$ because they are divisible by every nonzero ordinary integer.

\section{A fixed pair of class-sized Hahn subfields}\label{isc:hj:sec:classes}

\subsection{Explicit coordinates, not a class Hamel basis}

Use the exponents in \eqref{isc:hj:eq:innerexponents} for all ordinals and
define
\[
 A_*=\left\{\sum_{\alpha\in F}q_\alpha e_\alpha:
                  F\subseteq\On\text{ finite},\ q_\alpha\in\Q\right\},
 \qquad
 B_*=\left\{\sum_{\alpha\in F}q_\alpha f_\alpha:
                  F\subseteq\On\text{ finite},\ q_\alpha\in\Q\right\}.
\]
These are proper-class divisible subgroups of $\No$ with zero
intersection. Each exponent in $A_*+B_*$ has a unique \emph{finite}
coordinate expression in the $e_\alpha,f_\alpha$. Consequently its
coordinate functions are explicitly defined class maps.

Let
\[
 K_{0,*}=\K{k}{A_*},\qquad K_{1,*}=\K{k}{B_*},\qquad
 M_*=K_{0,*}K_{1,*},\qquad E_*=\K{k}{A_*+B_*}.
\]
The notation allows only set-sized Hahn supports. All field operations
are well defined by localization: finitely many input series involve a
set of exponent coordinates, and the relevant operations take place in
a set-sized divisible-group Hahn field. The same observation proves real
closedness for $k=\R$ and algebraic closedness for $k=\C$.
The intersection and linear-disjointness proof in \cref{isc:hj:prop:ld} is
coefficientwise and applies without change to every finite relation.

For each ordinal $\alpha$, the map extracting the negative of the
$e_\alpha$ coordinate of a $t$-exponent is an additive character,
vanishes on $B_*$, and takes value $1$ at $-e_\alpha$.
It therefore gives exactly the coordinate derivation used in the
set-sized proof. No extension of a class-independent family to a class
basis is required.

\begin{lemma}[The mixed-support proof is local]\label[lemma]{isc:hj:lem:classlocal}
The conclusion of \cref{isc:hj:main:mixed} applies to the class fields above
for any set of the coordinate indices, with coefficients of a putative
relation allowed anywhere in $M_*$. It also proves independence for a
class of witnesses whenever each finite subfamily satisfies the
exclusive infinitude condition.
\end{lemma}
\begin{proof}
A polynomial relation uses finitely many witnesses and finitely many
coefficients. Each coefficient belongs to
$K_{1,*}(u_1,\ldots,u_m)$ for finitely many $u_q\in K_{0,*}$.
Choose $m+1$ exclusive indices. The required characters are the explicit
coordinate functions just described. A finite homogeneous system over
the class field $K_{0,*}$ has a nonzero solution by the ordinary finite
Gaussian-elimination proof; it involves only finitely many field
operations. Coset independence is still valid, since the supports of
each of the finitely many summands are sets in disjoint cosets.
Every subsequent step is the finite polynomial calculation in
\cref{isc:hj:sec:mixed}.
\end{proof}

A more general class group can be handled by set localization when
suitable local characters exist. The explicit construction above is
preferable here because it displays them and avoids additional
foundational hypotheses about arbitrary class vector spaces.

\subsection{An ordinal-indexed family of witnesses}

For an ordinal $\alpha$, let
\[
 \Lambda_\alpha=\{(\omega\cdot_{\Ord}\alpha)+_{\Ord}n:n<\omega\}.
\]
These countably infinite ordinal blocks are pairwise disjoint. Define
\begin{equation}\label{isc:hj:eq:classwitness}
 W_\alpha=\sum_{n<\omega}
   \omega^{r_{(\omega\cdot_{\Ord}\alpha)+_{\Ord}n}}.
\end{equation}
Each $W_\alpha$ has countable, strictly decreasing positive support, so
is an omnific integer belonging to $E_*$.
\mtag{} The blocks are source~02's labels
$\beta_{\alpha,n}=\omega\cdot\alpha+n$ of \eqref{isc:cp:eq:labels},
here with $n\ge0$ instead of $n\ge1$.

\begin{proof}[Proof of \cref{isc:hj:main:class}]
The supports $\Lambda_\alpha$ are disjoint and infinite. Apply
\cref{isc:hj:lem:classlocal} to any finite selection of the $W_\alpha$.
This proves algebraic independence of the ordinal-indexed family over
$M_*$, for real or complex coefficients.

In addition, fix an infinite cardinal $\kappa$ and use the family
$\mathcal U_\kappa$ from \cref{isc:hj:sec:cardinal} on the first $\kappa$
indices. The same local proof shows that its $2^\kappa$ witnesses remain
independent over the \emph{full} class compositum $M_*$, not merely
over the small compositum from that section. The coefficients of a
relation were never restricted to the first $\kappa$ coordinates.

Finally, a set-sized transcendence basis of cardinality $\lambda$ would
bound every algebraically independent subset over $M_*$ by $\lambda$
(for a finite basis, by its finite size). This is the ordinary exchange
argument for algebraic dependence; it uses only finite polynomial
relations and remains valid over a class base field. The arbitrary
cardinal lower bounds contradict such a bound. Hence no set-sized
transcendence basis exists.
\end{proof}

The theorem is deliberately not written as
``$\trdeg(E_*/M_*)=\On$'': $\On$ is not a cardinal in the ordinary
set-sized sense of transcendence degree. The precise assertions are an
explicit independent class and independent sets of arbitrarily large
cardinality. Also, no order-group isomorphism $A_*\cong\No$ or
$B_*\cong\No$ has been proved or used.
\mtag{} In the last proof $\lambda$ is a cardinal, not the Euler weight.

\begin{remark}[A second route to the transcendence of $S$]\label[remark]{isc:hj:rem:secondroute}
\mtag{} The class-local proof also applies to the explicit copies of
\cref{isc:sec:explicit}, and gives a second proof that the witness
\eqref{isc:eq:explicitS} is transcendental over $M=K_0K_1$ in
$E=\K{k}{G_0+G_1}$, for $k=\R$ and $\C$. In the $t$-convention,
$S=\sum_{n\ge1}t^{a_n+b_n}$ with
\[
 a_n=-g_n=-\omega^{\psi(n)}\in G_0,\qquad b_n=ng\in G_1,
 \qquad g=\omega^{7/2}.
\]
Here $G_0\cap G_1=0$; the $a_n$ are distinct inner monomials, hence
$\Q$-linearly independent; the $b_n$ are distinct; and the $t$-exponents
increase, since consecutive differences are $g+g_n-g_{n+1}>0$ by
\eqref{isc:eq:separated}. The singleton family $\{\N_{>0}\}$ has
finite-exclusive infinitude. The characters are explicit, as in
\cref{isc:hj:lem:classlocal}: for $\gamma=\gamma_0+\gamma_1$ with
$\gamma_j\in G_j$, let $\chi_n(\gamma)$ be minus the coefficient of the inner
monomial $\omega^{\psi(n)}$ in $\gamma_0$. It is additive with values in
$\R\subseteq k$ (\cref{isc:hj:lem:euler} needs no more), vanishes on $G_1$,
and $\chi_n(a_m)=\delta_{nm}$. The proof of \cref{isc:hj:lem:classlocal},
with the class fields $K_0,K_1$ in place of $K_{0,*},K_{1,*}$, then gives
transcendence of $S$ over $M$. This route uses neither the iterated
coordinates of \cref{isc:lem:iterated} nor the coefficient-field bound of
\cref{isc:lem:finitecoeff}; unlike \cref{isc:thm:joinobstruction} it needs
characteristic zero and monomial coefficients $s_n=\omega^{g_n}$, on which
$\partial_{\chi_n}$ acts as $s_n\,\partial/\partial s_n$. It does not give
the differential transcendence of \cref{isc:cp:thm:omnific}. The
observation is the merge's; the source does not apply its theorem to $S$.
\end{remark}

\section{Ordinary maps versus preservation of Hahn sums}\label{isc:hj:sec:strong}

The maximal gap has a further structural consequence. This section
uses only the set-sized fields from \cref{isc:hj:main:cardinal} or
\cref{isc:hj:main:prime}. Put $E=\K{k}{A+B}$ and
$M=\K{k}{A}\K{k}{B}$ as before.

\subsection{The strong relative structure is rigid}

\begin{definition}\label[definition]{isc:hj:def:strong}
A field homomorphism on $E$ is \emph{strong} if it sends every Hahn
summable family to a Hahn summable family and commutes with its sum.
A derivation is strong under the analogous requirement. These are
requirements on the specified Hahn presentation, not on an unspecified
topology.
\end{definition}

\begin{proposition}[Strong relative rigidity]\label[proposition]{isc:hj:prop:strongrigidity}
Every strong field endomorphism of $E$ fixing $M$ pointwise is the
identity. Every strong derivation of $E$ annihilating $M$ is zero.
\end{proposition}
\begin{proof}
For $g=a+b\in A+B$, one has
$\omega^g=\omega^a\omega^b\in M$, and $k\subseteq M$.
Every $x\in E$ is the Hahn sum of its monomial terms
$c_g\omega^g$, each of which belongs to $M$.
A strong endomorphism fixing those terms fixes their sum. A strong
derivation annihilating those terms annihilates their sum.
\end{proof}

Notice that this is not automatic strongness of arbitrary field maps.
It is rigidity \emph{conditional on} preserving all the relevant Hahn
sums. The next results show how large the distinction can be.

\subsection{Ordinary relative derivations are abundant}

\begin{lemma}[Extending assignments on a transcendence basis]\label[lemma]{isc:hj:lem:derivebasis}
Let $E/M$ be a characteristic-zero extension and let $\mathcal T$ be a
transcendence basis. Every function $\varphi:\mathcal T\to E$ extends
uniquely to a derivation $\delta:E\to E$ annihilating $M$.
\end{lemma}
\begin{proof}
Algebraic independence of $\mathcal T$ gives a derivation on the polynomial
ring $M[\mathcal T]$ by the finite polynomial rule $\delta(y)=\varphi(y)$,
$y\in\mathcal T$; the quotient rule extends it uniquely to $M(\mathcal T)$.
Since $E/M(\mathcal T)$ is algebraic and characteristic zero, every minimal
polynomial is separable. For an algebraic element $z$ with minimal
polynomial $p(X)=\sum_ja_jX^j$, any extension must satisfy
\[
 \delta(z)=-\frac{\sum_j\delta(a_j)z^j}{p'(z)}.
\]
The denominator is nonzero. The standard quotient-ring calculation
shows this formula defines a derivation on each simple algebraic
extension. Iterating over finite algebraic extensions gives compatible
maps by uniqueness, and their union defines $\delta$ on $E$.
\end{proof}

\begin{theorem}[Maximal ordinary differential freedom]\label{isc:hj:thm:derivations}
For the fields in \cref{isc:hj:main:cardinal},
\[
 |\operatorname{Der}_M(E,E)|=2^{\,2^\kappa},
 \qquad
 \operatorname{Der}^{\mathrm{strong}}_M(E,E)=\{0\}.
\]
The vector space of relative K\"ahler differentials $\Omega_{E/M}$ has
an $E$-basis of cardinality $2^\kappa$. For the Archimedean example,
these quantities become $2^{\cc}$ and $\cc$, respectively.
\end{theorem}
\begin{proof}
A transcendence basis has cardinality $2^\kappa$ by
\cref{isc:hj:main:cardinal}. Choosing values $0$ or $1$ independently at its
members and applying \cref{isc:hj:lem:derivebasis} already gives
$2^{2^\kappa}$ distinct derivations. The total number of functions
$E\to E$ is
$(2^\kappa)^{2^\kappa}=2^{2^\kappa}$, so this is exact.
The strong assertion is \cref{isc:hj:prop:strongrigidity}.

For completeness, the symbols $dy$, $y\in\mathcal T$, span $\Omega_{E/M}$:
finite rational differentiation gives this for $M(\mathcal T)$, and
differentiating a separable minimal polynomial expresses each $dz$,
$z\in E$, in that span. They are independent because the derivations that
assign $1$ to one basis member and $0$ to all others extract their
individual coefficients. Hence their cardinality is the dimension.
The Archimedean calculation uses \cref{isc:hj:prop:archindependence}
instead.
\end{proof}

Extend the explicit independent prime family of \cref{isc:hj:main:prime}
to a transcendence basis. The lemma then gives, for each fixed $\eta$, a
relative derivation with $\delta(P_\eta)=1$ and $\delta(P_\theta)=0$ for
$\theta\ne\eta$. It annihilates every individual monomial in the
expansion of $P_\eta$ because those monomials lie in $M$. Thus it cannot
commute with that Hahn sum. This concrete consequence is not a paradox:
ordinary algebraic differentiation has no built-in rule for infinite
formal sums.

The finite annihilators used to prove independence are a different
object. They annihilate only a finitely generated portion of $M$,
chosen after a hypothetical polynomial relation is given. Attempting to
make a nonzero strong derivation annihilate \emph{all} of $M$ at once
would contradict \cref{isc:hj:prop:strongrigidity}.

\subsection{The complex field has many non-strong automorphisms}

\begin{theorem}[Hidden complex-field automorphisms]\label{isc:hj:thm:automorphisms}
For $k=\C$ in \cref{isc:hj:main:cardinal},
\[
 |\operatorname{Aut}(E/M)|=2^{\,2^\kappa},
\]
whereas its strong subgroup is trivial. In the Archimedean example,
$|\operatorname{Aut}(E/M)|=2^{\cc}$.
\end{theorem}
\begin{proof}
Let $M^{\mathrm{alg}}$ be the relative algebraic closure of $M$ in $E$.
It is algebraically closed. Choose a transcendence basis $\mathcal T$ of $E$
over $M^{\mathrm{alg}}$. By \cref{isc:hj:cor:algbase},
$|\mathcal T|=2^\kappa$. Every permutation of $\mathcal T$ defines an
automorphism of $M^{\mathrm{alg}}(\mathcal T)$ fixing $M^{\mathrm{alg}}$.
Such an automorphism extends to its algebraic closure $E$: the usual
extension-by-maximality argument extends embeddings across algebraic
elements by choosing a root of the transported minimal polynomial. The
resulting image is algebraically closed and contains the base field, so it
is all of $E$.

An infinite set of cardinality $\lambda$ has $2^\lambda$ permutations.
For the lower bound, partition it into $\lambda$ disjoint pairs and
independently swap or fix each pair; the upper bound is
$\lambda^\lambda=2^\lambda$. Thus the distinct permutations of $\mathcal T$
produce at least $2^{2^\kappa}$ distinct extensions. There are at most
$|E|^{|E|}=2^{2^\kappa}$ maps on $E$, giving equality.
The strong subgroup is trivial by \cref{isc:hj:prop:strongrigidity}.
The same proof with $|\mathcal T|=\cc$ handles the Archimedean example.
\end{proof}

These automorphisms fix both entire complex Hahn factor fields, hence
all mixed monomials, but nonidentity ones fail to preserve Hahn sums.
They are not claimed to preserve the chosen real part, complex
conjugation, omnific integers, order, or surreal exponentiation.
In particular the theorem is not an assertion about ordered
automorphisms of $\No$, nor about the automorphism group of $\Oz$.
Those stronger compatibility conditions are genuinely additional.

\begin{remark}[Omnific rings force the identity]\label[remark]{isc:hj:rem:omnificaut}
\mtag{} Combined with the omnific-preserving report, the rigidity of
\cref{isc:hj:prop:strongrigidity} excludes one combination of the
compatibilities just listed. Let $E_\R=\K{\R}{A+B}$ and $M_\R$ be the real
join and compositum of \cref{isc:hj:main:cardinal} or
\cref{isc:hj:main:prime}. By \lbl{opa:as:cor:setpairs}, applied with the
nonzero divisible set group $\Gamma=A+B$ and $t^\gamma=\omega^{-\gamma}$,
every field automorphism of $E_\R$ that stabilizes $\A{\R}{A+B}=E_\R\cap\Oz$
is strongly $\R$-linear; if it fixes $M_\R$, it is the identity by
\cref{isc:hj:prop:strongrigidity}. Now let $\sigma\in\operatorname{Aut}(E/M)$
for $k=\C$ commute with complex conjugation and stabilize
$\A{\C}{A+B}=E\cap\Oz[i]$. Then $\sigma$ maps the fixed field $E_\R$ of
conjugation onto itself and stabilizes $E_\R\cap\Oz$, so it is the
identity on $E_\R$; since $i\in M$, $\sigma$ is the identity. So a
nonidentity automorphism of \cref{isc:hj:thm:automorphisms} cannot both
commute with conjugation and preserve the Gaussian omnific ring of $E$; in
the real joins, no nonidentity automorphism fixing $M_\R$ preserves the
omnific ring. The argument depends on \lbl{opa:as:cor:setpairs}, an
unrefereed result of that report, and was not in the source; it answers
Research question~\ref{isc:hj:q:automorphisms} only for this combination
of conditions.
\end{remark}

\section{Scope, comparison, and failure boundaries}\label{isc:hj:sec:scope}

\subsection{What has been answered}

At the pinned repository snapshot, Research question~3 of this report
asks for an exact description of ordinary composita and for transcendence
degrees of full Hahn joins over them. The results here answer its
\emph{degree component} as follows.

\begin{center}
\small
\begin{tabular}{>{\raggedright\arraybackslash}p{0.30\textwidth}>{\raggedright\arraybackslash}p{0.27\textwidth}>{\raggedright\arraybackslash}p{0.34\textwidth}}
\toprule
Setting & Answer & Witnesses and additional structure\\
\midrule
Every infinite set cardinal $\kappa$ & Exact degree $2^\kappa$ &
$2^\kappa$ omnific integers; real and complex versions\\[5pt]
Countable Archimedean exponent group & Exact degree $\cc$ &
Explicit continuum family prime in both $\Oz$ and $\Oz[i]$\\[5pt]
One fixed pair of class Hahn subfields & No set-sized transcendence basis &
Explicit ordinal-indexed independent omnific family\\[5pt]
Relative maps in the complex set examples & $2^{2^\kappa}$ automorphisms &
Only the identity preserves all Hahn sums\\
\bottomrule
\end{tabular}
\end{center}

This report's one-witness separated-scale mechanism is not being
represented as a new result. The new argument works without convex
separation and supplies maximal-size independent families. The class
result uses explicit finite-coordinate groups and does not invoke this
report's independent full copies of $\No$.
\mtag{} Since the pin, source~02 has answered the degree component in the
separated case and the Laurent cardinal cases, also with maximal-size
families (\cref{isc:cp:thm:cardinal}), and has given class-sized families in
the join of two actual copies (\cref{isc:cp:thm:omnific}); the answers above
are those for non-separated examples (\cref{isc:hj:sec:answers,isc:hj:rem:compare}).

The comparison is targeted, not a certification of originality against
every file, pending archive, or paper. The repository catalogue was
consulted to identify neighboring topics, and this report was inspected
for its compositum result and research question.
Published primality was checked separately. The asserted new results
should still receive external mathematical review and a more exhaustive
priority check before any publication claim.

\subsection{What the proof does not imply}\label{isc:hj:sub:notimply}

\paragraph{An arbitrary infinite mixed sum need not be new.}
If $u=t^a\in K_0$ and $v=t^b\in K_1$ with $a+b>0$, then
\[
 \sum_{n\geq0}(uv)^n=\frac1{1-uv}\in M.
\]
Here all $a$-directions are rationally dependent. The example explains
why merely having infinitely many support terms cannot prove
transcendence.

\paragraph{Distinct cosets really matter.}
If $b_i=0$ for all $i$, an admissible sum of $t^{a_i+b_i}$ belongs to
$K_0$. More generally, a common $b$ factors out and leaves a $K_0$-series.
The noncancellation conclusion \eqref{isc:hj:eq:noncancel} would fail without
the distinct-coset hypothesis.

\paragraph{Finite support cannot work over the full compositum.}
A finite mixed sum already belongs to $M$. Finite-exclusive infinitude
supplies arbitrarily many differential directions to overcome the
unknown finite number of coefficient generators. A bound suited to only
one fixed number of generators is not enough for the full compositum.

\paragraph{Not all incidence families give independence.}
For example, if $U_3=U_1\mathbin{\dot\cup}U_2$, then
$F_{U_3}=F_{U_1}+F_{U_2}$. The support hypothesis is therefore a
substantive condition. Pairwise distinctness of supports alone is not
sufficient.

\paragraph{A large independent set need not already be a basis.}
Even when its cardinality equals the transcendence degree, it may be
properly extendible to a transcendence basis of the same cardinality.
The prime family is not asserted to be a basis without extension.

\paragraph{Characteristic zero is used twice.}
It supplies nontrivial coordinate characters into $k$ from a divisible
$\Q$-vector group and ensures a nonzero partial derivative for a
nonconstant polynomial. In characteristic $p$, every additive map from
a divisible group to $(k,+)$ vanishes, and inseparable polynomial
relations can have all partials zero. The present proof does not extend
verbatim.

\paragraph{Order and topology are separate structures.}
No conclusion about order-topological convergence of partial sums, no
continuity classification of the non-strong maps, and no commutation
with conjugation is proved. Hahn summability is the exact infinite-sum
notion used here. Archimedean \emph{exponent groups} do not make their
Hahn \emph{fields} Archimedean.

\paragraph{The full class of surreals is not the base field.}
Every witness belongs to $\No$, so it cannot be transcendental over
$\No$. The base in each theorem is the displayed ordinary compositum.
Likewise a prime in $\Oz[i]$ is a unit in the field $\No(i)$; primality
is always a statement in the ring specified.

\subsection{Logical independence of the ingredients}

The cardinal-degree results do not depend on primality, and neither the
set-sized nor the explicit class independence proof depends on this
report's class-embedding theorems, which the source treats as unreviewed.
The primality corollary requires the published one-row theorem in addition
to the independence proof. The automorphism count requires algebraic
closedness of the complex Hahn join but not preservation of its real
structure. These dependency distinctions are useful both for checking the
source and for reusing only the part needed in another project.

\subsection{What the addition does not claim}\label{isc:hj:sub:nonclaims}

\mtag{} This list collects every limitation stated by source~03, in its
article and delivered README, with its location here, followed by those
added in the merge.
\begin{enumerate}[label=J\arabic*.,leftmargin=2.6em,itemsep=2pt]
\item The arguments are written out but have not been refereed,
independently reviewed or checked by Lean; priority and novelty are not
certified; no named longstanding conjecture is claimed solved (title page
of the source; README).
\item The contribution proposed as new is the mixed-support criterion and
its combined maximality and arithmetic applications, not the classical
ingredients: Hahn fields, the support theorem, the closedness theorems and
Conway normal forms are standard foundations
(\cref{isc:hj:sec:foundations}; title page of the source).
\item Only the transcendence-degree part of Research question~3 is
answered, and only for the examples constructed. No characterization of
all ordinary composita and no membership test is given; the question is not
settled in every case (\cref{isc:hj:sec:intro,isc:hj:sec:scope}; Research
question~\ref{isc:hj:q:characterization}; README).
\item Primality is imported. \Cref{isc:hj:thm:classicalprime} is credited to
Pitteloud, Biljakovic--Kochetov--Kuhlmann and L'Innocente--Mantova; no
proof of it is claimed, no new one-row primality theorem is claimed, and
the independence argument does not reprove or replace the factorization
theory behind it. Its formalization is a separate project and must not be
accepted as checked because the independence kernel is
(\cref{isc:hj:sec:prime,isc:hj:main:prime}; \cref{isc:hj:sub:formal}).
\item Primality of a real element is not inferred to survive
complexification; the Gaussian assertion is proved separately
(\cref{isc:hj:sub:audit}).
\item The class result concerns explicit support subfields
$\K{k}{A_*}$, $\K{k}{B_*}$. It does not identify either with a full copy of
$\No$, uses none of this report's copy constructions, and proves or uses no
order-group isomorphism $A_*\cong\No$ or $B_*\cong\No$. It is not written as
$\trdeg(E_*/M_*)=\On$. No proper-class Hamel basis is selected; algebraic
independence of a class means independence of every finite subfamily
(\cref{isc:hj:main:class,isc:hj:sec:classes}).
\item The derivations are additive-character derivations of a Hahn field.
They are not the Berarducci--Mantova derivation and are not assumed to
respect a surreal exponential (\cref{isc:hj:sec:foundations}).
\item The false blanket identity $\Frac(\A{k}{G})=\K{k}{G}$ is not used.
\Cref{isc:hj:cor:arithmetic} is a failure of reconstruction by finite ring
operations and fractions, distinct from this report's descent failure for
intersections; neither argument authorizes replacing fixed-support omnific
fraction fields by full Hahn fields (\cref{isc:hj:sec:foundations,isc:hj:sec:arithmetic}).
\item The independent-family and almost-disjoint-family constructions are
standard; no novelty is claimed for them (\cref{isc:hj:sec:combinatorics}).
\item \Cref{isc:hj:prop:quantitative} gives no bound on the expression
length of an arbitrary Hahn series (\cref{isc:hj:sec:mixed}).
\item \Cref{isc:hj:prop:strongrigidity} is rigidity conditional on
preserving Hahn sums, not automatic strongness of arbitrary field maps. No
single strong derivation annihilates the full compositum; the finite
annihilators kill only a finitely generated part chosen after a relation is
given (\cref{isc:hj:sec:strong}; \cref{isc:hj:sub:audit}).
\item The automorphisms of \cref{isc:hj:thm:automorphisms} are not claimed
to preserve the chosen real part, complex conjugation, omnific integers,
order or surreal exponentiation. The theorem is not about ordered
automorphisms of $\No$ or about $\operatorname{Aut}(\Oz)$
(\cref{isc:hj:sec:strong}).
\item This report's one-witness separated-scale mechanism is not
represented as new. The repository comparison was targeted, not
exhaustive; the catalogue was consulted for neighbouring topics only.
External review and a more exhaustive priority check are needed before any
publication claim (\cref{isc:hj:sec:scope}; README).
\item An arbitrary infinite mixed sum need not be new (a geometric series
lies in $M$); the distinct-coset hypothesis is necessary; finite support
cannot work over the full compositum; not every incidence family gives
independence, and pairwise distinct supports are not enough
(\cref{isc:hj:sub:notimply}).
\item An independent set of the right cardinality need not be a
transcendence basis; the prime family is not asserted to be one
(\cref{isc:hj:sub:notimply}).
\item Characteristic zero is used twice; the proof does not extend
verbatim to characteristic $p$ (\cref{isc:hj:sub:notimply}; Research
question~\ref{isc:hj:q:charp}).
\item No order-topological convergence, continuity classification of the
non-strong maps or commutation with conjugation is proved; Archimedean
exponent groups do not make the Hahn fields Archimedean
(\cref{isc:hj:sub:notimply}).
\item $\No$ is not the base field; every witness lies in $\No$. A prime of
$\Oz[i]$ is a unit in $\No(i)$; primality is a statement in the ring
specified (\cref{isc:hj:sub:notimply}).
\item The cardinal results depend neither on primality nor on unreviewed
class-embedding theorems; the primality corollary needs the published
theorem; the automorphism count needs algebraic closedness of the complex
join (\cref{isc:hj:sec:scope}).
\item Finite rational-rank factors are not covered: the proof needs
infinitely many rationally independent directions, and the lower bound for
$\K{k}{\Q+\sqrt2\Q}$ is open. Prime witnesses are obtained only for the
countable Archimedean construction, not for uncountable $\kappa$. The
algebraic theorem allows $H\ne0$, but no example with $H\ne0$ is
constructed (Research questions~\ref{isc:hj:q:rankone},
\ref{isc:hj:q:primes} and~\ref{isc:hj:q:cores}).
\item The research questions are proposed by the arguments; they are not
asserted to be previously published open problems. Research
question~\ref{isc:hj:q:formal} is an implementation task
(\cref{isc:hj:sec:questions}).
\item The finite verification is not a proof of the infinite Hahn,
cardinality or primality assertions and does not test primality in $\Oz$.
The floating-point check is orientation only; monotonicity and the limit
are proved analytically. A numeric specialization typo found during
verification was corrected before release (\cref{isc:hj:sub:verification};
README).
\item No Lean source or kernel verification is claimed; the interfaces of
\cref{isc:hj:sub:formal} are proposals, not declarations of the
collection (\cref{isc:hj:sub:formal}).
\item The user-supplied encyclopedia page served as orientation only; a
publisher record for Pitteloud was consulted for bibliographic data only.
The collection catalogue's inclusion of a report is not taken as evidence
of correctness (Appendix~\ref{isc:hj:app:source}).
\item Finite checks, typesetting checks and written proofs are the only
evidence supplied; there has been no external review. The delivered
README's PDF inspection (a contact sheet of rendered pages) does not
verify the mathematics (Appendix~\ref{isc:hj:app:source}; README).
\end{enumerate}
\begin{enumerate}[label=JM\arabic*.,leftmargin=2.8em,itemsep=2pt]
\item Research question~\ref{isc:cp:q:amalgams} is answered only in its
transcendence-defect part and only for the explicit non-separated examples;
the image of the compositum and a denominator criterion stay open. Research
question~\ref{isc:cp:q:cores} is continued, not answered
(\cref{isc:hj:sec:answers}).
\item Source~03 cites no report of the collection except this one. Its
finite-exclusive infinitude is \cref{isc:cp:def:private}; its branch coding
is \lbl{tail:eq:binarycode} shifted by one; the square-root independence is
also used in \lbl{tail:thm:surreal}; the one-row primality is cited in the
Diophantine report; the degrees of \lbl{bst:cor:exactcard} and
\lbl{bst:cor:realgroup} are over a different base
(\cref{isc:hj:sec:collection}).
\item \Cref{isc:hj:rem:secondroute,isc:hj:rem:omnificaut} and the remarks
comparing the sources are observations of the merge, read only in the
merge. \Cref{isc:hj:rem:omnificaut} depends on the unrefereed
\lbl{opa:as:cor:setpairs}.
\item The source was written against the pin, before source~02 was merged;
its statements about this report are those of the pin. Where they are now
incomplete (the degree question in the separated case, the exact
membership test in the separated case, class families in actual copies),
\mtag{} notes say so.
\item No \texttt{isc:hj:} label has a Lean mapping in
\rpath{FORMALIZATION.md}. The proofs were read when the addition was
merged (Appendix~\ref{isc:hj:app:source}); that reading is not an
independent review. The wording of \cref{isc:hj:thm:classicalprime} was
compared with \cite[Section~1.3]{LM}; the other new references were not
checked.
\item The recorded run was repeated in the merge on a copy; the shipped
\rpath{code/03-hahn-joins-build.sh} does not run unchanged under the
shipped names, and an adapted copy would overwrite a PDF and a
verification record in its own directory (Appendix~\ref{isc:hj:app:files}).
\end{enumerate}

\section{Research questions of the third source}\label{isc:hj:sec:questions}

The questions below are proposed by the source's arguments. They are not
asserted to be previously published open problems. The first two are
especially direct next targets, and the final question is an
implementation task rather than a conjecture about truth.
\mtag{} The source's Question~$n$ is Research question~$22+n$ here.

\begin{question}[Finite rational-rank factors]\label{isc:hj:q:rankone}
Determine
\[
 \trdeg\bigl(\K{k}{\Q+\sqrt2\Q}/
              \K{k}{\Q}\K{k}{\sqrt2\Q}\bigr),\qquad k=\R,\C.
\]
Is it $\cc$? The ambient cardinality gives that upper bound, but the
coordinate-annihilator proof has only one rational direction from each
factor and therefore does not prove the lower bound. A new obstruction
that survives finite differential rank is needed.
\end{question}
\mtag{} This is the smallest dense case of
Research question~\ref{isc:cp:q:amalgams}. The rank test of
\cref{isc:cp:thm:rank} does not apply, since neither group is convex in
the sum.

\begin{question}[Exact characterization of the ordinary compositum]\label{isc:hj:q:characterization}
Find a necessary and sufficient condition, expressed in coefficients or
supports, for an element of $\K{k}{A+B}$ to lie in
$\K{k}{A}\K{k}{B}$. The mixed-support criterion supplies a large
family outside even the relative algebraic closure, but not a
membership test. The separated coefficient-field condition discussed
in this report (\cref{isc:lem:finitecoeff}) is a useful comparison point,
not a replacement for such a characterization in interleaved groups.
\end{question}
\mtag{} Since the pin, the separated case has an exact test, the finite
coefficient-rank denominator criterion (\cref{isc:cp:thm:rank,isc:cp:cor:surrealrank}),
and the coefficient-field condition is known not to be sufficient
(\cref{isc:cp:cor:fginsufficient}). For interleaved groups the question is
the first part of Research question~\ref{isc:cp:q:amalgams}, and stays open.

\begin{question}[Weaker incidence hypotheses]\label{isc:hj:q:incidence}
Can finite-exclusive infinitude be replaced by an exact rank condition
on the incidence functions $1_U$ modulo finitely supported functions?
Finite linear identities among these functions already give polynomial
relations among the series. Determine whether a suitable higher-order
rank condition captures all algebraic relations for coordinate-
independent mixed exponents.
\end{question}

\begin{question}[Prescribed shared Hahn cores]\label{isc:hj:q:cores}
For $H\ne0$, give sharp sufficient conditions on
$\Gamma_0/H$ and $\Gamma_1/H$ for the existence of a well-ordered mixed
support satisfying \cref{isc:hj:main:mixed} with $2^\kappa$ incidence patterns.
The algebraic theorem already allows a nontrivial intersection; the
remaining problem is the ordered construction of enough suitable
representatives and the corresponding optimal cardinal upper bound.
\end{question}
\mtag{} It continues Research questions~3 and~\ref{isc:cp:q:cores}; the
latter asks for tail resilience over a prescribed core, which is not
addressed here.

\begin{question}[Full surreal copies rather than support subfields]\label{isc:hj:q:copies}
For pairs of full monomial-preserving copies of $\No$ over a prescribed
Hahn core, determine when their exponent images contain the coordinate
configurations of \cref{isc:hj:main:mixed} at every set cardinal.
A positive result would transfer the no-set-basis conclusion to those
full-copy diagrams. The explicit class groups in \cref{isc:hj:sec:classes}
resolve a different case and should not be silently substituted for
full copies.
\end{question}
\mtag{} For the explicit separated copies of \cref{isc:sec:explicit}, with
trivial core, the no-set-basis conclusion is already known by another
route: no set generates the join over the compositum even algebraically
(\cref{isc:cp:cor:nosetgen}); and one configuration of
\cref{isc:hj:main:mixed} is exhibited there by \cref{isc:hj:rem:secondroute}.
Prescribed nonzero cores remain open (Research question~\ref{isc:cp:q:cores}).

\begin{question}[Uncountable prime witnesses]\label{isc:hj:q:primes}
For $\kappa>\aleph_0$, can the $2^\kappa$ independent omnific integers
in \cref{isc:hj:main:cardinal} be replaced by primes in $\Oz$, or in both
$\Oz$ and $\Oz[i]$? The one-row theorem treats the countable
real-exponent construction, not arbitrary uncountable nested supports.
The question asks for a genuinely compatible primality mechanism, not
merely an independent family with constant term one.
\end{question}

\begin{question}[Compatibility of the hidden complex automorphisms]\label{isc:hj:q:automorphisms}
Which automorphisms counted in \cref{isc:hj:thm:automorphisms} can commute
with conjugation, preserve the real Hahn subfield, or preserve its omnific
subring? The argument by permutations of a transcendence basis imposes
none of these constraints. A classification would measure precisely how
much additional rigidity each distinguished structure contributes.
\end{question}
\mtag{} \emph{Partial status.} By \cref{isc:hj:rem:omnificaut}, which rests
on \lbl{opa:as:cor:setpairs}, only the identity commutes with conjugation
and preserves the Gaussian omnific ring of the join. The other combinations
stay open.

\begin{question}[More than two factor fields]\label{isc:hj:q:factors}
Develop a mixed-support criterion for
$K_1\cdots K_r$ inside the full Hahn field on
$\Gamma_1+\cdots+\Gamma_r$. One can annihilate all but one coefficient
field and finitely many generators of the remaining field; what coset
or tensor-independence hypothesis ensures that the surviving mixed
terms do not cancel? This formulation isolates the exact place where
the two-field proof would have to be strengthened.
\end{question}
\mtag{} Here the indices $1,\ldots,r$ are the source's; they number $r$
arbitrary factors. The question continues
Research question~\ref{isc:cp:q:scales}, which asks the differential
version for three or more separated copies with commuting Euler
derivations.

\begin{question}[Positive characteristic]\label{isc:hj:q:charp}
Find an analogue using higher or iterative derivations, or identify a
sharp obstruction caused by divisible exponent groups and Frobenius.
A meaningful statement must treat inseparable algebraic relations and
cannot simply replace the ground field in \cref{isc:hj:main:mixed}.
\end{question}

\begin{question}[A small formal verification kernel]\label{isc:hj:q:formal}
Formalize the support-coset lemma, additive-character Hahn derivations,
and the finite annihilator argument in Lean. These three pieces should
be separated from the normal-form identification with $\No$ and from
published primality. A verified generic Hahn theorem would already
cover the main new mathematical mechanism; \cref{isc:hj:sub:formal} lists
suggested interfaces.
\end{question}
\mtag{} It overlaps Research questions~10 and~\ref{isc:cp:q:formal} in the
coset and Euler-derivation kernel; none of the proposed interfaces exists
in the collection's Lean library.

\section{Conclusion of the third source}\label{isc:hj:sec:conclusion}

An ordinary compositum has access to arbitrary series in each factor,
but only finitely many such series in any one algebraic expression.
The mixed-support theorem exploits that finiteness without making a
finiteness assumption on their supports. A finite homogeneous system
removes the coefficient dependence, and distinct exponent cosets keep
the desired derivative nonzero.

This yields maximal transcendence gaps at every infinite set cardinal,
a prime-witness version with interleaved real exponents, and a fixed
class example with no set-sized transcendence basis. It also separates
ordinary algebraic symmetries from symmetries of the Hahn presentation:
in the complex examples the former are as numerous as possible while
the relative strong subgroup is trivial. The finite-rank compositum
problem and compatibility with full surreal-copy structures remain
particularly clear next challenges.

\section{Proof audit, finite verification and a formalization roadmap of the third source}
\label{isc:hj:sec:verify}

\mtag{} This section prints the source's Appendices~A, B and~C.

\subsection{Dependency ledger}\label{isc:hj:sub:audit}

\begin{center}
\small
\begin{tabular}{>{\raggedright\arraybackslash}p{0.25\textwidth}>{\raggedright\arraybackslash}p{0.41\textwidth}>{\raggedright\arraybackslash}p{0.25\textwidth}}
\toprule
Result & Essential dependencies & Evidence in the source\\
\midrule
Mixed-support independence & Hahn multiplication; set basis extension;
finite linear algebra; characteristic zero & Full proof in
\cref{isc:hj:sec:lemmas,isc:hj:sec:mixed}\\[4pt]
Maximal cardinal gap & Mixed theorem; independent incidence family;
explicit monomial coordinates & Full proof in
\cref{isc:hj:sec:combinatorics,isc:hj:sec:cardinal}\\[4pt]
Archimedean gap & Mixed theorem; square-root independence; binary branches &
Full proof in \cref{isc:hj:sec:archimedean}\\[4pt]
Omnific primality & Classical one-row prime theorem; constant-term pullback;
set localization & External theorem isolated; application proved in
\cref{isc:hj:sec:prime}\\[4pt]
Class-sized example & Explicit class coordinates; finite local proof &
Full argument in \cref{isc:hj:sec:classes}; no class basis selection\\[4pt]
Relative derivations and complex automorphisms & Transcendence-degree
results; separability; algebraic-closure extension & Proofs in
\cref{isc:hj:sec:strong}\\
\bottomrule
\end{tabular}
\end{center}

\subsection{Potential failure points that have been checked}

\paragraph{Support existence.}
The abstract theorem explicitly assumes a well-ordered mixed
$t$-support. The nested example supplies a strictly increasing negative
$t$-support indexed by $\kappa$, and the prime example supplies an
increasing sequence with supremum zero. A supremum is not required to
belong to a Hahn support.

\paragraph{Sign convention.}
Positive Conway exponents are negative $t$-exponents. The candidates
$\Xi_U$ and $P_\eta$ therefore lie on the omnific side of the support
cut, not on the infinitesimal side.

\paragraph{Coefficient-field membership.}
Each entry $\partial_i(u_q)$ lies in $K_0$, not merely in the ambient join.
This is why the homogeneous system can be solved with its coefficients
$c_i$ in $K_0$ and why the final coset-independence step is valid.
Solving over the entire join instead would not establish the needed
noncancellation.

\paragraph{Dependence of the annihilator on the relation.}
The annihilator is chosen after all the coefficients of one polynomial
have been placed in a finitely generated field over $K_1$.
There is no claim that a single strong derivation annihilates the full
compositum; \cref{isc:hj:sec:strong} proves that it cannot do so
nontrivially.

\paragraph{Nonzero polynomial partial.}
The relation is chosen of minimal total degree over a fixed coefficient
field. A nonzero formal partial has smaller degree, so its evaluation
cannot vanish. Merely choosing any polynomial relation would not
justify this step.

\paragraph{Primality after complexification.}
Primality of a real element need not survive an arbitrary ring
extension. We do not infer the Gaussian assertion from the real one.
Instead we apply the characteristic-zero one-row theorem separately
with complex coefficients and then descend along the constant-term
map to $\Z[i]$.

\paragraph{Class localization.}
Any factorization identity in the full omnific ring uses a set of
support exponents. Enlarging that set by $\R$ still yields a set group,
so the external theorem applies within its stated set-sized scope.
The class independence theorem uses explicit coordinate characters and
finite systems rather than a proper-class choice of basis.

\paragraph{Cardinality and transcendence.}
The lower bound is supplied by proved algebraic independence, not by
counting different series. The upper bound counts all coefficient
functions on a group of cardinality $\kappa$. No continuum hypothesis
or assumption that $\kappa$ is regular is used.

\subsection{Strengthening the multiquadratic justification}

For clarity, the finite Galois fact used in \cref{isc:hj:sec:archimedean}
can be derived without appealing to an infinite construction. If $F$ has
characteristic not two and $d$ is not a square in $F$, then an element
$r\in F^\times$ becomes a square in $F(\sqrt d)$ only if either $r$ or
$r/d$ is a square in $F$. Indeed, writing a square root as $u+v\sqrt d$
gives $2uv=0$; the two possibilities give exactly those alternatives.
Induction shows that the kernel of
\[
 \Q^\times/(\Q^\times)^2\longrightarrow
 \Q(\sqrt{p_1},\ldots,\sqrt{p_m})^\times/
 \bigl(\Q(\sqrt{p_1},\ldots,\sqrt{p_m})^\times\bigr)^2
\]
is generated by the square classes of the $p_j$. Unique prime
factorization says a new prime has a square class outside this subgroup,
so every adjoining step doubles the degree. The finite extension has
the expected $2^m$ monomial basis, and independently changing the signs
of its generators preserves the multiplication relations. This supplies
the sign-changing automorphisms used to prove linear independence of
the square roots.
\mtag{} In this paragraph $F$ is a field of characteristic not two, $r$ an
element and $p_1,\ldots,p_m$ distinct primes listed afresh, not the
enumeration $p_0<p_1<\cdots$ of \eqref{isc:hj:eq:realexponents}.

\subsection{A finite exact verification model}\label{isc:hj:sub:verification}

The source's program, shipped as \rpath{code/03-hahn-joins-verify.py},
checks finite algebraic and combinatorial instances. These checks are not a
proof of the infinite Hahn assertions and do not test primality in $\Oz$.
Their purpose is to catch mistakes in the annihilator calculation, the
incidence coding, and the explicit finite formulas.

For a concrete annihilator, let $x_1,x_2,x_3$ be independent formal
variables and put
\[
 u_1=x_1+x_2+x_3,\qquad u_2=x_1x_2+x_3^2,
 \qquad \partial_j=x_j\frac{\partial}{\partial x_j}.
\]
The column vectors of the two derivative lists are
\[
 v=(x_1,x_2,x_3),\qquad
 w=(x_1x_2,x_1x_2,2x_3^2).
\]
Their cross product gives
\[
 \begin{aligned}
 c_1&=2x_2x_3^2-x_1x_2x_3,\\
 c_2&=x_1x_2x_3-2x_1x_3^2,\\
 c_3&=x_1x_2(x_1-x_2).
 \end{aligned}
\]
Thus $\partial_{\mathbf c}=\sum c_j\partial_j$ annihilates $u_1$ and $u_2$
exactly. If $y_1,y_2,y_3$ are additional independent variables and
$F=x_1y_1+x_2y_2+x_3y_3$, then
\[
 \partial_{\mathbf c}(F)=c_1x_1y_1+c_2x_2y_2+c_3x_3y_3\ne0.
\]
At $(x_1,x_2,x_3)=(1,2,3)$ this becomes
$30y_1-24y_2-6y_3$, visibly nonzero.
This is the finite polynomial shadow of the coset argument.
\mtag{} In this model $F$ is a polynomial, not a Hahn series.

The verification script uses exact symbolic arithmetic, enumerates
small Boolean-pattern instances, checks exclusive binary tails for
several distinct eventually periodic branches, and verifies the displayed
cross-product identities and quotient-rule calculation. Numerical
approximations to the first prime exponents are included only as an
orientation check; the monotonicity and limit are proved analytically
above. The actual output is shipped as
\rpath{data/03-hahn-joins-verification.txt}.
\mtag{} Its five PASS lines are: the exact cross-product annihilator for two
generators and a quotient; the nonvanishing mixed derivative, with the
specialization $30y_1-24y_2-6y_3$; all $34{,}112$ Boolean patterns on up to
four distinct subsets of a four-point universe (the finite analogue of
\cref{isc:hj:lem:independentfamily}; the coordinate $\xi$ is not modelled);
eight distinct eventually periodic branches through depth~32, each with at
least $29$ exclusive nodes, and the codes $\{1,3,7,15\}$, $\{2,6,14,30\}$ of
the two constant branches; and the decrease of the first ten exponents
$\rho_n$ in floating point. The rerun is in Appendix~\ref{isc:hj:app:files}.

\subsection{A formalization roadmap}\label{isc:hj:sub:formal}

No Lean source or kernel verification is claimed for the new theorem.
The following are proposed interfaces, not declarations alleged to
exist in the repository.

\paragraph{Hahn support layer.}
Prove that multiplying a series supported in a subgroup $\Gamma'$ by a
monomial of exponent $\beta$ places its support in $\Gamma'+\beta$.
Deduce finite linear independence for monomials in distinct cosets.
This part needs neither divisibility nor characteristic zero.

\paragraph{Character derivation layer.}
For an additive character $\chi:\Gamma\to k$, define the coefficient
map $x_\gamma\mapsto\chi(\gamma)x_\gamma$ on Hahn series. Establish
support containment, the coefficientwise Leibniz rule, and commutation
with summable families. Then construct the coordinate characters from
linear independence in $\Gamma/\Gamma_1$ using a set vector-space basis.

\paragraph{Finite annihilator layer.}
Package an $m\times(m+1)$ matrix over $K_0$ and obtain a nonzero kernel
vector. Form the finite linear combination of derivations and prove that
it annihilates the intermediate field generated by $K_1$ and the $u_q$.
Keep its coefficients in $K_0$ throughout; this is a useful type-level
invariant that prevents an invalid noncancellation step.

\paragraph{Algebraic independence layer.}
Use minimal total degree, or an equivalent separable minimal-polynomial
argument, to formalize the contradiction. The finite-exclusive support
condition can be abstracted as the ability to select $m+1$ indices for
any finite relation. Cardinal amplification should be a separate theorem.

\paragraph{Surreal and arithmetic bridges.}
Only after the generic Hahn theorem is verified should one transport
it through Conway normal form and add omnific support conditions.
The one-row prime theorem is a distinct formalization project and must
not be silently accepted as checked merely because the independence
kernel is checked. The constant-term pullback lemma itself is an
accessible independent commutative-algebra target.

\mtag{} The coset layer overlaps item~(1) of \cref{isc:sub:formal} and the
support and base-change layer of \cref{isc:cp:sub:formal}; the character
derivation layer overlaps its derivation layer.

\appendix
\section{Provenance of this report}\label{isc:app:report}

\subsection{One manuscript}\label{isc:app:onesource}

This report is built from \emph{one} manuscript: the archive
\texttt{Surreal\_Independent\_Copies}, manuscript~06 of batch~31 of the
collection, with main file \texttt{independent\_surreal\_copies.tex}, a
25-page PDF, dated 23 September 2026 and pinned to repository commit
\nolinkurl{bcac55ae6fd3f2b354e568ba1d2e94d496a2cc59}. It was placed in
commit \texttt{9d28e28} as a new report. It answers no named question of
the collection, and its subject, independence, disjointness and descent
for full copies of $\No$, is the subject of no existing report
(\cref{isc:sec:collection}). There was no second source. Nothing was
merged and nothing was selected away: every result, proof, example,
question and limitation of the manuscript is printed here.
\mtag{} This subsection describes the report as written in batch~31. In
batch~34 a second manuscript was added as source~02
(Appendix~\ref{isc:cp:app:source}), and in batch~35 a third as source~03
(Appendix~\ref{isc:hj:app:source}); the report is now built from three
manuscripts.

The package contained the article, its PDF and a README, and nothing else:
no code, recorded data, audit or checksum list. The article was renamed
\texttt{article.tex} and rewritten as described below; the README was
replaced by the collection README; the delivered PDF is not shipped, and
\texttt{article.pdf} is a build of this text.

The editorial changes are these.
\begin{enumerate}[label=(\arabic*)]
\item Every label carries the prefix \texttt{isc:}. The source's 62 labels
are kept, unchanged after the prefix. Twenty were added:
\lbl{isc:sec:conventions}, \lbl{isc:sec:collection},
\lbl{isc:sec:nonclaims}, \lbl{isc:rem:base}, \lbl{isc:rem:floor},
\lbl{isc:warn:omega}, \lbl{isc:rem:join}, \lbl{isc:rem:scaling},
\lbl{isc:rem:infinite}, \lbl{isc:sub:model}, \lbl{isc:sub:bounded},
\lbl{isc:sub:formal}, \lbl{isc:sub:notmade}, \lbl{isc:app:report},
\lbl{isc:app:onesource}, \lbl{isc:app:pinned}, \lbl{isc:app:literature},
\lbl{isc:app:checked}, \lbl{isc:app:files} and \lbl{isc:app:numbers}. Labels
on existing environments change no number.
\item The title page names the report as a single-source report of the
collection, gives the batch and the pin, and extends the status paragraph.
\Cref{isc:sec:conventions,isc:sec:collection}, the ledger of
\cref{isc:sec:nonclaims} and this appendix are new. Sentences marked
\wtag{} were added at the end of \cref{isc:subsec:provenance}; after
\cref{isc:prop:classification}, \cref{isc:cor:parameters},
\cref{isc:thm:intersectionclassification}, \eqref{isc:eq:explicitouter},
\cref{isc:lem:denominator}, \cref{isc:prop:uniform} and
\cref{isc:prop:boundedfrac}; at the start of \cref{isc:sec:gaps}; after
\cref{isc:cor:factorial}; after research questions 1, 2, 6, 7, 8, 9 and~10;
and in \cref{isc:sub:notmade}.
\end{enumerate}
No mathematical statement was changed and no symbol was renamed: the
collisions are between reports, or between local letters of different
sections, and \cref{isc:sec:conventions} fixes the meanings. Where the edit
had to choose: the manuscript was placed as a new report, not as an
addition to the omnific-preserving report, because it answers none of that
report's questions and its spine (independence, exact intersections and
descent) is not that report's (stabilizers, classification and
coefficient motion). It was placed among the surreal reports because its
fields are subfields of $\No$ and $\SC$. The title is kept.

\subsection{Repository statements at the pin}\label{isc:app:pinned}

The source compared itself with the catalogue, the omnific-preserving
report and the bounded-support report at its pin, through links to that
revision (\cite{RepoCatalogue,RepoAuto,RepoBounded}). Checked at the
placement commit:
\begin{itemize}[leftmargin=*,itemsep=3pt]
\item At the pin the omnific-preserving report was the merged report of
four manuscripts subtitled ``Convex-scale stabilizers, definable constants,
nondefinable monomials, and algebraic-parameter rigidity''. It had no part
on embeddings and no \texttt{opa:as:} label. The source's description of it
as concerning stabilizers and rigidity was accurate. Its Part~III, on
automatic strongness and proper embeddings, was written in
\texttt{20c4c9f}, after the pin.
\item The bounded-support report is unchanged between the pin and the
placement commit (its last change, \texttt{ce6715b}, precedes the pin). The
source's description of it as proposing stronger independent families
stands.
\item The large-cardinal report was placed in \texttt{21375f8}, after the
pin, and is not cited by the source.
\item No report of the collection supplies \cref{isc:main:copies},
\cref{isc:thm:sliceddisjoint}, \cref{isc:main:fractions} for copies, or
\cref{isc:main:join}; the overlaps are those of \cref{isc:sec:collection}.
\end{itemize}

\subsection{References checked for this report}\label{isc:app:literature}

\begin{itemize}[leftmargin=*,itemsep=3pt]
\item \emph{L'Innocente--Mantova} \cite{LM}. The arXiv page lists version~5
of 22 January 2024 and the journal DOI
\nolinkurl{10.1016/j.aim.2024.109513}; the volume number was not checked.
In the arXiv text, Fact~2.1.1 states that a Hahn field is algebraically
closed, respectively real closed, if and only if its coefficient field is
algebraically closed, respectively real closed, and its group is divisible.
Section~2.3 introduces the surreal numbers as a Hahn field. Proposition~2.4.5
is stated for the Hahn field of series with supports of size below a
cardinal $\kappa$: the fraction field of $\Z+K((G^{<0}))$ is the whole field
if and only if $\cf(G)\ge\kappa$ (or $G=0$ and $\Frac\Z=K$), and in
particular $\Frac(\Oz)=\No$. For a set-sized $G$ and $\kappa>\lvert G\rvert$ this is
the full Hahn field, so its proof, with a lacunary series along a cofinal
sequence, gives the non-membership obstruction cited in
\cref{isc:subsec:provenance} and \cref{isc:sec:gaps}. It proves
non-membership in the fraction field, not transcendence. The citations are
therefore accurate.
\item \emph{Ehrlich--Kaplan} \cite{EK}. The author-hosted PDF (37 pages) has
the stated title and authors; only its first page was read.
\item Gonshor \cite{Gonshor}, van den Dries--Ehrlich \cite{vdDE}, Marker
\cite{Marker} and the encyclopedia page \cite{Wiki} were not checked.
\end{itemize}

\subsection{Proofs read when the report was written}\label{isc:app:checked}

The proofs of \cref{isc:lem:avoid,isc:lem:extend}, \cref{isc:thm:group},
\cref{isc:thm:lift}, \cref{isc:prop:arithmeticlift},
\cref{isc:prop:classification}, \cref{isc:sec:slices}, \cref{isc:main:copies},
\cref{isc:thm:intersectionclassification,isc:thm:boolean},
\cref{isc:cor:choices}, \cref{isc:sec:explicit,isc:sec:join},
\cref{isc:sec:fractions,isc:sec:gaps} and \cref{isc:thm:setversion} were
read line by line when the report was written. No error was found. The
explicit formulas were rechecked: $\Phi_0(\omega)=\omega^{\omega^{1/2}}$,
$\Phi_1(\omega)=\omega^{\omega^{7/2}}$, $g=h_1(1)=\omega^{7/2}$, and the
consecutive exponents of \eqref{isc:eq:explicitS} differ by
$g+g_n-g_{n+1}>0$. One understatement is recorded after
\cref{isc:prop:uniform}. This reading is not an independent proof review.

\subsection{Theorem numbers}\label{isc:app:numbers}

Numbers are those of this build and of the delivered PDF, which agree:
labels were added and no environment was inserted before an existing one.
Remarks and warnings share the theorem counter, so the classification of
monomial-preserving embeddings is Proposition~\ref{isc:prop:classification}
(after Remark~\ref{isc:rem:floor}), and the saturated-group theorem is
Theorem~\ref{isc:thm:setversion} (after a definition).

\subsection{Files}\label{isc:app:files}

\begin{center}\small
\begin{tabular}{@{}ll@{}}
\toprule
File & Role \\
\midrule
\rpath{article.tex} & this report (delivered as \texttt{independent\_surreal\_copies.tex})\\
\rpath{article.pdf} & a build of this text, not the delivered PDF\\
\rpath{README.md} & the collection README, replacing the delivered one\\
\bottomrule
\end{tabular}
\end{center}

The source shipped no verification program, recorded run or audit, so
there is nothing to rerun. The delivered README reports a build with
pdfTeX~1.40.26 and a visual inspection of rendered pages and contact
sheets; those renderings are not shipped.
\mtag{} This table lists the files of source~01. The seven shipped files of
source~02, with the prefix \texttt{02-composita-}, and their rerun are in
Appendix~\ref{isc:cp:app:files}; the three of source~03, with the prefix
\texttt{03-hahn-joins-}, in Appendix~\ref{isc:hj:app:files}.

\subsection{Source 02 (batch 34)}\label{isc:cp:app:source}

\mtag{} Source~02 is manuscript~05 of batch~34 of the collection: the
archive \nolinkurl{Surreal_Composita_Research}, delivered in commit
\texttt{a4611ad} and placed in commit \texttt{a7a435f} as an addition to
this report with local number~02. Its main file \texttt{article.tex} (a
standalone article with an internal bibliography) is titled \emph{Beyond
Composita of Surreal Copies: Exact rank tests, differential transcendence,
and omnific witnesses}, dated 23 September 2026, with the author line
``AI-assisted research draft prepared for Vladimir Reshetnikov''; its PDF
has 29 pages. It is pinned to commit
\nolinkurl{7b256e0792df7718995b816ad12a77bf56731f8f}, at which this report's
\texttt{article.tex} was blob \nolinkurl{7d5cfe644e7c902f2b594179c337a6c633350fc3}.
The placement chose an addition rather than a new report because the
manuscript answers this report's Research question~3 in the separated case
and extends its witness \eqref{isc:eq:explicitS}. Its manuscript, README and
PDF are not shipped; its code, recorded data and two audits are shipped with
the prefix \texttt{02-composita-} (Appendix~\ref{isc:cp:app:files}).

\emph{What it contributed.} The finite coefficient-rank test for composita
(\cref{isc:cp:thm:rank}), the base-change lemmas of
\cref{isc:cp:sec:disjoint}, the rapid coefficient-degree theorem and its
continuum family (\cref{isc:cp:thm:lacunary,isc:cp:cor:continuum}), the
failure of the finite coefficient-field criterion
(\cref{isc:cp:cor:fginsufficient}), the private-block theorem and tree
amplification over the finite-generation envelope
(\cref{isc:cp:thm:generic,isc:cp:cor:tree,isc:cp:prop:local,isc:cp:cor:sandwich}),
the exact Laurent degrees $|F|^{\aleph_0}$ (\cref{isc:cp:thm:cardinal}), the
ordinal family $Y_\alpha$ of purely infinite omnific integers with
independent Euler jets (\cref{isc:cp:thm:omnific,isc:cp:cor:threshold}),
set-parameter tail resilience (\cref{isc:cp:thm:resilience,isc:cp:cor:nosetgen}),
the arithmetic and Gaussian consequences
(\cref{isc:cp:thm:arithmetic,isc:cp:thm:complex}), twelve research questions
and a formalization decomposition.

\emph{Editorial changes.}
\begin{enumerate}[label=(\arabic*)]
\item Every label of the source carries the prefix \texttt{isc:cp:}, with
the source's name after it, except that its \texttt{sec:conventions} is
\lbl{isc:cp:sec:foundations} and its appendix labels are replaced by
\lbl{isc:cp:sec:verify} and \lbl{isc:cp:sub:formal}. The merge also labels
the source's unlabelled definition of coefficient rank, example, remarks
and questions, and its new subsections and appendix subsections.
\item Symbols are renamed as in \cref{isc:cp:sec:conventions}; no
normalization changed.
\item The source's re-derivation of the explicit copies and of the
re-expansion (its Subsections~8.1--8.2) is printed once, in
\cref{isc:sec:explicit,isc:sec:join}, with the correspondence in
\cref{isc:cp:sub:copies}; the source's added observations are kept there.
Its duplicate of \cref{isc:lem:denominator} is kept inside
\cref{isc:cp:lem:clearing}, where it serves a new statement.
\item The source's citations of this report by section number are replaced
by cross-references to the same sections. Its bibliography entries
\cite{KKS,KM,BM,StacksLD} are added to this report's; its entry for
L'Innocente--Mantova is this report's \cite{LM}. These four references
were not checked in the merge.
\item Text marked \mtag{} was added in the merge: in
\cref{isc:cp:sec:intro} (the first paragraph, the note on $Y_0$, and
\cref{isc:cp:sec:answers,isc:cp:sec:conventions,isc:cp:sec:collection}),
after \cref{isc:cp:thm:rank,isc:cp:lem:coset,isc:cp:cor:continuum,isc:cp:ex:mechanism,isc:cp:thm:omnific,isc:cp:lem:localizeP},
in \cref{isc:cp:sub:copies}, in the construction of the $a_{\alpha,n}$, in
\cref{isc:cp:sub:nonclaims}, after Research
questions~\ref{isc:cp:q:cores}, \ref{isc:cp:q:bm},
\ref{isc:cp:q:normalization} and~\ref{isc:cp:q:formal}, in
\cref{isc:cp:sec:verify}, and in this appendix. In source~01's text, \mtag{}
notes were added to the title page, to \cref{isc:sec:collection} (the status
of the research questions), after \cref{isc:cor:parameters} (with the
reciprocal note for the dilation report), after
\cref{isc:lem:finitecoeff}, after the proof of \cref{isc:main:join}, after
Research question~3, at the start of \cref{isc:sec:nonclaims}, in item~W1,
and in Appendices~\ref{isc:app:onesource} and~\ref{isc:app:files}. No
statement of source~01 was changed.
\item Cross-references to lemmas, propositions, corollaries, remarks and
the warning were printed as ``theorem'' before the merge: the environments
share the theorem counter, and cleveref named them by the counter. The
merge adds cleveref's type hint to the labels of these environments (26 of
source~01 and 33 of source~02; for example the label of
\cref{isc:lem:finitecoeff} now carries the hint \texttt{[lemma]}) and a
cleveref name for the warning. No label was renamed and no number changed; only the printed names
of those references are corrected.
\end{enumerate}

\emph{Where the merge chose.} The addition is printed after the ledger of
non-claims and before the appendix, so that no existing section, theorem,
equation or question number changes; its sections keep the source's
numbering shifted by fourteen. The outer group is renamed $\Gamma$ rather
than keeping $H$, because $H$ is this report's common core and the source
uses it in that sense too. The explicit copies are not reprinted.

\emph{Checks in the merge.} The proofs of
\cref{isc:cp:thm:rank,isc:cp:lem:basechange,isc:cp:lem:coset,isc:cp:cor:transfer,isc:cp:thm:lacunary,isc:cp:lem:partial,isc:cp:thm:generic,isc:cp:cor:tree,isc:cp:prop:local,isc:cp:thm:cardinal,isc:cp:thm:omnific,isc:cp:lem:tailcoeff,isc:cp:thm:resilience,isc:cp:lem:clearing,isc:cp:thm:arithmetic,isc:cp:thm:complex}
were read line by line when the addition was merged; no error was found.
The following were rechecked: $g=\omega^{7/2}$ and $b=\omega^{15/4}$ lie in
$G_1$ (inner exponents in $(3,4)$), $\lambda(g)=1$, $\lambda(b)=0$,
$\Delta\tau=\tau$ and $\Delta\omega^b=0$ from \eqref{isc:cp:eq:Delta};
$\beta_{0,n}=n$, so $g_{0,n}=\omega^{\psi(n)}$ and
$Y_0=\omega^{\omega^{15/4}}S$ with $S$ as in \eqref{isc:eq:explicitS}; the
consecutive exponents of \eqref{isc:cp:eq:Ynormal} differ by
$g+g_{\alpha,n}-g_{\alpha,n+1}>0$; and the source's numbering (its
Theorem~8.2 is the omnific theorem, its Theorem~10.4 the surcomplex one)
agrees with its delivered PDF. This reading is not an independent proof
review. The source's statements about this report are those of its pin;
the parts it used are unchanged since.

\subsection{Files of source 02, and their rerun}\label{isc:cp:app:files}

\mtag{} The seven shipped files of source~02 are these.
\begin{center}\small
\begin{tabular}{@{}p{0.40\textwidth}p{0.22\textwidth}p{0.30\textwidth}@{}}
\toprule
Shipped path & Delivered name & Role\\
\midrule
\rpath{02-composita-SOURCE\_AUDIT.md} & \texttt{SOURCE\_AUDIT.md} & pin,
source scope, priority limits\\
\rpath{02-composita-PROOF\_AUDIT.md} & \texttt{PROOF\_AUDIT.md} &
hypotheses, critical steps, nonclaims\\
\rpath{code/02-composita-verify.py} & \texttt{verify.py} & exact finite
checks (SymPy)\\
\rpath{code/02-composita-build.sh} & \texttt{build.sh} & runs the checks,
then builds the source's PDF\\
\rpath{data/02-composita-verification.json} & \texttt{verification.json} &
recorded run: 2,530 checks, 21 categories\\
\rpath{data/02-composita-requirements.txt} & \texttt{requirements.txt} &
\texttt{sympy==1.14.0}\\
\rpath{data/02-composita-BUILD\_AUDIT.json} & \texttt{BUILD\_AUDIT.json} &
build and PDF inspection record of the delivered PDF\\
\bottomrule
\end{tabular}
\end{center}
All seven are byte-identical to the delivery. The manuscript, its README
and its PDF are not shipped; the build audit therefore describes a PDF (29
pages, 499{,}044 bytes, pdfTeX~1.40.26) that is not in the collection.

The program writes a file only when given \texttt{--output}; otherwise it
prints its JSON report. The delivered \texttt{build.sh} changes into its own
directory and runs \texttt{python3 verify.py --output verification.json}
before building \texttt{article.tex}. Under the shipped names it finds
neither file in \rpath{code/}, and an adapted copy would write its record
next to the script, not over the shipped data file. Rerun the checks on a
copy, or with an explicit output path outside the report:
\begin{verbatim}
python code/02-composita-verify.py --output <scratch>/verification.json
\end{verbatim}
In the merge the program was run on a copy with Python~3.14.4 and
SymPy~1.14.0. It passed all 2,530 checks with the same 21 category counts
as the recorded run; the only difference in the JSON record is the Python
version (3.13.5 in the recorded run).

\subsection{Source 03 (batch 35)}\label{isc:hj:app:source}

\mtag{} Source~03 is manuscript~01 of batch~35 of the collection: the
archive \nolinkurl{Surreal_Hahn_Joins_Research} (inner directory
\nolinkurl{maximal_transcendence_hahn_joins}), delivered in commit
\texttt{267b910} and placed in commit \texttt{dec8d56} as an addition to
this report with local number~03. Its main file \texttt{article.tex} (a
standalone article with an embedded bibliography) is titled
\emph{Maximal Transcendence in Hahn Joins: Mixed-support independence and
prime omnific integers}, dated 23 September 2026, with the author line
``AI-assisted research manuscript''; its PDF has 27 pages (cover, contents
and 23 numbered pages). It is pinned to commit
\nolinkurl{f388f950f44664f21f50f04b8ac74dab56c8ff2f}, at which this
report's \texttt{article.tex} was blob
\nolinkurl{7d5cfe644e7c902f2b594179c337a6c633350fc3}: the same snapshot as
source~02's pin, before source~02 was merged in \texttt{b9f964c}. The
placement chose an addition rather than a new report because the
manuscript takes Research question~3 as its target and answers the
transcendence-defect part of Research question~\ref{isc:cp:q:amalgams} for
explicit examples. Its manuscript, README and PDF are not shipped; its
program, build script and recorded run are shipped with the prefix
\texttt{03-hahn-joins-} (Appendix~\ref{isc:hj:app:files}).

\emph{What it contributed.} The mixed-support independence criterion
(\cref{isc:hj:main:mixed}) with its lemmas, generator-sensitive form and
stable variants (\cref{isc:hj:sec:lemmas,isc:hj:sec:mixed}); the incidence
and branch families (\cref{isc:hj:sec:combinatorics}); exact degree
$2^\kappa$ for interleaved real closed and algebraically closed factors
with omnific witnesses and the arithmetic-join gap
(\cref{isc:hj:main:cardinal,isc:hj:cor:realclosuregap,isc:hj:cor:arithmetic});
the Archimedean family $P_\eta$, independent and prime in $\Oz$ and $\Oz[i]$
(\cref{isc:hj:main:prime,isc:hj:prop:archindependence,isc:hj:lem:pullback,isc:hj:prop:fullprime,isc:hj:cor:nonassociate});
the class example (\cref{isc:hj:main:class,isc:hj:lem:classlocal}); strong
rigidity and the counts of relative derivations and automorphisms
(\cref{isc:hj:sec:strong}); ten research questions; a dependency audit, a
finite verification model and a formalization roadmap.

\emph{Sources and provenance, as the source states them (its Appendix~D).}
The repository comparison is pinned to the commit above and was made on
23 September 2026. The primary target was this report's
\texttt{article.tex}, its question headed \emph{An exact description of
ordinary composita} (Research question~3) being the precise starting point.
The collection catalogue \rpath{README.md} at the pin was consulted; it
distinguishes drafts, review and formalization status, and the source does
not infer correctness from inclusion in it or use its unreviewed embedding
theorems as black boxes. For the published prime input the source checked
L'Innocente--Mantova's arXiv version~5 of 22 January 2024, especially its
Section~1.3 and its references to Pitteloud and Biljakovic--Kochetov--Kuhlmann,
which are distinct from the proposed new arguments; a short publisher
record for Pitteloud was consulted to verify bibliographic data. A
user-supplied Wikipedia page served as orientation only, not as a proof
source. The source gives the build commands (three \texttt{pdflatex} passes,
then \texttt{python verify.py}) and says that finite checks, typesetting
checks, mathematical proofs and external review are separate forms of
evidence, of which only the first three are provided.

\emph{Editorial changes.}
\begin{enumerate}[label=(\arabic*)]
\item Every label of the source carries the prefix \texttt{isc:hj:}, with
the source's name after it, except that \texttt{sec:arithmetic} is kept as
\lbl{isc:hj:sec:arithmetic} on the same subsection, the appendix labels
\texttt{app:audit}, \texttt{app:verification} and \texttt{app:formal} become
\lbl{isc:hj:sub:audit}, \lbl{isc:hj:sub:verification} and
\lbl{isc:hj:sub:formal}, \texttt{app:sources} becomes this appendix
\lbl{isc:hj:app:source}, and \texttt{main:mixed}, \texttt{main:cardinal},
\texttt{main:prime}, \texttt{main:class} become
\lbl{isc:hj:main:mixed} etc. The merge also labels the source's unlabelled
definitions, questions, conclusion and new subsections, and gives every
lemma, proposition, corollary, definition and remark label a cleveref type
hint.
\item Symbols are renamed as in \cref{isc:hj:sec:conventions}; no
normalization changed. The source's Theorems~A--D print as
Theorems~D--G, and its Questions~1--10 as Research questions~23--32.
\item The source's citations of this report are replaced by
cross-references to the same places; its hard-coded section numbers in the
dependency ledger and the audit are replaced by cross-references. Its
bibliography entries \cite{Conway,Hahn,Kaplansky,Pitteloud,BKK} are added
to this report's; its entries for Gonshor and L'Innocente--Mantova are this
report's \cite{Gonshor,LM}; its entries for this report and for the
catalogue at its pin are replaced by cross-references and by
\rpath{README.md}. Only the wording of \cref{isc:hj:thm:classicalprime} was
compared with \cite{LM} in the merge; the new entries were not checked.
\item Text marked \mtag{} was added in the merge: in
\cref{isc:hj:sec:intro} (the first paragraph and
\cref{isc:hj:sec:answers,isc:hj:sec:conventions,isc:hj:sec:collection,isc:hj:sub:compare}),
after the fraction-field sentence of \cref{isc:hj:sec:foundations}, after
\cref{isc:hj:prop:ld,isc:hj:lem:euler,isc:hj:def:exclusive}, after
\cref{isc:hj:cor:algbase}, after \eqref{isc:hj:eq:branchsets}, after
\cref{isc:hj:cor:arithmetic}, after \cref{isc:hj:lem:realgroups}, after
\cref{isc:hj:thm:classicalprime,isc:hj:lem:pullback}, after
\eqref{isc:hj:eq:classwitness}, after the proof of \cref{isc:hj:main:class},
\cref{isc:hj:rem:secondroute,isc:hj:rem:omnificaut}, in
\cref{isc:hj:sec:scope} (after the table, and
\cref{isc:hj:sub:nonclaims}), at the head of \cref{isc:hj:sec:questions}
and after Research
questions~\ref{isc:hj:q:rankone}, \ref{isc:hj:q:characterization},
\ref{isc:hj:q:cores}, \ref{isc:hj:q:copies}, \ref{isc:hj:q:automorphisms},
\ref{isc:hj:q:factors} and~\ref{isc:hj:q:formal}, in
\cref{isc:hj:sec:verify}, and in this appendix. In the text of sources~01
and~02, batch-35 \mtag{} notes were added to the title page, to
\cref{isc:sec:collection} (status of the research questions), after the
proof of \cref{isc:main:join}, after Research questions~3,
\ref{isc:cp:q:amalgams}, \ref{isc:cp:q:cores}, \ref{isc:cp:q:scales}
and~\ref{isc:cp:q:formal}, in item~1 of \cref{isc:cp:sec:answers}, at the
start of \cref{isc:sec:nonclaims}, in items~W1 and~M3, and in
Appendices~\ref{isc:app:onesource} and~\ref{isc:app:files}. The
\mtag{} tag marks text of both merges; the notes naming source~03 are
those of batch~35. No statement of sources~01 or~02 was changed.
\end{enumerate}

\emph{Where the merge chose.} The addition is printed after source~02's
sections and before the appendix, so that no existing section, theorem,
equation or question number changes; its sections keep the source's
numbering shifted by twenty-seven, its lettered theorems continue the
lettered counter, and its questions continue the question counter. The
compositum is renamed $M$ and the factors are indexed $0,1$, so that the
second route of \cref{isc:hj:rem:secondroute} reads with $\Gamma_0=G_0$,
$\Gamma_1=G_1$ and $K_0,K_1$ as in \cref{isc:sec:join}. The source's
$\K{k}{G}$ definitions of \cref{isc:hj:sec:intro,isc:hj:sec:foundations}
are printed, not replaced by references, because they state the source's
exactness claims; they agree with \cref{isc:sec:introduction,isc:sec:fractions}.

\emph{Checks in the merge.} The proofs of
\cref{isc:hj:lem:cosets,isc:hj:prop:ld,isc:hj:lem:euler,isc:hj:lem:characters,isc:hj:lem:annihilator,isc:hj:main:mixed,isc:hj:prop:quantitative,isc:hj:cor:weights,isc:hj:cor:algbase,isc:hj:lem:independentfamily,isc:hj:lem:branches,isc:hj:lem:groups,isc:hj:main:cardinal,isc:hj:cor:arithmetic,isc:hj:lem:realgroups,isc:hj:prop:archindependence,isc:hj:lem:pullback,isc:hj:prop:fullprime,isc:hj:cor:nonassociate,isc:hj:lem:classlocal,isc:hj:main:class,isc:hj:prop:strongrigidity,isc:hj:lem:derivebasis,isc:hj:thm:derivations,isc:hj:thm:automorphisms}
were read line by line when the addition was merged; no error was found.
The following were rechecked: the mixed exponents are pairwise distinct
under hypothesis~(i); the cosets $b_i+\Gamma_0$ are distinct under~(ii);
$r_\alpha-r_\beta$ has leading exponent $-\iota_0(\alpha)$ for
$\alpha<\beta$, and $2\cdot_{\Ord}\alpha$ is strictly increasing in
$\alpha$; the codes of words of length $m$ fill $[2^m-1,2^{m+1}-2]$, with
$\mathfrak B_{000\cdots}=\{1,3,7,15,\ldots\}$ and
$\mathfrak B_{111\cdots}=\{2,6,14,30,\ldots\}$; $\rho_0=2^{-3/2}+3^{-3/2}\approx
0.5460$; the cross product and the specialization $30y_1-24y_2-6y_3$ of
\cref{isc:hj:sub:verification}; the count
$\sum_{r=1}^4\binom{16}{r}2^r=34{,}112$ of the recorded Boolean patterns;
and the source's numbering (Lemma~3.1 to Theorem~10.5, Question~1) against
a build of its text, which agrees with its delivered PDF. The interleaving
used in \cref{isc:hj:sec:answers} was checked: $e_0=1>f_0=\omega^{-1}>
e_1=\omega^{-2}$, so neither of $A_\kappa$, $B_\kappa$ is infinitesimal
relative to the other. This reading is not an independent proof review.
The source's statements about this report are those of its pin; the
sections it used (\cref{isc:sec:join,isc:sec:questions}) are unchanged since,
and the later source~02 is taken into account by the \mtag{} notes.

\subsection{Files of source 03, and their rerun}\label{isc:hj:app:files}

\mtag{} The three shipped files of source~03 are these.
\begin{center}\small
\begin{tabular}{@{}>{\raggedright\arraybackslash}p{0.40\textwidth}>{\raggedright\arraybackslash}p{0.18\textwidth}>{\raggedright\arraybackslash}p{0.34\textwidth}@{}}
\toprule
Shipped path & Delivered name & Role\\
\midrule
\rpath{code/03-hahn-joins-verify.py} & \texttt{verify.py} & exact finite
checks (SymPy); prints its report, writes no file\\
\rpath{code/03-hahn-joins-build.sh} & \texttt{build.sh} & three
\texttt{pdflatex} passes, then \texttt{python3 verify.py > verification.txt}\\
\rpath{data/03-hahn-joins-verification.txt} & \texttt{verification.txt} &
recorded run: five PASS lines, Python~3.13.5, SymPy~1.14.0\\
\bottomrule
\end{tabular}
\end{center}
All three are byte-identical to the delivery. The manuscript, its README
and its PDF are not shipped. No requirements file was delivered; the
program needs Python~3 and SymPy (the recorded run used SymPy~1.14.0).

The program prints to standard output and writes no file. The delivered
\texttt{build.sh} changes into its own directory, runs \texttt{pdflatex}
three times on \texttt{article.tex}, overwriting \texttt{article.pdf}
there, and then runs \texttt{python3 verify.py > verification.txt},
rewriting that record. Under the shipped names it stops at once, as a
test on a copy confirmed: there is no \texttt{article.tex} in
\rpath{code/}, the first \texttt{pdflatex} run fails (leaving a
\texttt{texput.log}), and the script runs with \texttt{set -e}, so the
checks are never reached. An adapted copy would write \texttt{verification.txt} next
to the script, not over the shipped data file. Rerun the checks on a copy,
or redirect the output outside the report:
\begin{verbatim}
python code/03-hahn-joins-verify.py > <scratch>/verification.txt
\end{verbatim}
In the merge the program was run on a copy with Python~3.14.4 and
SymPy~1.14.0. All five checks passed; apart from line endings, the output
differs from the recorded run only in its second line, the Python version
(3.13.5 recorded).

\begin{thebibliography}{99}

\bibitem{Gonshor}
H.~Gonshor, \emph{An Introduction to the Theory of Surreal Numbers},
London Mathematical Society Lecture Note Series 110, Cambridge University
Press, 1986. Standard background on normal forms and surreal arithmetic.

\bibitem{vdDE}
L.~van den Dries and P.~Ehrlich, Fields of surreal numbers and
exponentiation, \emph{Fundamenta Mathematicae} \textbf{167} (2001), 173--188.
\href{https://doi.org/10.4064/fm167-2-3}{doi:10.4064/fm167-2-3}.
In particular pp.~174--176 for the class framework, set-cut property,
normal forms, and Hahn subfields.

\bibitem{LM}
S.~L'Innocente and V.~Mantova, A factorisation theory for generalised power
series and omnific integers, \emph{Advances in Mathematics} \textbf{442}
(2024), article 109513.
\href{https://arxiv.org/abs/1710.07304}{arXiv:1710.07304}, version 5,
22 January 2024. See Fact~2.1.1 for Hahn-field closedness, Section~2.3 for
normal forms, and Proposition~2.4.5 for the cofinality/fraction-field
obstruction. The arithmetic-intersection application here is distinct from
that published proposition.

\bibitem{Marker}
D. Marker, \emph{Model Theory: An Introduction},
Graduate Texts in Mathematics, vol.~217, Springer, 2002.
\href{https://doi.org/10.1007/b98860}{doi:10.1007/b98860}.

\bibitem{EK}
P.~Ehrlich and E.~Kaplan, \emph{Surreal ordered exponential fields},
author-hosted manuscript.
\url{https://elliotakaplan.github.io/Surreal_ordered_exponential_fields.pdf}.
Consulted as a comparison for the additional exponential-embedding
requirements, not as an input to any proof above.

\bibitem{RepoCatalogue}
VladimirReshetnikov/Surreal, documentation catalogue and project overview,
repository snapshot \snapshot, consulted 23 September 2026.
\repo{docs/README.md}; \repo{README.md}.
The catalogue distinguishes AI-assisted manuscripts from the exact scope
of checked formalization.

\bibitem{RepoAuto}
Surreal project, \emph{Omnific-Preserving Automorphisms: Convex-scale
stabilizers, definable constants, nondefinable monomials, and
algebraic-parameter rigidity}, AI-assisted merged research draft,
23 September 2026, at the same pinned snapshot.
\repo{docs/surreal/omnific-preserving-automorphisms/article.tex}.
Used for scope and nonduplication comparison, not as an unproved dependency.

\bibitem{RepoBounded}
Surreal project, \emph{Bounded-Support Hahn Arithmetic: Cofinal descent,
optimal support thresholds, and maximal transcendence}, AI-assisted
research draft, 22 September 2026, at the same pinned snapshot.
\repo{docs/surreal/transcendence-over-bounded-support/article.tex}.
Its main theorems propose stronger transcendence and descent results;
\cref{isc:sec:gaps} gives the needed one-witness argument independently.

\bibitem{Wiki}
Wikipedia contributors, \emph{Surreal number}, user-supplied background
page, consulted 23 September 2026.
\url{https://en.wikipedia.org/wiki/Surreal_number}.
This reference is orientation only; the proofs rely on the named
mathematical foundations above rather than on this page.


\bibitem{KKS}
L.~S.~Krapp, S.~Kuhlmann and M.~Serra, \emph{Generalised power series
determined by linear recurrence relations},
\href{https://arxiv.org/abs/2206.04126v3}{arXiv:2206.04126v3}, 2025.
Cited by source~02 for the literature comparison with recurrence-defined
series substructures, not as a proof of its compositum theorems.

\bibitem{KM}
S.~Kuhlmann and M.~Matusinski, \emph{Hardy type derivations on generalized
series fields},
\href{https://arxiv.org/abs/0903.2197v4}{arXiv:0903.2197v4}, 2012.
\href{https://doi.org/10.1016/j.jalgebra.2011.11.024}{doi:10.1016/j.jalgebra.2011.11.024}.
Cited by source~02 for the background distinction between series
derivations and Hardy-type conditions.

\bibitem{BM}
A.~Berarducci and V.~Mantova, \emph{Surreal numbers, derivations and
transseries}, \emph{Journal of the European Mathematical Society}
\textbf{20} (2018), 339--390.
\href{https://arxiv.org/abs/1503.00315v3}{arXiv:1503.00315v3}.
\href{https://doi.org/10.4171/JEMS/769}{doi:10.4171/JEMS/769}.

\bibitem{StacksLD}
The Stacks Project Authors, \emph{The Stacks Project}, Section~9.27,
``Linearly disjoint extensions'', Tag~09IC.
\url{https://stacks.math.columbia.edu/tag/09IC}.

\bibitem{Conway}
J.~H.~Conway, \emph{On Numbers and Games}, Academic Press, London, 1976.
Cited by source~03.

\bibitem{Hahn}
H.~Hahn, \emph{\"Uber die nichtarchimedischen Gr\"o\ss ensysteme},
Sitzungsberichte der Kaiserlichen Akademie der Wissenschaften,
Mathematisch-Naturwissenschaftliche Klasse, Abteilung IIa
\textbf{116} (1907), 601--655. Cited by source~03.

\bibitem{Kaplansky}
I.~Kaplansky, \emph{Maximal fields with valuations},
Duke Mathematical Journal \textbf{9} (1942), 303--321. Cited by source~03.

\bibitem{Pitteloud}
D.~Pitteloud, \emph{Existence of prime elements in rings of generalized
power series}, Journal of Symbolic Logic \textbf{66}(3) (2001), 1206--1216.
\href{https://doi.org/10.2307/2695102}{doi:10.2307/2695102}.
Cited by source~03 for the one-row prime theorem.

\bibitem{BKK}
D.~Biljakovic, M.~Kochetov and S.~Kuhlmann, \emph{Primes and irreducibles
in truncation integer parts of real closed fields}, in A.~Enayat,
I.~Kalantari and M.~Moniri (eds.), \emph{Logic in Tehran}, Lecture Notes in
Logic 26, A K Peters, 2006, 42--64; Theorem~4.12, as cited in
\cite[\S1.3]{LM}.
\href{https://doi.org/10.1017/9781316755747.006}{doi:10.1017/9781316755747.006}.
Cited by source~03 for the transfer of one-row primality.

\end{thebibliography}
\end{document}
